% Description of computer networking — ICT Upper Sixth — Cours signé OnBuch+
% Official MINESEC syllabus. Compiler avec : tectonic ICT02-description-of-computer-networking.tex
\newcommand{\DOCMATIERE}{ICT}
\newcommand{\DOCNIVEAU}{Upper Sixth}
\newcommand{\DOCMODULE}{Upper Sixth — ICT}
\newcommand{\DOCLECON}{2}
\newcommand{\DOCTITRE}{Description of computer networking}
% OnBuch+ — préambule « cours signés » ENRICHI (compilé avec tectonic / XeLaTeX)
% Système visuel « Pop » d'OnBuch+ : fond crème (jamais blanc), bordures encre
% épaisses, ombres « dures » décalées, titres Archivo Black en capitales.
% Version MATHÉMATIQUES : graphes (pgfplots/tikz), boîtes pédagogiques
% (définition, propriété, méthode, attention, le savais-tu, exercice résolu,
% exercice, TP, bilan), page de garde et signature OnBuch+ ancrée en bas.
%
% Macros attendues dans main.tex AVANT \input{preamble} :
%   \DOCMATIERE  \DOCNIVEAU  \DOCMODULE  \DOCLECON (numéro)  \DOCTITRE
\documentclass[11pt]{article}

\usepackage{fontspec}
\usepackage[english]{babel}
\usepackage[a4paper,margin=2cm,top=3.4cm,bottom=2.8cm,headheight=1.4cm,footskip=1.3cm]{geometry}
\usepackage{amsmath,amssymb,amsfonts}
\usepackage{mathtools} % \xmapsto, \coloneqq... souvent utilisés par les modèles
\usepackage{cancel} % simplifications barrées (bilans de forces)
% \abs{z} (module, valeur absolue) et \norm{u} (norme), utilisés par les modèles
\providecommand{\abs}{}\let\abs\relax\DeclarePairedDelimiter{\abs}{\lvert}{\rvert}
\providecommand{\norm}{}\let\norm\relax\DeclarePairedDelimiter{\norm}{\lVert}{\rVert}
\usepackage[table]{xcolor} % [table] : fournit aussi \rowcolors (alternance de lignes)
\usepackage{pagecolor}
\usepackage{tikz}
\usetikzlibrary{arrows.meta,positioning,calc,shapes.geometric,shapes.misc,shapes.arrows,decorations.pathmorphing,decorations.pathreplacing,decorations.markings,patterns,shadows,mindmap,trees,fit,backgrounds,matrix,intersections,angles,quotes}
\usepackage{pgfplots}
\usepackage[version=4]{mhchem} % équations nucléaires \ce{^{235}_{92}U}
\usepackage[european,siunitx]{circuitikz} % schémas électriques
\pgfplotsset{compat=1.18}
\usepgfplotslibrary{fillbetween,groupplots} % aires entre courbes (calcul intégral)
\usepackage{enumitem}
\usepackage{fancyhdr}
% [most] charge toutes les bibliothèques tcolorbox (listings, minted, raster,
% documentation...) et gonfle inutilement la mémoire TeX sur un document long
% (~300 boîtes) : on ne charge que ce qui est réellement utilisé.
\usepackage[skins,breakable]{tcolorbox}
\usepackage{graphicx}
\usepackage{colortbl}
\usepackage{booktabs}
\usepackage{tabularx}
\usepackage{diagbox} % case d'en-tête diagonale des tableaux à double entrée
% Colonnes extensibles centrée / gauche / droite (C, L, R), souvent utilisées par les modèles
\newcolumntype{C}{>{\centering\arraybackslash}X}
\newcolumntype{L}{>{\raggedright\arraybackslash}X}
\newcolumntype{R}{>{\raggedleft\arraybackslash}X}
\usepackage{array}
\usepackage{multicol}
\usepackage{siunitx}
% parse-numbers=false : accepte aussi \SI{2,5 \times 10^{-3}}{...}
\sisetup{locale=UK,output-decimal-marker={.},group-minimum-digits=4,parse-numbers=false}
\DeclareSIUnit\ohm{\ensuremath{\symup{\Omega}}} % U+2126 absent de la police
\DeclareSIUnit\division{div} % réglages d'oscilloscope (ms/div, V/div)
\DeclareSIUnit\tr{tr} % tours (vitesse de rotation, ex. \SI{1500}{\tr\per\minute})
\usepackage{qrcode}
% \degree : commande fréquemment utilisée par les modèles pour un angle ou une
% température en dehors de \SI{}{} (non fournie par les paquets chargés).
\providecommand{\degree}{^{\circ}}
% \qed (fin de démonstration), utilisé par les modèles sans amsthm
\providecommand{\qed}{\hfill$\square$}
% \barre : double barre des valeurs interdites dans un tableau de variations (tkz-tab)
\providecommand{\barre}{\ensuremath{\|}}
% environnement proof (amsthm non chargé) utilisé par les modèles
\ifdefined\proof\else\newenvironment{proof}[1][Proof]{\par\noindent\textit{#1.}\ }{\qed\par}\fi
\usepackage{lastpage}
\usepackage{hyperref}
\hypersetup{hidelinks}

% ============================================================
% Palette « Pop » OnBuch+ — mêmes hex que la classe P (plus_ui.dart)
% ============================================================
\definecolor{poporange}{HTML}{F59321}
\definecolor{popdark}{HTML}{E07A0C}
\definecolor{popcream}{HTML}{FFF6E8}
\definecolor{popink}{HTML}{1C1714}
\definecolor{poppurple}{HTML}{7A5AE0}
\definecolor{popgreen}{HTML}{2E9E6A}
\definecolor{poppink}{HTML}{E0457B}
\definecolor{popblue}{HTML}{2F7DE1}
\definecolor{popgold}{HTML}{C98A12}
\definecolor{popmuted}{HTML}{8A7F76}
\definecolor{popline}{HTML}{EADFD2}
\definecolor{popgray}{HTML}{8A7F76}
\colorlet{popgrey}{popgray}
% Alias tolérants (le modèle nomme parfois les couleurs différemment)
\colorlet{popwhite}{white}
\colorlet{popblack}{popink}
\colorlet{popred}{poppink}
\colorlet{popyellow}{popgold}
% Teintes claires pour fonds de tableaux / zones
\colorlet{popblueL}{popblue!10!white}
\colorlet{poporangeL}{poporange!14!white}
\colorlet{popgreenL}{popgreen!12!white}
\colorlet{poppurpleL}{poppurple!12!white}
\colorlet{poppinkL}{poppink!12!white}
\colorlet{popgoldL}{popgold!14!white}

\pagecolor{popcream}

% ============================================================
% Typographie
% ============================================================
\defaultfontfeatures{Ligatures=TeX}
\setmainfont{PlusJakartaSans-Medium.ttf}[Path=../../../fonts/,BoldFont=PlusJakartaSans-Bold.ttf]
% Maths en sans-serif (Fira Math) avec chiffres et lettres droites de Plus
% Jakarta, pour que les formules se fondent dans le texte.
\usepackage[math-style=ISO,bold-style=ISO]{unicode-math}
\setmathfont{FiraMath-Regular.otf}[Path=../../../fonts/]
\setmathfont{PlusJakartaSans-Medium.ttf}[Path=../../../fonts/,range={up/num,up/{Latin,latin},\mathup}]
% (icomma retiré : décimales avec point en anglais)
% Caractères mathématiques Unicode absents des polices de texte (Plus Jakarta,
% Archivo Black) : rendus par leur équivalent mathématique, en texte comme en
% formule — sinon ils s'affichent en carré vide (ex. « Arithmétique dans ℤ »).
\usepackage{newunicodechar}
\newunicodechar{ℝ}{\ensuremath{\mathbb{R}}}
\newunicodechar{ℤ}{\ensuremath{\mathbb{Z}}}
\newunicodechar{ℂ}{\ensuremath{\mathbb{C}}}
\newunicodechar{ℕ}{\ensuremath{\mathbb{N}}}
\newunicodechar{ℚ}{\ensuremath{\mathbb{Q}}}
\newunicodechar{𝔻}{\ensuremath{\mathbb{D}}}
\newunicodechar{α}{\ensuremath{\alpha}}
\newunicodechar{β}{\ensuremath{\beta}}
\newunicodechar{γ}{\ensuremath{\gamma}}
\newunicodechar{δ}{\ensuremath{\delta}}
\newunicodechar{ε}{\ensuremath{\varepsilon}}
\newunicodechar{θ}{\ensuremath{\theta}}
\newunicodechar{λ}{\ensuremath{\lambda}}
\newunicodechar{μ}{\ensuremath{\mu}}
\newunicodechar{σ}{\ensuremath{\sigma}}
\newunicodechar{τ}{\ensuremath{\tau}}
\newunicodechar{φ}{\ensuremath{\varphi}}
\newunicodechar{ψ}{\ensuremath{\psi}}
\newunicodechar{ω}{\ensuremath{\omega}}
\newunicodechar{Σ}{\ensuremath{\Sigma}}
% majuscules produites par \MakeUppercase dans les titres (α-aminés -> Α)
\newunicodechar{Α}{\ensuremath{\alpha}}
\newunicodechar{Β}{\ensuremath{\beta}}
\newunicodechar{Γ}{\ensuremath{\Gamma}}
\newunicodechar{Δ}{\ensuremath{\Delta}}
\newunicodechar{Ε}{\ensuremath{\varepsilon}}
\newunicodechar{Θ}{\ensuremath{\Theta}}
\newunicodechar{Λ}{\ensuremath{\Lambda}}
\newunicodechar{Μ}{\ensuremath{\mu}}
\newunicodechar{Τ}{\ensuremath{\tau}}
\newunicodechar{Φ}{\ensuremath{\Phi}}
\newunicodechar{Ψ}{\ensuremath{\Psi}}
\newunicodechar{Ω}{\ensuremath{\Omega}}
\newunicodechar{↦}{\ensuremath{\mapsto}}
\newunicodechar{∈}{\ensuremath{\in}}
\newunicodechar{∉}{\ensuremath{\notin}}
\newunicodechar{∧}{\ensuremath{\wedge}}
\newunicodechar{⇔}{\ensuremath{\Leftrightarrow}}
\newunicodechar{⟺}{\ensuremath{\Longleftrightarrow}}
\newunicodechar{⇒}{\ensuremath{\Rightarrow}}
\newunicodechar{⟹}{\ensuremath{\Longrightarrow}}
\newunicodechar{∘}{\ensuremath{\circ}}
\newunicodechar{∪}{\ensuremath{\cup}}
\newunicodechar{∩}{\ensuremath{\cap}}
\newunicodechar{≡}{\ensuremath{\equiv}}
\newunicodechar{⊂}{\ensuremath{\subset}}
\newunicodechar{⊥}{\ensuremath{\perp}}
\newunicodechar{∃}{\ensuremath{\exists}}
\newunicodechar{∀}{\ensuremath{\forall}}
\newunicodechar{∛}{\ensuremath{\sqrt[3]{\,}}}
\newunicodechar{∜}{\ensuremath{\sqrt[4]{\,}}}
\newunicodechar{⃗}{\ensuremath{{}^{\to}}}
\newunicodechar{ⁿ}{\ifmmode^{n}\else\textsuperscript{\ensuremath{n}}\fi}
\newunicodechar{ˣ}{\ifmmode^{x}\else\textsuperscript{\ensuremath{x}}\fi}
\newunicodechar{ᵇ}{\ifmmode^{b}\else\textsuperscript{\ensuremath{b}}\fi}
\newunicodechar{ʳ}{\ifmmode^{r}\else\textsuperscript{\ensuremath{r}}\fi}
\newunicodechar{ᵘ}{\ifmmode^{u}\else\textsuperscript{\ensuremath{u}}\fi}
\newunicodechar{ʸ}{\ifmmode^{y}\else\textsuperscript{\ensuremath{y}}\fi}
\newunicodechar{ᵃ}{\ifmmode^{a}\else\textsuperscript{\ensuremath{a}}\fi}
\newunicodechar{ᵉ}{\ifmmode^{e}\else\textsuperscript{\ensuremath{e}}\fi}
\newunicodechar{ᵏ}{\ifmmode^{k}\else\textsuperscript{\ensuremath{k}}\fi}
\newunicodechar{ᵖ}{\ifmmode^{p}\else\textsuperscript{\ensuremath{p}}\fi}
\newunicodechar{ᴺ}{\ifmmode^{N}\else\textsuperscript{\ensuremath{N}}\fi}
\newunicodechar{ᶜ}{\ifmmode^{c}\else\textsuperscript{\ensuremath{c}}\fi}
\newunicodechar{ʰ}{\ifmmode^{h}\else\textsuperscript{\ensuremath{h}}\fi}
\newunicodechar{ᵠ}{\ifmmode^{q}\else\textsuperscript{\ensuremath{q}}\fi}
\newunicodechar{ₙ}{\ifmmode_{n}\else\textsubscript{\ensuremath{n}}\fi}
\newunicodechar{ₐ}{\ifmmode_{a}\else\textsubscript{\ensuremath{a}}\fi}
\newunicodechar{ᵢ}{\ifmmode_{i}\else\textsubscript{\ensuremath{i}}\fi}
\newunicodechar{ᵣ}{\ifmmode_{r}\else\textsubscript{\ensuremath{r}}\fi}
\newunicodechar{ₖ}{\ifmmode_{k}\else\textsubscript{\ensuremath{k}}\fi}
\newunicodechar{ᵦ}{\ifmmode_{\beta}\else\textsubscript{\ensuremath{\beta}}\fi}
\newunicodechar{ₓ}{\ifmmode_{x}\else\textsubscript{\ensuremath{x}}\fi}
\newfontfamily\popdisplay{ArchivoBlack-Regular.ttf}[Path=../../../fonts/]
\newfontfamily\poplogo{SpaceGrotesk-Bold.ttf}[Path=../../../fonts/]
\newfontfamily\popsemi{PlusJakartaSans-SemiBold.ttf}[Path=../../../fonts/]
\newfontfamily\popxbold{PlusJakartaSans-ExtraBold.ttf}[Path=../../../fonts/]
\newfontfamily\popxboldls{PlusJakartaSans-ExtraBold.ttf}[Path=../../../fonts/,LetterSpace=3.0]

\setlist[enumerate]{leftmargin=1.7em,itemsep=5pt,topsep=4pt}
\setlist[itemize]{leftmargin=1.7em,itemsep=4pt,topsep=4pt,label={\color{poporange}\textbullet}}
\setlength{\parindent}{0pt}
\setlength{\parskip}{5pt}
\renewcommand{\arraystretch}{1.25}

% pgfplots : style Pop par défaut
\pgfplotsset{
  every axis/.append style={axis line style={popink,line width=1pt},tick style={popink},
    grid style={popline,line width=0.5pt},label style={font=\small},tick label style={font=\footnotesize}},
  popcurve/.style={poporange,line width=1.8pt,no markers,smooth},
}
% Même style disponible dans un \draw TikZ simple (hors pgfplots)
\tikzset{popcurve/.style={poporange,line width=1.8pt}}
\tikzset{popgrid/.style={popline,very thin}}
\colorlet{popcurve}{poporange} % le nom du style est parfois utilisé comme couleur

% ============================================================
% Pill / squiggle
% ============================================================
\newcommand{\poppill}[3]{%
  \tikz[baseline=-0.55ex]\node[draw=popink,line width=1.2pt,rounded corners=2.6pt,%
    fill=#2,inner xsep=6.5pt,inner ysep=3pt]{%
    \popxboldls\fontsize{7}{7}\selectfont\color{#3}\MakeUppercase{#1}};%
}
\newcommand{\popsquiggle}[1]{%
  \tikz[baseline=-0.4ex]{
    \draw[#1,line width=2.6pt,line cap=round]
      plot [smooth] coordinates {(0,0.09) (0.34,0.01) (0.68,0.21) (1.02,0.01) (1.36,0.10)};
  }%
}
% Mot-clé surligné (marqueur orange sous le texte)
\newcommand{\cle}[1]{{\popxbold\color{popdark}#1}}

% ============================================================
% En-tête / pied
% ============================================================
\pagestyle{fancy}
\fancyhf{}
\renewcommand{\headrulewidth}{0pt}
\renewcommand{\footrulewidth}{0pt}
\lhead{\raisebox{-0.15ex}{\poplogo\fontsize{13}{13}\selectfont\color{popink}OnBuch\color{poporange}+}}
\rhead{\poppill{\DOCMATIERE\ \textbullet\ \DOCNIVEAU\ \textbullet\ Chapter \DOCLECON}{white}{popink}}
\lfoot{\popxbold\fontsize{7.6}{7.6}\selectfont\color{popmuted}\MakeUppercase{Signed course OnBuch+}\ \textbullet\ {\popsemi\DOCTITRE}}
\rfoot{\tikz[baseline=-0.6ex]\node[rounded corners=8.5pt,fill=popink,minimum height=17pt,minimum width=17pt,inner xsep=5pt,inner sep=0]{\popxbold\fontsize{8}{8}\selectfont\color{popcream}\thepage\,/\,\pageref*{LastPage}};}

% ============================================================
% Titres
% ============================================================
\newcommand{\chaptitre}[2]{%
  \begin{tcolorbox}[enhanced,colback=poporange,colframe=popink,boxrule=2.6pt,arc=11pt,
      drop shadow=popink,left=15pt,right=15pt,top=12pt,bottom=11pt,before skip=3pt,after skip=18pt]
    \poppill{#2}{popink}{popcream}\par\vspace{9pt}
    {\raggedright\hyphenpenalty=10000\exhyphenpenalty=10000\popdisplay\fontsize{22}{24}\selectfont\color{popcream}\MakeUppercase{#1}\par}
  \end{tcolorbox}
}
% Section numérotée : pastille orange avec le numéro + titre Archivo
\newcounter{popsec}
\newcommand{\coursec}[1]{%
  \stepcounter{popsec}\setcounter{popsub}{0}%
  \par\vspace{16pt}\noindent
  \begin{minipage}{\linewidth}
  \tikz[baseline=-0.7ex]\node[draw=popink,line width=1.6pt,fill=poporange,rounded corners=4pt,minimum size=22pt,inner sep=0]{\popdisplay\fontsize{12}{12}\selectfont\color{popcream}\thepopsec};%
  \hspace{8pt}{\popdisplay\fontsize{15}{17}\selectfont\color{popink}\MakeUppercase{#1}}\par
  \vspace{4pt}\noindent{\color{poporange}\rule{36pt}{3.6pt}}
  \end{minipage}\par\nopagebreak\vspace{8pt}
}
\newcounter{popsub}
\newcommand{\courssub}[1]{%
  \stepcounter{popsub}%
  \par\vspace{9pt}\noindent{\popxbold\fontsize{11.5}{13}\selectfont\color{popink}{\color{poporange}\thepopsec.\thepopsub}\hspace{6pt}#1}\par\nopagebreak\vspace{4pt}
}

% ============================================================
% Boîtes pédagogiques — même recette : fond blanc, bordure épaisse colorée,
% ombre dure encre, badge plein. Titre optionnel [#1].
% ============================================================
\tcbset{popbase/.style={breakable,enhanced,colback=white,boxrule=2.3pt,arc=9pt,
  shadow={3pt}{-3pt}{0pt}{popink},left=14pt,right=14pt,top=11pt,bottom=11pt,before skip=11pt,after skip=15pt,
  fontupper=\popsemi\fontsize{10.6}{14.5}\selectfont\color{popink},
  fontlower=\popsemi\fontsize{10.6}{14.5}\selectfont\color{popink}}}
% Coche dessinée (le glyphe ✓ n'existe pas dans les polices du document)
\newcommand{\popcheck}[1]{\tikz[baseline=-0.5ex]\draw[#1,line width=1.6pt,line cap=round,line join=round](-0.13,0.0)--(-0.04,-0.09)--(0.13,0.1);}
\providecommand{\coche}{\popcheck{popgreen}}
\providecommand{\resultat}[1]{\par\smallskip\noindent\fcolorbox{popgreen}{popgreenL}{\parbox{\dimexpr\linewidth-2\fboxsep-2\fboxrule\relax}{\textbf{Result.} #1}}\par\smallskip}
\newcommand{\popboxhead}[3]{\poppill{#1}{#2}{white}\ifx\relax#3\relax\else\hspace{6pt}{\popxbold\fontsize{10.6}{12}\selectfont\color{popink}#3}\fi\par\vspace{7pt}}

\newtcolorbox{aretenir}[1][]{popbase,colframe=popgreen,before upper={\popboxhead{Key points}{popgreen}{#1}}}
\newtcolorbox{exemplebox}[1][]{popbase,colframe=poporange,before upper={\popboxhead{Exemple}{poporange}{#1}}}
\newtcolorbox{definition}[1][]{popbase,colframe=popblue,before upper={\popboxhead{Definition}{popblue}{#1}}}
\newtcolorbox{propriete}[1][]{popbase,colframe=poppurple,before upper={\popboxhead{Property}{poppurple}{#1}}}
\newtcolorbox{methode}[1][]{popbase,colframe=popgold,before upper={\popboxhead{Method}{popgold}{#1}}}
\newtcolorbox{attention}[1][]{popbase,colframe=poppink,before upper={\popboxhead{Warning}{poppink}{#1}}}
\newtcolorbox{experience}[1][]{popbase,colframe=popblue,colback=popblueL,before upper={\popboxhead{Experiment}{popblue}{#1}}}
% « Le savais-tu ? » : carte violette pleine (contexte, culture, Cameroun)
\newtcolorbox{savaistu}[1][]{popbase,colback=poppurple,colframe=popink,
  fontupper=\popsemi\fontsize{10.4}{14.2}\selectfont\color{white},
  before upper={\poppill{Did you know?}{white}{poppurple}\ifx\relax#1\relax\else\hspace{6pt}{\popxbold\color{white}#1}\fi\par\vspace{7pt}}}
% Exercice résolu : énoncé en haut, \tcblower puis la solution sur fond crème
\newcounter{popexo}\newcounter{popexr}
\newtcolorbox{exoresolu}[1][]{popbase,colframe=poporange,colbacklower=popcream,
  segmentation style={popink,line width=1.2pt,dashed},
  before upper={\stepcounter{popexr}\popboxhead{Worked example \thepopexr}{poporange}{#1}},
  before lower={\poppill{Solution}{popink}{popcream}\par\vspace{6pt}}}
% Exercice (non résolu en place) — la correction est dans la partie Corrigés
\newtcolorbox{exercice}[1][]{popbase,colframe=popink,
  before upper={\stepcounter{popexo}\popboxhead{Exercise \thepopexo}{popink}{#1}}}
% Corrigé
\newtcolorbox{corrige}[1][]{popbase,colframe=popgreen,colback=popgreenL,
  before upper={\popboxhead{Solution}{popgreen}{#1}}}
% Objectifs / prérequis / bilan
\newtcolorbox{objectifs}{popbase,colframe=popink,colback=poporangeL,
  before upper={\popboxhead{Chapter objectives}{popink}{}}}
\newtcolorbox{prerequis}{popbase,colframe=popink,colback=popblueL,
  before upper={\popboxhead{Prerequisites}{popblue}{}}}
\newtcolorbox{bilan}{popbase,colframe=popink,colback=white,boxrule=2.8pt,
  before upper={\popboxhead{Fiche bilan}{poporange}{L'essentiel à connaître pour l'examen}}}
% Figure encadrée avec légende
\newtcolorbox{popfigure}[1][]{enhanced,breakable=false,colback=white,colframe=popline,boxrule=1.4pt,arc=8pt,
  left=10pt,right=10pt,top=10pt,bottom=8pt,before skip=10pt,after skip=12pt,halign=center,
  lower separated=false,halign lower=center,
  fontlower=\popsemi\fontsize{9}{11.5}\selectfont\color{popmuted},
  #1}
\newcommand{\legende}[1]{\par\vspace{6pt}{\popsemi\fontsize{9}{11.5}\selectfont\color{popmuted}#1}}

% ============================================================
% Signature OnBuch+
% ============================================================
\newcommand{\poplogoinline}[1]{{\poplogo\fontsize{#1}{#1}\selectfont\color{popink}OnBuch\color{poporange}+}}
\newcommand{\signatureonbuch}{%
  \begin{tcolorbox}[enhanced,colback=popink,colframe=popink,boxrule=2.2pt,arc=10pt,
      drop shadow=poporange,left=14pt,right=14pt,top=10pt,bottom=10pt,before skip=0pt,after skip=0pt]
    \tikz[baseline=-0.7ex]\node[circle,fill=popgreen,draw=popcream,line width=1.3pt,minimum size=22pt,inner sep=0]{\popcheck{white}};%
    \hspace{9pt}\begin{minipage}[c]{0.62\linewidth}
      {\popxbold\fontsize{12}{13}\selectfont\color{popcream}Signed\ \textbullet\ {\poplogo OnBuch\color{poporange}+}}\par\vspace{2pt}
      {\raggedright\popsemi\fontsize{8}{10.5}\selectfont\color{popline}\MakeUppercase{OnBuch+ teaching team}\\\MakeUppercase{Follows the official MINESEC syllabus}\par}
    \end{minipage}\hfill
    {\poplogo\fontsize{22}{22}\selectfont\color{popcream}OnBuch\color{poporange}+}
  \end{tcolorbox}%
}

% ============================================================
% Page de garde
% #1 titre · #2 matière · #3 niveau · #4 module · #5 n° leçon · #6 durée ·
% #7 description / objectifs (texte ou itemize)
% ============================================================
\newcommand{\pagegarde}[7]{%
  \thispagestyle{empty}
  \newgeometry{a4paper,margin=1.7cm,top=1.6cm,bottom=1.4cm}%
  \begin{tikzpicture}[remember picture,overlay]
    \fill[poporange!18] ([xshift=-3.2cm,yshift=-3.6cm]current page.north east) circle (4.6cm);
    \fill[poppurple!14] ([xshift=2.4cm,yshift=7.6cm]current page.south west) circle (3.8cm);
  \end{tikzpicture}%
  {\poplogo\fontsize{30}{30}\selectfont\color{popink}OnBuch\color{poporange}+}\hfill
  \raisebox{0.6ex}{\poppill{Signed course \textbullet\ Premium}{popink}{popcream}}\par
  \vspace{1pt}\popsquiggle{poppurple}\par
  \vspace{8pt}
  \begin{tcolorbox}[enhanced,colback=poporange,colframe=popink,boxrule=2.8pt,arc=13pt,
      drop shadow=popink,left=18pt,right=18pt,top=12pt,bottom=13pt,after skip=10pt]
    \poppill{#2 \textbullet\ #3}{popink}{popcream}\hspace{5pt}\poppill{#4}{white}{popink}\par\vspace{8pt}
    {\popdisplay\fontsize{14}{14}\selectfont\color{popink}CHAPTER #5}\par\vspace{4pt}
    {\raggedright\hyphenpenalty=10000\exhyphenpenalty=10000\popdisplay\fontsize{25}{27}\selectfont\color{popcream}\MakeUppercase{#1}\par}
    \vspace{8pt}
    {\popxbold\fontsize{9.5}{9.5}\selectfont\color{popink}\MakeUppercase{Suggested time: #6}}
  \end{tcolorbox}
  \begin{tcolorbox}[enhanced,colback=white,colframe=popink,boxrule=2.2pt,arc=10pt,
      drop shadow=popink,left=16pt,right=16pt,top=10pt,bottom=10pt,after skip=8pt]
    \poppill{In this chapter}{poppurple}{white}\par\vspace{5pt}
    \popsemi\fontsize{9.3}{12.6}\selectfont\color{popink}#7
  \end{tcolorbox}
  \noindent\begin{minipage}[t]{0.487\linewidth}
    \begin{tcolorbox}[enhanced,colback=popgreen,colframe=popink,boxrule=2.2pt,arc=10pt,
        drop shadow=popink,left=11pt,right=11pt,top=10pt,bottom=10pt,height=3.1cm,valign=center]
      \tcbox[colback=white,colframe=popink,boxrule=1.3pt,arc=5pt,boxsep=0pt,nobeforeafter,
        left=3pt,right=3pt,top=3pt,bottom=3pt,valign=center]{\qrcode[height=1.9cm]{https://play.google.com/store/apps/details?id=com.luvvix.onbuch&hl=en}}%
      \hspace{8pt}\begin{minipage}[c]{0.5\linewidth}
        {\popdisplay\fontsize{11}{12.5}\selectfont\color{white}\MakeUppercase{Download OnBuch}}\par\vspace{4pt}
        {\raggedright\popsemi\fontsize{8.4}{11}\selectfont\color{white}Free \textbullet\ Google Play\\AI tutor, past papers, results\par}
      \end{minipage}
    \end{tcolorbox}
  \end{minipage}\hfill
  \begin{minipage}[t]{0.487\linewidth}
    \begin{tcolorbox}[enhanced,colback=poppurple,colframe=popink,boxrule=2.2pt,arc=10pt,
        drop shadow=popink,left=11pt,right=11pt,top=10pt,bottom=10pt,height=3.1cm,valign=center]
      \tikz[baseline=-0.7ex]\node[circle,fill=white,draw=popink,line width=1.3pt,minimum size=1.6cm,inner sep=0]{\popdisplay\fontsize{24}{24}\selectfont\color{poppurple}+};%
      \hspace{8pt}\begin{minipage}[c]{0.6\linewidth}
        {\popdisplay\fontsize{11}{12.5}\selectfont\color{white}\MakeUppercase{Go OnBuch+}}\par\vspace{4pt}
        {\raggedright\popsemi\fontsize{8.4}{11}\selectfont\color{white}All signed courses, unlimited AI tutor, PDF sheets — OnBuch+ tab in the app\par}
      \end{minipage}
    \end{tcolorbox}
  \end{minipage}
  \vspace{8pt}
  \popcontenu
  \vfill
  \signatureonbuch
  \restoregeometry
  \newpage
}

% Rangée « Ce cours contient » (page de garde)
\newcommand{\poptile}[3]{% #1 couleur #2 gros texte #3 légende
  \begin{tcolorbox}[enhanced,colback=#1,colframe=popink,boxrule=2pt,arc=9pt,shadow={2.5pt}{-2.5pt}{0pt}{popink},
      left=6pt,right=6pt,top=9pt,bottom=9pt,halign=center,valign=center,height=1.75cm,width=\linewidth,nobeforeafter]
    {\popdisplay\fontsize{12.5}{13}\selectfont\color{white}\MakeUppercase{#2}}\par\vspace{3pt}
    {\centering\popsemi\fontsize{7.8}{9.5}\selectfont\color{white}#3\par}
  \end{tcolorbox}}
\newcommand{\popcontenu}{%
  \par\noindent{\popxboldls\fontsize{7.5}{7.5}\selectfont\color{popmuted}THIS SIGNED COURSE CONTAINS}\par\vspace{7pt}
  \noindent\begin{minipage}[t]{0.235\linewidth}\poptile{popblue}{Course}{complete, illustrated, step by step}\end{minipage}\hfill
  \begin{minipage}[t]{0.235\linewidth}\poptile{poporange}{Methods}{and worked examples}\end{minipage}\hfill
  \begin{minipage}[t]{0.235\linewidth}\poptile{poppink}{Exercises}{exam style + solutions}\end{minipage}\hfill
  \begin{minipage}[t]{0.235\linewidth}\poptile{popgreen}{Summary sheet}{the essentials to remember}\end{minipage}\par}

% Sommaire
\newtcolorbox{sommaire}{enhanced,colback=white,colframe=popink,boxrule=2.2pt,arc=10pt,
  drop shadow=popink,left=18pt,right=18pt,top=13pt,bottom=13pt,before skip=6pt,after skip=20pt,
  before upper={\poppill{Contents}{poppurple}{white}\par\vspace{9pt}\popsemi\fontsize{10.6}{14.5}\selectfont\color{popink}}}

% Signature de fin de cours — poussée tout en bas de la dernière page
\newcommand{\findecours}{%
  \par\vfill\nopagebreak
  \begin{center}\popsquiggle{poporange}\end{center}\vspace{4pt}
  \signatureonbuch
}
\AtBeginDocument{\def\llbracket{\mathopen{[\mkern-3.5mu[}}\def\rrbracket{\mathclose{]\mkern-3.5mu]}}} % intervalles d'entiers (glyphe ⟦ absent de la police)
% Glyphes absents de Fira Math : redéfinis à partir de glyphes disponibles
\AtBeginDocument{%
  \def\setminus{\mathbin{\backslash}}\let\smallsetminus\setminus
  \def\vdots{\mathord{\vbox{\baselineskip3.2pt\lineskiplimit0pt\kern4pt\hbox{.}\hbox{.}\hbox{.}}}}%
  \def\ddots{\mathinner{\mkern1mu\raise6.5pt\hbox{.}\mkern2mu\raise3.6pt\hbox{.}\mkern2mu\raise0.7pt\hbox{.}\mkern1mu}}%
  \def\triangle{\mathord{\scalebox{0.9}{$\symup{\Delta}$}}}%
  \def\nabla{\mathord{\raisebox{\depth}{\scalebox{1}[-1]{$\symup{\Delta}$}}}}%
  \def\checkmark{\popcheck{popdark}}%
  \def\star{\mathbin{\ast}}\def\bigstar{\mathord{\ast}}%
  \def\diamond{\mathbin{\scalebox{0.6}{$\Diamond$}}}%
  \def\bigcup{\mathop{\mathchoice{\raisebox{-0.3em}{\scalebox{1.7}{$\cup$}}}{\scalebox{1.25}{$\cup$}}{\cup}{\cup}}\displaylimits}%
  \def\bigcap{\mathop{\mathchoice{\raisebox{-0.3em}{\scalebox{1.7}{$\cap$}}}{\scalebox{1.25}{$\cap$}}{\cap}{\cap}}\displaylimits}%
  \def\lesseqqgtr{\mathrel{\lessgtr}}\def\bigoplus{\mathop{\oplus}\displaylimits}%
}
\providecommand{\ssi}{if and only if} % abréviation fréquente
\AtBeginDocument{\providecommand{\wideparen}{\overparen}} % arc de cercle

% Fonctions usuelles que les modèles utilisent sans que amsmath les fournisse
\providecommand{\arsinh}{\operatorname{arsinh}}\providecommand{\arcosh}{\operatorname{arcosh}}
\providecommand{\artanh}{\operatorname{artanh}}\providecommand{\arsech}{\operatorname{arsech}}
\providecommand{\arcsch}{\operatorname{arcsch}}\providecommand{\arcoth}{\operatorname{arcoth}}
\providecommand{\arcsinh}{\operatorname{arsinh}}\providecommand{\arccosh}{\operatorname{arcosh}}
\providecommand{\arctanh}{\operatorname{artanh}}\providecommand{\sech}{\operatorname{sech}}
\providecommand{\csch}{\operatorname{csch}}\providecommand{\arcsec}{\operatorname{arcsec}}
\providecommand{\arccsc}{\operatorname{arccsc}}\providecommand{\arccot}{\operatorname{arccot}}
\providecommand{\sgn}{\operatorname{sgn}}\providecommand{\lcm}{\operatorname{lcm}}
\providecommand{\Var}{\operatorname{Var}}\providecommand{\Cov}{\operatorname{Cov}}

%% FIGFIX-BEGIN v4 — correctifs globaux des figures (docs/onbuch-generation/tools/figfix)
\usepackage{adjustbox}\usepackage{etoolbox}
% toute figure TikZ/pgfplots plus large que la ligne est réduite à la largeur disponible
\BeforeBeginEnvironment{tikzpicture}{\begin{adjustbox}{max width=\linewidth}}
\AfterEndEnvironment{tikzpicture}{\end{adjustbox}}
% légendes pgfplots lisibles par-dessus les courbes
\pgfplotsset{every axis/.append style={legend style={fill=white,fill opacity=0.92,text opacity=1,draw=black!15,inner sep=3pt}}}
% étiquettes d'arêtes (syntaxe edge["..."]) avec fond blanc
\tikzset{every edge quotes/.append style={fill=white,inner sep=1.2pt}}
%% FIGFIX-END
\begin{document}

\pagegarde{Description of computer networking}{ICT}{Upper Sixth}{Upper Sixth — ICT}{2}{2 periods}{Ce chapitre explore la complexité inhérente aux réseaux informatiques modernes et les défis opérationnels qui en découlent, avant de proposer des activités d'intégration pour concevoir, déployer et superviser des solutions réseau robustes.
\vspace{4pt}
\begin{multicols}{2}\small
\begin{itemize}
\item Analyser les facteurs de complexité dans les architectures réseau (topologies, protocoles, adressage, virtualisation).
\item Évaluer les impacts des défis opérationnels : performance, disponibilité, sécurité et gestion de la configuration.
\item Appliquer des méthodologies de conception réseau (hiérarchique, modulaire, redondante) pour répondre à des besoins donnés.
\item Mettre en œuvre des outils de supervision, de diagnostic et de dépannage (SNMP, NetFlow, analyseurs de paquets).
\item Concevoir un plan de gestion des changements et de continuité de service (backup, restauration, documentation).
\item Réaliser une activité d'intégration complète : de l'étude de cas à la simulation d'un réseau d'entreprise camerounais.
\end{itemize}\end{multicols}}

\begin{sommaire}
\begin{multicols}{2}
\begin{enumerate}
\item 1. Facteurs de complexité dans les réseaux modernes
\item 2. Défis opérationnels : performance, disponibilité, sécurité, gestion
\item 3. Conception, déploiement et supervision d'un réseau d'entreprise
\item 4. Activités d'intégration : étude de cas, simulation et rapport
\item Integration activity
\item Methods and worked examples
\item Exercises
\item Solutions to the exercises
\item Summary sheet
\end{enumerate}
\end{multicols}
\end{sommaire}

\begin{objectifs}
\begin{itemize}
\item Analyser les facteurs de complexité dans les architectures réseau (topologies, protocoles, adressage, virtualisation).
\item Évaluer les impacts des défis opérationnels : performance, disponibilité, sécurité et gestion de la configuration.
\item Appliquer des méthodologies de conception réseau (hiérarchique, modulaire, redondante) pour répondre à des besoins donnés.
\item Mettre en œuvre des outils de supervision, de diagnostic et de dépannage (SNMP, NetFlow, analyseurs de paquets).
\item Concevoir un plan de gestion des changements et de continuité de service (backup, restauration, documentation).
\item Réaliser une activité d'intégration complète : de l'étude de cas à la simulation d'un réseau d'entreprise camerounais.
\item Rédiger un rapport technique structuré présentant l'architecture, les choix techniques et les résultats de tests.
\item Préparer les épreuves du GCE Advanced Level en résolvant des problèmes types sur la complexité et l'exploitation réseau.
\end{itemize}
\end{objectifs}

\begin{prerequis}
\begin{itemize}
\item Modèle OSI / TCP‑IP et encapsulation des données (Lower Sixth).
\item Adressage IPv4 / IPv6, sous‑réseaux, VLSM, CIDR.
\item Principes de base des LAN (Ethernet, commutateurs, VLAN) et WAN (routage statique/dynamique).
\item Notions élémentaires de sécurité réseau (pare‑feu, ACL, VPN).
\end{itemize}
\end{prerequis}

\newpage

\coursec{Factors of Complexity in Modern Networks}

In a cybercafé in Yaoundé, a connection outage cuts off access to a dozen workstations during an e-sports competition, revealing how the complexity of a local network can paralyse a daily activity. Students discover that behind every click lie protocols, equipment, and management strategies that determine the fluidity or the breakdown of service. This introductory situation highlights a fundamental truth: as networks grow in size, heterogeneity, and criticality, their complexity explodes, turning simple connectivity into a major engineering challenge. This section dissects the anatomical roots of that complexity.

\courssub{Diversity of Physical and Logical Topologies}

The topology of a network defines the geometric arrangement of links and nodes. In a modern enterprise, we rarely encounter a single pure topology; instead, we face a \cle{hybrid topology} combining several models to optimise cost, redundancy, and performance.

\begin{definition}[Network Topology]
The \textbf{physical topology} describes the actual cabling layout (cables, fibres, wireless links). The \textbf{logical topology} describes the path that data frames/packets follow, which may differ from the physical layout (e.g., a physical star wired to a switch but a logical bus due to a hub, or a logical ring over a physical mesh via MPLS).
\end{definition}

\begin{popfigure}
\begin{tikzpicture}[scale=0.8, every node/.style={font=\small}]
% --- Star Topology ---
\begin{scope}[xshift=0cm]
\node[draw, circle, fill=popblue!20, minimum size=1.2cm] (center) at (0,0) {Switch};
\foreach \angle/\lbl in {90/PC1, 150/PC2, 210/PC3, 270/Printer, 330/AP} {
    \node[draw, rectangle, fill=popgreen!20, minimum width=1cm, minimum height=0.6cm] (n) at (\angle:2.5) {\lbl};
    \draw[popline, thick] (center) -- (n) node[midway, sloped, above, font=\tiny] {UTP Cat6};
}
\node[align=center, font=\bfseries] at (0,-3.5) {Star Topology};
\node[align=center, font=\tiny, text width=3cm] at (0,-4.1) {Central point of failure.\par Easy to troubleshoot.};
\end{scope}

% --- Full Mesh Topology ---
\begin{scope}[xshift=7cm]
\foreach \i/\angle in {1/90, 2/162, 3/234, 4/306, 5/18} {
    \node[draw, circle, fill=poporange!20, minimum size=0.8cm] (n\i) at (\angle:2) {R\i};
}
\foreach \i in {1,...,5} {
    \pgfmathtruncatemacro{\j}{mod(\i,5)+1}
    \draw[popline, thick] (n\i) -- (n\j);
    \pgfmathtruncatemacro{\k}{mod(\i+1,5)+1}
    \draw[popline, thick, dashed] (n\i) -- (n\k);
}
\node[align=center, font=\bfseries] at (7,-3.5) {Full Mesh Topology};
\node[align=center, font=\tiny, text width=3cm] at (7,-4.1) {Maximum redundancy.\par $N(N-1)/2$ links $\to$ costly.};
\end{scope}

% --- Spine-Leaf Topology ---
\begin{scope}[xshift=14cm]
% Spine switches
\node[draw, rectangle, fill=poppurple!20, minimum width=1.5cm, minimum height=0.8cm] (s1) at (0,3) {Spine 1};
\node[draw, rectangle, fill=poppurple!20, minimum width=1.5cm, minimum height=0.8cm] (s2) at (3,3) {Spine 2};
% Leaf switches
\foreach \x/\lbl/\i in {0/L1/1, 1.5/L2/2, 3/L3/3, 4.5/L4/4} {
    \node[draw, rectangle, fill=poppink!20, minimum width=1.2cm, minimum height=0.7cm] (leaf\i) at (\x,0) {\lbl};
    \draw[popline, thick] (s1) -- (leaf\i);
    \draw[popline, thick] (s2) -- (leaf\i);
}
% Servers/Hosts under leaves
\foreach \x/\lbl/\i in {0/S1/1, 1.5/S2/2, 3/S3/3, 4.5/S4/4} {
    \node[draw, rectangle, fill=popgreen!10, minimum width=0.8cm, minimum height=0.5cm, font=\tiny] (srv\i) at (\x,-1.2) {\lbl};
    \draw[popline] (leaf\i) -- (srv\i);
}
\node[align=center, font=\bfseries] at (2.25,-3.5) {Spine-Leaf (Clos) Topology};
\node[align=center, font=\tiny, text width=3.5cm] at (2.25,-4.1) {Predictable latency (2 hops).\par Horizontal scalability.\par Standard in Data Centres.};
\end{scope}
\end{tikzpicture}
\legende{Comparison of three fundamental topologies: Star (Access layer), Full Mesh (Core redundancy), and Spine-Leaf (Modern Data Centre).}
\end{popfigure}

\textbf{Star Topology} dominates the \textbf{Access Layer} (end-user connection). It is cheap and simple but creates a single point of failure at the central switch.
\textbf{Full/Partial Mesh} appears in the \textbf{Core Layer} or WAN links where redundancy justifies the $O(N^2)$ cabling cost.
\textbf{Spine-Leaf (Clos Network)} is the modern standard for Data Centres (e.g., hosting MTN Mobile Money servers or government e-services platforms). Every Leaf connects to every Spine. This guarantees a constant \textbf{2-hop latency} between any two servers (Leaf $\to$ Spine $\to$ Leaf) and allows horizontal scaling by adding Leaf switches without rewiring the Spine.

\begin{attention}[Topology Mismatch]
A common design error is imposing a Star topology at the Core layer to save ports. This creates a bottleneck and a single point of failure for the entire organisation. Always match topology to the hierarchical layer: Access=Star, Distribution=Partial Mesh, Core=Full Mesh or Spine-Leaf.
\end{attention}

\courssub{Multiplicity of Routing and Switching Protocols}

A network is not just cables; it is a distributed operating system running dozens of protocols simultaneously. Complexity arises from the \textbf{interaction} between Layer 2 (Switching) and Layer 3 (Routing) protocols.

\begin{center}
\begin{tabularx}{\linewidth}{|l|X|X|X|}
\hline
\rowcolor{popblueL}
\textbf{Domain} & \textbf{Protocol} & \textbf{Type} & \textbf{Complexity Driver} \\
\hline
\textbf{Switching (L2)} & STP (802.1D) & Loop prevention & Slow convergence (30-50s), blocks links. \\
& RSTP (802.1w) & Fast convergence & Port roles/states complexity, edge port config. \\
& MSTP (802.1s) & Multi-instance & Mapping VLANs to instances (MSTI), CST interaction. \\
& VLAN (802.1Q) & Segmentation & Trunk negotiation (DTP), Native VLAN mismatch, VTP pruning. \\
\hline
\textbf{Routing (L3)} & OSPF & Link-State (IGP) & Area design (Backbone Area 0), LSA flooding, DR/BDR election, metric tuning. \\
& EIGRP & Distance Vector (Cisco) & Feasible Successor logic, Stuck-in-Active (SIA), manual summarization. \\
& BGP & Path Vector (EGP) & Policy-based routing, Attribute manipulation (Weight, Local Pref, AS-Path, MED), Route Reflectors, Confederations. \\
\hline
\end{tabularx}
\end{center}

\textbf{Interaction Example:} A link failure triggers STP/RSTP reconvergence (L2). If the link was a routed link (L3), OSPF detects the neighbour loss (Dead Timer) and recalculates SPF. If the L2 convergence is slower than L3 timers, routing adjacencies flap, causing \textbf{route flapping} and CPU spikes. This \cle{protocol interaction} is a primary source of "grey failures" — the network is up, but performance is degraded.

\begin{exemplebox}[OSPF Area Design in a Regional Administration]
Consider the \textbf{Ministry of Secondary Education (MINESEC)} connecting its Regional Delegations (10 regions) to the Central Administration in Yaoundé.
\begin{itemize}
\item \textbf{Area 0 (Backbone):} Yaoundé Core Routers (CR1, CR2).
\item \textbf{Area 1 (Centre):} Yaoundé Distribution.
\item \textbf{Area 2 (Littoral):} Douala Delegation.
\item \textbf{Area 3 (West):} Bafoussam Delegation.
\item ... up to Area 10 (Far North).
\end{itemize}
Each region is a \textbf{Stub Area} or \textbf{Totally Stubby Area} to prevent external routes (Internet, Cloud) from flooding the WAN links (often limited bandwidth VSAT or Fibre). Only a default route is injected by the ABR (Area Border Router). This reduces the LSDB (Link State Database) size on regional routers, saving RAM and CPU.
\end{exemplebox}

\begin{popfigure}
\begin{tikzpicture}
\begin{axis}[
    width=\linewidth, height=7cm,
    xlabel={Number of Routers in OSPF Domain ($N$)},
    ylabel={LSAs / SPF Complexity ($O(N \log N)$)},
    grid=major, grid style={popline},
    legend pos=north west,
    xmin=10, xmax=500,
    ymin=0, ymax=120000,
    tick label style={font=\small},
    label style={font=\small},
    legend style={font=\small}
]
\addplot[popblue, thick, mark=*] coordinates {
    (10, 500) (50, 8000) (100, 25000) (200, 70000) (300, 110000)
};
\addlegendentry{Estimated LSA Count (Full Mesh Area 0)}

\addplot[poporange, thick, mark=square*] coordinates {
    (10, 200) (50, 1500) (100, 4000) (200, 12000) (300, 22000)
};
\addlegendentry{With Hierarchical Areas (Sum of Areas)}

\addplot[popgreen, thick, dashed, mark=triangle*] coordinates {
    (10, 100) (50, 500) (100, 1200) (200, 3000) (300, 5500)
};
\addlegendentry{Stub/Totally Stubby Areas (Regional)}
\end{axis}
\end{tikzpicture}
\legende{Growth of OSPF Link State Advertisements (LSAs) and SPF calculation load vs. Network Size. Hierarchical design (Areas) and Stub areas drastically reduce complexity.}
\end{popfigure}

\begin{propriete}[Metcalfe's Law vs. Management Cost]
\textbf{Metcalfe's Law} states the value of a telecommunications network is proportional to the square of the number of connected users ($V \propto n^2$).
However, the \textbf{management complexity} (configuration states, interaction paths, troubleshooting scenarios) grows approximately as $C \propto n \log n$ for a well-designed hierarchical network, but can approach $O(n^2)$ or worse in a flat, poorly segmented network.
\textit{Examinable:} Be able to explain why hierarchical design (Core/Distribution/Access) changes the complexity scaling from quadratic to quasi-linear.
\end{propriete}

\courssub{Addressing Management: IPv4/IPv6 Transition, Dual-Stack, NAT, DHCPv6}

The exhaustion of IPv4 public addresses forces a coexistence strategy that significantly increases operational complexity.

\begin{definition}[Dual Stack]
\textbf{Dual Stack} means every network interface (router, server, PC, smartphone) runs IPv4 and IPv6 protocol stacks simultaneously. The OS uses \textbf{Happy Eyeballs (RFC 8305)} algorithm to prefer IPv6 but fallback quickly to IPv4 if latency/failure occurs.
\end{definition}

\begin{methode}[Deploying Dual-Stack in an Enterprise LAN]
\begin{enumerate}
\item \textbf{Address Plan:} Define IPv6 /48 prefix from ISP (e.g., AFRINIC allocation). Subnet into /64 per VLAN.
\item \textbf{Core/Distribution:} Enable IPv6 unicast routing (`ipv6 unicast-routing`). Configure OSPFv3 or BGP for IPv6 address families.
\item \textbf{Access Layer (DHCPv6 / SLAAC):}
    \item \textbf{SLAAC (Stateless):} Router Advertisement (RA) with `Managed Config Flag = 0`, `Other Config Flag = 0`. Hosts self-assign via EUI-64 or Privacy Extensions (RFC 4941).
    \item \textbf{DHCPv6 Stateful:} RA `Managed Config Flag = 1`. DHCPv6 server assigns addresses, DNS, NTP. Required for strict asset tracking (common in Cameroonian admin networks).
    \item \textbf{DHCPv6 Stateless:} RA `Managed Config Flag = 0`, `Other Config Flag = 1`. SLAAC for address, DHCPv6 for DNS/NTP only.
\item \textbf{Security:} Implement RA Guard, DHCPv6 Guard, ND Inspection on switches to prevent rogue RAs (a major attack vector in IPv6).
\item \textbf{NAT64/DNS64:} For IPv6-only hosts accessing IPv4-only legacy servers (still common in some gov. apps).
\end{enumerate}
\end{methode}

\begin{attention}[NAT44 vs NAT64 vs NPTv6]
\begin{itemize}
\item \textbf{NAT44 (IPv4 NAT):} Overload (PAT) hides thousands of private IPs behind one public IP. Breaks end-to-end principle, complicates VoIP/IPsec (requires ALG).
\item \textbf{NAT64:} Translates IPv6 $\to$ IPv4. Allows IPv6-only clients to reach IPv4 servers. Requires DNS64 to synthesize AAAA records from A records.
\item \textbf{NPTv6 (Network Prefix Translation, RFC 6296):} Stateless, 1:1 prefix translation (e.g., Internal ULA `fd00:1::/64` $\leftrightarrow$ External GUA `2001:db8:1::/64`). Preserves end-to-end reachability and port semantics. \textbf{Preferred for IPv6 multihoming} without BGP PI space.
\end{itemize}
\textbf{Trap:} Do not configure NAT66 (IPv6-to-IPv6 NAT) as a standard practice; it violates IPv6 design philosophy and breaks IPsec AH/ESP integrity checks. Use NPTv6 or Unique Local Addresses (ULA) internally with proper routing.
\end{attention}

\courssub{Virtualisation and Overlay: VXLAN, NVGRE, SDN, NFV}

Physical network hardware (boxes) is being abstracted into software. This \textbf{decoupling} introduces a new layer of complexity: the \cle{Overlay Network}.

\begin{definition}[Overlay vs Underlay]
The \textbf{Underlay} is the physical IP fabric (Spine-Leaf running BGP/OSPF) providing IP reachability between all nodes. The \textbf{Overlay} is a virtual topology (tunnels) built on top, decoupling tenant addressing (VNI) from physical addressing (VTEP IP).
\end{definition}

\begin{center}
\begin{tabularx}{\linewidth}{|l|X|X|}
\hline
\rowcolor{popblueL}
\textbf{Technology} & \textbf{Encapsulation} & \textbf{Key Characteristic} \\
\hline
\textbf{VXLAN (RFC 7348)} & UDP Header (Dst Port 4789) + VXLAN Header (24-bit VNI) & \textbf{Industry Standard}. 16M segments (VNI) vs 4k VLANs. Runs over any L3 fabric. Supported by Linux, Open vSwitch, Cisco, Juniper, Arista. \\
\textbf{NVGRE} & GRE Header + 24-bit VSID (Virtual Subnet ID) & Microsoft / Hyper-V centric. No UDP header $\to$ harder load balancing (ECMP) on L4 headers. \\
\textbf{GENEVE} & UDP + Extensible TLV Options & Designed to supersede VXLAN/NVGRE. Used by OVN/OVS, AWS, Azure. \\
\hline
\end{tabularx}
\end{center}

\textbf{SDN (Software-Defined Networking)} separates the \textbf{Control Plane} (Controller: OpenDaylight, ONOS, Cisco APIC, OpenStack Neutron) from the \textbf{Data Plane} (Switches: OpenFlow, P4 programmable ASICs).
\textbf{NFV (Network Functions Virtualization)} moves network appliances (Firewalls, Load Balancers, Routers, CGNAT, DPI) from proprietary hardware to \textbf{VNFs (Virtual Network Functions)} running on standard servers (COTS) orchestrated by MANO (MANagement and Orchestration - ETSI NFV MANO).

\begin{experience}[Visualising VXLAN Overlay in a Cameroonian Data Centre]
\textbf{Scenario:} CAMTEL hosts two government tenants: \texttt{MINSANTE} (Health) and \texttt{MINESEC} (Education) on the same physical Spine-Leaf fabric.
\begin{enumerate}
\item \textbf{Underlay:} Leaf switches run BGP EVPN (Ethernet VPN) as control plane for VXLAN. Spines are route reflectors.
\item \textbf{Tenant MINSANTE:} VNI 10100 (L2VNI) for VLAN 100 extension across sites. VRF \texttt{VRF-MINSANTE} with L3VNI 50001 for inter-VLAN routing.
\item \textbf{Tenant MINESEC:} VNI 20100, VRF \texttt{VRF-MINESEC}, L3VNI 50002.
\item \textbf{Isolation:} Despite sharing physical ports on Leaf 1, traffic from MINSANTE Server (VNI 10100) is encapsulated in UDP/4789 towards Leaf 2. MINESEC traffic uses VNI 20100. They never see each other's frames.
\item \textbf{Mobility:} A VM migrates from Leaf 1 to Leaf 3 (vMotion). The Controller updates the EVPN MAC/IP route (Type 2) advertising the new VTEP IP (Leaf 3). Traffic follows instantly. No STP convergence needed.
\end{enumerate}
\textbf{Complexity Cost:} Troubleshooting now requires checking: Underlay routing (BGP/OSPF) $\to$ VTEP reachability $\to$ EVPN Control Plane (Route Types 2, 3, 5) $\to$ VXLAN Encapsulation/Decapsulation counters $\to$ Tenant VRF routing.
\end{experience}

\courssub{Multi-Vendor Interoperability and Standards}

No modern network is single-vendor. A typical Cameroonian ISP or large enterprise (e.g., SONARA, SNH, University of Yaoundé I) mixes Cisco (Core), Huawei (Access/Metro), Fortinet (Security), Aruba (Wireless), and Linux servers (Monitoring).

\begin{definition}[Interoperability]
The ability of equipment from different vendors to work together correctly using open standards. It relies on strict adherence to \textbf{IEEE} (L1/L2: 802.3, 802.1Q, 802.1X, 802.11), \textbf{IETF} (L3+: BGP, OSPF, IS-IS, MPLS, VXLAN, BFD, SNMP, NETCONF/YANG), and \textbf{ITU-T} (G.8032 ERPS, OTN, PWE3) standards.
\end{definition}

\begin{aretenir}[Key Interoperability Checkpoints]
\begin{itemize}
\item \textbf{MTU/MSS:} VXLAN adds 50+ bytes. Ensure Underlay MTU $\ge$ 1600 (pref. 9000 Jumbo Frames). Configure `ip tcp adjust-mss` on L3 interfaces.
\item \textbf{ECMP Hashing:} Vendor A hashes on Src/Dst IP only; Vendor B includes L4 ports. Asymmetric hashing breaks stateful firewalls/load balancers. Configure symmetric hashing or consistent hashing algorithms.
\item \textbf{BGP Capabilities:} Ensure both ends negotiate \textbf{Multiprotocol Extensions (MP-BGP)} for IPv6, VPNv4, EVPN. Mismatch in `address-family` capability = session up but no routes exchanged.
\item \textbf{Time Synchronisation:} NTP/PTP (IEEE 1588v2) mandatory for log correlation, BFD, SAML/SSO, Kerberos. Use Stratum 1 GPS clocks (ANTIC provides time reference via \texttt{ntp.antic.cm}).
\item \textbf{Telemetry:} Move from SNMP polling (slow, heavy) to \textbf{Streaming Telemetry (gNMI/gRPC, NETCONF)} for real-time visibility.
\end{itemize}
\end{aretenir}

\courssub{Scalability: From Cybercafé LAN to Regional Administration WAN}

Complexity is not static; it scales with the organisation. Understanding the \textbf{scalability boundaries} of technologies is crucial for the GCE exam and real design.

\begin{center}
\begin{tabularx}{\linewidth}{|l|X|X|X|}
\hline
\rowcolor{popblueL}
\textbf{Scale} & \textbf{Typical Context (Cameroon)} & \textbf{Key Technologies} & \textbf{Complexity Drivers} \\
\hline
\textbf{Micro} & Cybercafé, Small Office (SOHO) & 1 Router (ISP CPE), 1 Switch (Unmanaged/Smart), WiFi AP. NAT, DHCPv4. & Low. Single broadcast domain. Flat topology. \\
\textbf{Small Enterprise} & Local NGO, School, SME & L2 Switches (VLANs), Router/Firewall (UTM), Wireless Controller (or Cloud-managed AP). OSPF Single Area. & VLAN design, Inter-VLAN routing (Router-on-Stick vs SVI), Basic Security (ACL, Port-Security). \\
\textbf{Medium Enterprise} & Regional Hospital, University Faculty, Bank Branch Network & Hierarchical (Core/Dist/Access). Stackwise/VSS/MC-LAG. OSPF Multi-Area / EIGRP. WAN: MPLS L3VPN or SD-WAN. & Route Summarization, Redistribution, QoS (Voice/Video), High Availability (HSRP/VRRP/GLBP), Wireless Roaming. \\
\textbf{Large / Carrier} & MTN/Orange Core, CAMTEL Backbone, MINPOSTEL GovNet & Spine-Leaf DC, MPLS Core (LDP/RSVP-TE/SR-MPLS), BGP (Internet Peering, VPNv4/6, EVPN), Segment Routing (SR-MPLS/SRv6). & Massive Scale (100k+ routes), Fast Reroute (FRR/TI-LFA), Traffic Engineering (SR-TE), Automation (Ansible, NetBox, Python/Go), Security (BGPsec, RPKI, FlowSpec). \\
\hline
\end{tabularx}
\end{center}

\begin{exoresolu}[Design Evolution: From Flat LAN to Hierarchical WAN]
\textbf{Statement:} A startup in Douala (``TechHub237'') grows from 20 to 500 employees across 3 buildings in 2 years. Their initial flat network (single /24, one L2 switch stack) suffers broadcast storms and IP conflicts. Propose the evolution steps.

\textbf{Detailed Solution:}
\begin{enumerate}
\item \textbf{Segmentation (L2):} Implement VLANs per department (Admin, Dev, Servers, Guests, IoT, Voice). VLAN 1 unused (shutdown). Native VLAN = 999 (dedicated, unused).
\item \textbf{Addressing (L3):} Migrate from 192.168.1.0/24 to RFC 1918 \texttt{10.0.0.0/16} supernet.
    \begin{itemize}
    \item Building A (Core): 10.0.0.0/20 (Core Links, Server Farm).
    \item Building B (Dist/Access): 10.0.16.0/20.
    \item Building C (Dist/Access): 10.0.32.0/20.
    \item WAN/VPN: 10.0.255.0/24.
    \end{itemize}
\item \textbf{Hierarchy:} Deploy \textbf{Collapsed Core} (Core + Distribution in same pair of switches, e.g., Cisco Catalyst 9500 / Huawei CloudEngine 6800) in Building A (MDF). Access switches in Building B \& C (IDFs) uplink via 10G/25G fibre (LACP EtherChannel) to Core.
\item \textbf{Routing:} OSPF Single Area (Area 0) initially. Core runs SVIs (L3 interfaces) for VLANs. Access switches remain L2 (or L3 Lite with static routes if budget tight).
\item \textbf{Redundancy:} Dual-homed Access switches (MC-LAG / vPC / M-LAG) to both Core switches. No STP blocking $\to$ all links active.
\item \textbf{Services:} DHCP Relay (`ip helper-address`) on Core SVIs pointing to central DHCP Server (Windows/Linux/Kea). DNS internal (AD integrated). NTP from Core.
\item \textbf{Security:} 802.1X (Dot1x) for wired + MAB for printers/IoT. Guest VLAN / Auth-Fail VLAN. Firewall (FortiGate/Palo Alto) at Internet edge with SSL Inspection.
\item \textbf{IPv6:} Request /48 from ISP (MTN/Orange Business). Deploy Dual-Stack on Core. SLAAC + Stateless DHCPv6 for workstations. Stateful DHCPv6 for Servers.
\item \textbf{Automation:} Config backup (Oxidized/NetBox), Push changes via Ansible. Monitoring: Prometheus + Grafana (SNMP Exporter / gNMI).
\end{enumerate}
\textbf{Result:} Broadcast domains limited to /24 (254 hosts). Failure domain limited to Access switch pair. Troubleshooting time reduced from hours to minutes. Ready for SD-WAN branch expansion.
\end{exoresolu}

\begin{savaistu}[Cameroon Context: ANTIC and the National Cybersecurity Strategy]
The \textbf{National Agency for Information and Communication Technologies (ANTIC)} mandates security standards for critical infrastructure (Decree 2011/... on cybersecurity). This forces network architects to integrate:
\begin{itemize}
\item \textbf{Network Segmentation:} Strict separation of IT (Information Technology) and OT (Operational Technology) networks in SONARA, CDC, AES-SONEL.
\item \textbf{Logging \& SIEM:} Centralised logs (Syslog/NetFlow/IPFIX) retained 12 months.
\item \textbf{Encryption:} IPsec/MACsec on inter-site links. TLS 1.3 for management planes (NETCONF/RESTCONF/HTTPS).
\item \textbf{IPv6 Readiness:} Government portals (\texttt{gov.cm}) must be IPv6 accessible (ANTIC directive).
\end{itemize}
Ignoring these regulatory constraints is a design failure, not just a technical oversight.
\end{savaistu}

\begin{exercice}[Complexity Analysis]
\begin{enumerate}
\item A network engineer configures OSPF Area 0 with 150 routers in a full mesh of adjacencies (via a Layer 2 fabric). Estimate the number of LSA Type 1 (Router LSAs) and the number of adjacencies. Why is this design discouraged?
\item Compare VXLAN and NVGRE regarding ECMP load balancing in the Underlay. Which field in the outer header allows VXLAN to hash flows better?
\item A company uses Dual-Stack. Windows 10 clients prefer IPv6 but the IPv6 path has 200ms RTT vs 20ms for IPv4. Explain the Happy Eyeballs mechanism (RFC 8305) behaviour in this scenario.
\item Define "Route Leaking" in an OSPF Multi-Area design. When is it necessary? What is the risk of creating routing loops?
\item In a Spine-Leaf fabric with 4 Spines and 48 Leaves, each Leaf has 48x 100G ports (32 Downlink to servers, 16 Uplink to Spines). Calculate the \textbf{oversubscription ratio} (Downlink bandwidth : Uplink bandwidth). Is it 1:1, 2:1, or 3:1?
\end{enumerate}
\end{exercice}

\begin{corrige}[Exercise 1]
\begin{enumerate}
\item \textbf{LSAs:} 150 Router LSAs (one per router). \textbf{Adjacencies:} Full mesh $\to N(N-1)/2 = 150 \times 149 / 2 = 11,175$ adjacencies. Each adjacency exchanges DBD/LSU/LSR. \textbf{Discouraged because:} SPF calculation runs on 150-node graph (heavy CPU), LSA flooding churn is massive on link flap, Single Point of Failure (Area 0 partition splits domain). \textbf{Solution:} Hierarchical Areas (Area 1, 2, 3...) connected to Area 0.
\item \textbf{VXLAN} uses UDP Destination Port 4789. The source port is a hash of the inner packet headers (flow entropy). This allows Underlay ECMP (which typically hashes on L4 Src/Dst Port) to distribute different VXLAN flows across different links. \textbf{NVGRE} uses GRE (Protocol 47) with no L4 port. Underlay ECMP hashing on IP headers only $\to$ all flows between same VTEP pair take same path $\to$ poor load balancing.
\item \textbf{Happy Eyeballs (RFC 8305):} Client initiates IPv6 and IPv4 connection attempts in parallel (or IPv6 first with short delay, e.g., 250ms). Whichever completes the TCP handshake first wins. Since IPv4 RTT (20ms) $\ll$ IPv6 RTT (200ms), the IPv4 SYN-ACK arrives first. The client abandons the IPv6 attempt and uses IPv4. This prevents user-perceived latency due to misconfigured/slow IPv6 paths.
\item \textbf{Route Leaking:} Redistributing routes from one OSPF Area (or Level 1 IS-IS) into another via an ABR (or L1/L2 router), typically using `summary-address` with `not-advertise` exceptions or `distribute-list` / `route-map`. \textbf{Necessary:} When a specific subnet in a Stub/Totally Stubby Area must be reachable by other areas (e.g., a shared server farm in a branch). \textbf{Risk:} If leaked back into the originating area with a better metric, or if two ABRs leak same routes without coordination $\to$ Routing Loops / Suboptimal routing. Use Route Tags to prevent re-advertisement.
\item \textbf{Downlink:} 32 $\times$ 100G = 3200 Gbps. \textbf{Uplink:} 16 $\times$ 100G = 1600 Gbps. \textbf{Ratio:} 3200 : 1600 = \textbf{2:1 (Oversubscribed)}. Non-blocking (1:1) would require 32 uplinks. 3:1 would be 10-11 uplinks. 2:1 is standard for general compute; 1:1 for storage/ML; 3:1 for web hosting.
\end{enumerate}
\end{corrige}

\coursec{Operational Challenges: Performance, Availability, Security, Management}

In the previous section, we saw that modern networks have become \cle{complex} systems: thousands of heterogeneous devices, virtualised services, mobile users, links to the cloud. This complexity is not merely a theoretical problem: it translates, every day, in a bank in Douala or a ministry in Yaoundé, into very concrete questions. Why does the videoconference stutter when the bandwidth seems sufficient? What happens if the main router fails on a Monday morning? How can we know that an intruder has been exploring the network for three weeks? Who modified the switch configuration last Thursday, and why?

These questions define the four great \cle{operational challenges} of running a network: \cle{performance}, \cle{availability}, \cle{security} and \cle{management} (in the sense of day-to-day administration and monitoring). This section studies them one by one, with their metrics, their mechanisms and their best practices.

\courssub{Measuring network performance}

\textbf{Observation.} Two companies pay for the same ``100 Mbit/s fibre'' subscription from an operator. In the first, VoIP calls are perfect; in the second, voices sound robotic and dropouts are frequent. The advertised bandwidth is identical: so bandwidth alone does not explain the difference. Network performance is \emph{multidimensional}.

\begin{definition}[Network performance quantities]
The performance of a network is measured mainly using four quantities:
\begin{itemize}
\item \cle{bandwidth}: the maximum amount of data transmissible per unit of time (in bit/s);
\item \cle{latency}: the one-way (or round-trip, RTT) delay of a packet between source and destination (in ms);
\item \cle{jitter}: the variation of latency from one packet to the next;
\item \cle{packet loss rate}: the proportion of emitted packets that never arrive (in \%).
\end{itemize}
\end{definition}

The intuition is that of a motorway: bandwidth is the number of lanes, latency is the travel time, jitter is the regularity of that time, and packet loss represents the vehicles that never arrive. A file-download application is mostly sensitive to bandwidth; VoIP and videoconferencing are mostly sensitive to \emph{latency}, \emph{jitter} and \emph{losses}, even with modest bandwidth.

\begin{exemplebox}[Orders of magnitude to know]
For voice over IP, one typically targets a one-way latency below $150\ \mathrm{ms}$, jitter below $30\ \mathrm{ms}$ and loss below $1\ \%$. Beyond that, conversation becomes painful. A geostationary satellite link already imposes roughly $250$ to $300\ \mathrm{ms}$ of one-way latency: VoIP over it is mediocre whatever the bandwidth.
\end{exemplebox}

\begin{definition}[Quality of Service (QoS)]
\cle{QoS} (\emph{Quality of Service}) denotes the set of mechanisms used to \textbf{prioritise critical traffic} and to guarantee differentiated performance levels according to the applications. The main families of mechanisms are:
\begin{itemize}
\item \cle{IntServ}: end-to-end resource reservation, flow by flow (RSVP protocol) --- precise but poorly scalable, rarely deployed at large scale;
\item \cle{DiffServ}: packets are marked (DSCP field) into service classes, and each router applies a per-class treatment (priority queues, guaranteed queues) --- this is the dominant model in enterprise networks;
\item \cle{policing}: the rate of a flow is measured and the excess is \emph{dropped} (or re-marked);
\item \cle{shaping}: traffic is \emph{smoothed} by delaying excess packets in a buffer, so as to respect a maximum rate without brutal losses.
\end{itemize}
\end{definition}

\begin{exemplebox}[QoS in a connected secondary school]
A school's network shares one Internet link between the computer room, the administration and the videoconference in the staff room. Without QoS, a massive download in the computer room saturates the link and degrades the videoconference. With DiffServ, voice/video is marked in the priority class (EF, \emph{Expedited Forwarding}) and downloads in the ``best effort'' class: under congestion, videoconference packets go first.
\end{exemplebox}

\begin{attention}[Confusing bandwidth with perceived quality]
Increasing bandwidth fixes neither latency, nor jitter, nor losses. Before buying extra capacity, one must \emph{measure} (for example with the ping utility for latency, or with jitter probes) to identify the real bottleneck. This is a classic --- and costly --- operational mistake.
\end{attention}

\courssub{High availability: never depend on a single device}

\textbf{Observation.} A mobile-money payment portal must work day and night. If a single router connects the site to the outside world, its failure cuts \emph{the whole} service. \cle{High availability} means designing the network so that no single failure interrupts the service: this is the principle of \cle{redundancy}.

\begin{definition}[Availability]
The \cle{availability} $A$ of a system is the proportion of time during which it actually delivers the expected service. It is expressed using two indicators:
\begin{itemize}
\item the \cle{MTBF} (\emph{Mean Time Between Failures}): the average time between two failures;
\item the \cle{MTTR} (\emph{Mean Time To Repair}): the average repair time (including detection, diagnosis and return to service).
\end{itemize}
\end{definition}

\begin{propriete}[Availability -- MTBF -- MTTR relationship]
For a system in steady state, availability is given by
\[ A = \frac{\mathrm{MTBF}}{\mathrm{MTBF} + \mathrm{MTTR}}. \]
This formula is \textbf{examinable}: it shows that availability can be improved in two ways --- increasing the MTBF (reliable equipment, redundancy) or \emph{reducing the MTTR} (monitoring, spare parts, procedures, saved configurations). The derivation is immediate: over a long cycle, the system is up for MTBF and down for MTTR, so the useful fraction of time is the ratio above.
\end{propriete}

\begin{exoresolu}[Computing availability and downtime]
A mail server of a company in Bonabéri fails on average 2 times per year, and each incident lasts on average 4 hours (detection included). Take a year of $8760$ hours.
\begin{enumerate}
\item Compute the MTBF and the MTTR.
\item Deduce the availability $A$.
\item What would be gained if the MTTR dropped to 30 minutes thanks to monitoring and saved configurations?
\end{enumerate}
\tcblower
\textbf{Solution.}
\begin{enumerate}
\item Two failures per year give an average cycle of $\dfrac{8760}{2} = 4380$ hours. The MTTR is $4$ h, so
\[ \mathrm{MTBF} = 4380 - 4 = 4376\ \mathrm{h}, \qquad \mathrm{MTTR} = 4\ \mathrm{h}. \]
\item Availability:
\[ A = \frac{4376}{4376 + 4} = \frac{4376}{4380} \approx 0.99909 \]
that is about $99.91\ \%$, i.e. about $8$ hours of downtime per year.
\item With $\mathrm{MTTR} = 0.5$ h, the cycle is $4380$ h and $\mathrm{MTBF} = 4379.5$ h:
\[ A' = \frac{4379.5}{4380} \approx 0.99989 \]
that is about $99.99\ \%$, i.e. about $1$ hour of downtime per year. Dividing the MTTR by 8 divided the annual downtime by 8: \textbf{reducing repair time is often the most cost-effective lever}.
\end{enumerate}
\end{exoresolu}

\begin{popfigure}
\begin{center}
\begin{tikzpicture}
\begin{axis}[
    width=0.86\linewidth, height=6.2cm,
    xlabel={Annual availability (\%)},
    ylabel={Annual downtime (hours)},
    xmin=99.0, xmax=99.999,
    ymode=log,
    xtick={99, 99.9, 99.99, 99.999},
    xticklabels={99, 99.9, 99.99, 99.999},
    ytick={0.0876, 0.876, 8.76, 87.6},
    yticklabels={0.088, 0.876, 8.76, 87.6},
    grid=both, grid style={popline, dashed},
    legend style={at={(0.03,0.03)}, anchor=south west, font=\small},
]
\addplot[popblue, very thick, mark=*, mark size=1.5pt] coordinates {
    (99, 87.6) (99.5, 43.8) (99.9, 8.76) (99.95, 4.38) (99.99, 0.876) (99.999, 0.0876)
};
\addlegendentry{Downtime $= 8760 \times (1-A)$}
\addplot[poporange, very thick, only marks, mark=square*, mark size=2.5pt] coordinates {
    (99.9, 8.76) (99.99, 0.876)
};
\node[poporange, font=\small, anchor=west] at (axis cs:99.9,8.76) {\ 99.9\,\% : $\approx$ 8 h 45/yr};
\node[poporange, font=\small, anchor=west] at (axis cs:99.99,0.876) {\ 99.99\,\% : $\approx$ 53 min/yr};
\end{axis}
\end{tikzpicture}
\end{center}
\legende{Impact of availability on annual downtime (8760-hour year, logarithmic scales). Moving from ``three nines'' to ``four nines'' divides downtime by ten: each additional ``nine'' costs more and more in redundancy.}
\end{popfigure}

\textbf{High-availability mechanisms.} Three techniques are on the syllabus:

\begin{itemize}
\item \cle{Gateway redundancy (VRRP, HSRP)}: two (or more) routers share a \emph{virtual IP address} used as the default gateway by the hosts. The ``active'' router answers; if it fails, the ``standby'' router takes over within seconds, without reconfiguring any host. HSRP is a Cisco proprietary protocol; VRRP is the open standard (RFC 5798).
\item \cle{Link aggregation (LACP)}: several physical links between two devices are bundled into a single logical link (IEEE 802.3ad standard). Bandwidth adds up and, if one cable is cut, traffic continues over the remaining links.
\item \cle{Fast convergence}: after a failure, routing and switching protocols must recompute paths \emph{quickly}. Fast-converging configurations are preferred (for example OSPF with tuned timers, RSTP rather than the historical STP which could take 30 to 50 seconds) to limit the interruption to a few seconds or less.
\end{itemize}

\begin{aretenir}[High availability: the essentials]
High availability rests on three pillars: \textbf{eliminating single points of failure} (device and link redundancy), \textbf{converging quickly} after a failure, and \textbf{reducing the MTTR} (monitoring, documentation, backups). Availability is computed as $A = \dfrac{\mathrm{MTBF}}{\mathrm{MTBF}+\mathrm{MTTR}}$.
\end{aretenir}

\courssub{Operational security: protecting the network every day}

Security is not limited to installing a firewall. It is a \emph{daily discipline}: partition, harden, log.

\textbf{Segmentation: DMZ and micro-segmentation.} Never place all devices in the same flat network. The \cle{DMZ} (\emph{demilitarized zone}) is a buffer subnetwork, placed between the Internet and the internal network, hosting the servers exposed to the public (website, e-mail). Thus, if a web server is compromised, the attacker is not yet \emph{inside} the internal network. \cle{Micro-segmentation} pushes the idea further: the internal network is divided into small zones (VLANs, internal firewalls) so that each service can communicate only with what it strictly needs. A compromised machine can then no longer move freely (``lateral movement'').

\begin{methode}[The ``Zero Trust'' approach applied to the internal network]
The \cle{Zero Trust} model starts from the principle that \emph{no} device or user is trusted by default, even inside the network. Typical implementation:
\begin{enumerate}
\item \textbf{Inventory} the legitimate devices, applications and flows of the network.
\item \textbf{Micro-segment}: create zones by function (management, accounting, guests, servers) with filtering rules between zones.
\item \textbf{Authenticate systematically}: every access (user or machine) is verified, for example via 802.1X on switch ports.
\item \textbf{Apply least privilege}: allow only the strictly necessary flows (precise ports and destinations).
\item \textbf{Verify continuously}: log all flows, analyse abnormal behaviour and revoke access as soon as doubt arises.
\end{enumerate}
\end{methode}

\textbf{Device hardening.} A router or switch delivered ``out of the box'' is vulnerable: default passwords, useless services enabled (Telnet, HTTP, SNMP v1 in clear text). \cle{Hardening} consists in disabling everything unnecessary, changing default credentials, enforcing SSH instead of Telnet, restricting administration to a dedicated management network and keeping firmware up to date.

\textbf{Centralised logging.} Every device generates \cle{logs}: connections, errors, link state changes. If these logs stay on each device, they are unusable --- and an attacker can erase them. They are therefore sent in real time to a central server via the \cle{Syslog} protocol. Above that, a \cle{SIEM} platform (\emph{Security Information and Event Management}) aggregates, correlates and analyses the logs of the whole estate to detect incidents (for example: ``200 failed logins in 2 minutes on the accounting server'').

\begin{savaistu}[Security monitoring in Cameroon]
In Cameroon, \cle{ANTIC} (the National Agency for Information and Communication Technologies) watches over the security of the State's information systems and raises awareness among companies about good practices: logging, access management, backups. Large operators such as MTN and Orange Cameroon run security operations centres (SOC) that apply precisely the principles of this section: segmentation, SIEM and continuous monitoring.
\end{savaistu}

\courssub{Configuration management: mastering change}

\textbf{Observation.} A technician modifies ``just one line'' on a switch on a Friday evening. On Monday, a whole building is out of service, and nobody knows exactly what was changed. \cle{Configuration management} aims to make every modification \emph{traced, reversible and reproducible}.

\begin{itemize}
\item \cle{Versioning (Git)}: device configuration files are stored in a Git repository. Every modification is a dated and signed \emph{commit}; two versions can be compared (\emph{diff}) and a rollback takes seconds.
\item \cle{Automation (Ansible, Netmiko)}: instead of connecting by hand to 50 switches, the desired configuration is described in a file (an Ansible \emph{playbook}) or a Python script (the \cle{Netmiko} library, based on SSH), and the tool applies it identically everywhere. Fewer human errors, fast deployment, repeatability.
\item \cle{Compliance (CIS benchmarks)}: the \emph{CIS Benchmarks} are public references of secure configuration good practices, device by device. Tools automatically check that real configurations remain compliant with these references and flag any drift.
\end{itemize}

\begin{attention}[The tragedy of the unsaved configuration]
Neglecting configuration backups causes a critical loss of time in case of failure: a router replaced ``blank'' must be reconfigured from memory, under pressure, with hours of interruption --- the MTTR explodes. Minimal good practice: automatic daily backup of configurations to a server (or a Git repository), and regular testing of the \emph{restore}. A backup never tested is not a backup.
\end{attention}

\courssub{Proactive monitoring: seeing the problem before the user}

\cle{Monitoring} (supervision) consists in continuously collecting measurements about the network in order to detect anomalies \emph{before} users complain.

\begin{itemize}
\item \cle{SNMP v3} (\emph{Simple Network Management Protocol}): the monitoring station periodically polls the devices (CPU load, per-interface traffic, temperature). Version 3 adds authentication and encryption --- versions 1 and 2c, in clear text, must be banned.
\item \cle{NetFlow / IPFIX}: routers export statistics about \emph{flows} (who talks to whom, on which port, what volume). This is the ideal tool to identify what saturates a link --- for example discovering that an infected host is massively sending data outside.
\item \cle{Streaming telemetry (gNMI)}: instead of polling devices every 5 minutes, the devices \emph{push} their measurements continuously (every few seconds) to the collector. Much finer and more reactive monitoring.
\item \cle{Alerting (Prometheus / Alertmanager)}: Prometheus stores metrics as time series and evaluates rules; Alertmanager sends alerts (e-mail, SMS, team messaging) when a threshold is crossed, for example ``WAN interface at 95 \% for 10 minutes''.
\end{itemize}

\courssub{Incident management: the ITIL method}

Even well designed, a network experiences incidents. The \cle{ITIL} framework (\emph{Information Technology Infrastructure Library}) structures their handling by distinguishing three processes:

\begin{itemize}
\item \cle{Incident management}: restore service \emph{as fast as possible} (a workaround is enough). Example: switch over to the backup 4G link.
\item \cle{Problem management}: find and eliminate the \emph{root cause} to avoid recurrence. Example: the main fibre link is cut every month by roadworks --- the duct must be relocated.
\item \cle{Change management}: every modification (update, new VLAN) is planned, assessed (risks, rollback plan) and approved \emph{before} being applied.
\end{itemize}

Two tools complete the approach: the \cle{post-mortem}, a blameless written report produced after every major incident (timeline, root cause, corrective actions), and \cle{runbooks}, step-by-step procedures written in advance for known incidents (``if the CAMTEL link goes down, do this''), which drastically reduce the MTTR.

\begin{popfigure}
\begin{center}
\begin{tikzpicture}[
    node distance=0.55cm and 0.9cm,
    etape/.style={draw=popblue, fill=popblueL, rounded corners=3pt, thick, minimum width=2.6cm, minimum height=1.05cm, align=center, font=\small},
    note/.style={font=\scriptsize\itshape, text=popmuted, align=center},
    fleche/.style={-{Stealth[length=3mm]}, very thick, popdark}
]
\node[etape, align=center] (det){\textbf{1. Detection}\\ monitoring, alert,\\ user call};
\node[etape, right=of det, align=center] (diag){\textbf{2. Diagnosis}\\ qualification,\\ localisation, cause};
\node[etape, right=of diag, align=center] (res){\textbf{3. Resolution}\\ fix or\\ workaround};
\node[etape, right=of res, align=center] (clo){\textbf{4. Closure}\\ validation, post-mortem,\\ runbook update};
\draw[fleche] (det) -- (diag);
\draw[fleche] (diag) -- (res);
\draw[fleche] (res) -- (clo);
\draw[fleche, poporange, dashed] (diag.south) .. controls +(0,-1.1) and +(0,-1.1) .. node[note, below]{root cause unknown: escalate / open a \emph{problem}} (det.south);
\draw[fleche, popgreen, dashed] (clo.north) .. controls +(0,0.9) and +(0,0.9) .. node[note, above]{feedback: improve monitoring} (det.north);
\end{tikzpicture}
\end{center}
\legende{Life cycle of an incident in the ITIL logic: detection, diagnosis, resolution, closure. The dashed loops show that diagnosis may lead to opening a \emph{problem}, and that closure feeds continuous improvement.}
\end{popfigure}

\begin{exercice}[Availability and architecture choice]
A company in Douala runs an e-commerce website. Its single access router has an MTBF of $2920$ h (three failures per year on average) and an MTTR of $6$ h.
\begin{enumerate}
\item Compute the current availability and the corresponding annual downtime (8760-hour year).
\item The director hesitates between buying a second router (VRRP redundancy) or hiring an on-call technician who would bring the MTTR down to $1$ h. Compute the availability in the second case and comment.
\item Name two elements, other than hardware, that reduce the MTTR.
\end{enumerate}
\end{exercice}

\begin{corrige}[Exercise: Availability and architecture choice]
\begin{enumerate}
\item $A = \dfrac{2920}{2920 + 6} = \dfrac{2920}{2926} \approx 0.99795$, i.e. about $99.8\ \%$. Annual downtime: $8760 \times (1 - 0.99795) \approx 18$ h per year (three 6-hour failures).
\item With $\mathrm{MTTR} = 1$ h: $A' = \dfrac{2920}{2921} \approx 0.99966$, i.e. about $99.97\ \%$, that is only $3$ h of downtime per year. Reducing the MTTR divides downtime by 6 without changing the hardware: it is often the best cost/effectiveness ratio. VRRP redundancy remains superior (failover takes seconds), but it is more expensive; ideally both are combined.
\item For example: monitoring with alerting (immediate detection), saved and versioned configurations (fast restore), runbooks, spare parts on site, trained staff.
\end{enumerate}
\end{corrige}

\begin{aretenir}[Key points of this section]
\begin{itemize}
\item Performance is measured by bandwidth, latency, jitter and loss; QoS (DiffServ, policing, shaping) protects critical traffic.
\item Availability is $A = \dfrac{\mathrm{MTBF}}{\mathrm{MTBF}+\mathrm{MTTR}}$; it is improved by redundancy (VRRP/HSRP, LACP) and above all by reducing the MTTR.
\item Operational security rests on segmentation (DMZ, micro-segmentation, Zero Trust), hardening and centralised logging (Syslog, SIEM).
\item Configurations are versioned (Git), deployed automatically (Ansible, Netmiko) and checked against references (CIS).
\item Proactive monitoring (SNMP v3, NetFlow/IPFIX, gNMI, Prometheus) detects before the user; ITIL structures the handling of incidents, problems and changes.
\end{itemize}
\end{aretenir}

\coursec{Conception, déploiement et supervision d'un réseau d'entreprise}

In the previous section we identified the operational challenges that any serious network must face: performance, availability, security and manageability. Knowing the problems is not enough — the network engineer must now \emph{build} an infrastructure that answers them. This section takes you through the complete life cycle of an enterprise network: from the first interview with the client to the moment a Grafana dashboard lights up green in the operations room. We will follow a concrete Cameroonian case study throughout: \textbf{CamAgro SARL}, an agro-industrial company with its headquarters in Yaoundé and two branches in Douala and Bafoussam.

\courssub{Gathering the requirements (recueil des besoins)}

Every successful network begins not with cables, but with \textbf{questions}. A network designed without a requirements study is like a house built without a plan: it may stand, but it will never fit the family.

\begin{methode}[Conducting the requirements study]
\begin{enumerate}
\item \textbf{Count the sites and users.} How many buildings, branches, remote workers? For CamAgro: 1 headquarters (Yaoundé, 180 users), 2 branches (Douala: 60 users, Bafoussam: 35 users).
\item \textbf{Identify the critical applications.} Which applications must never stop? For CamAgro: the ERP (stock and accounting), VoIP telephony, video surveillance, and mobile-money payment terminals at the Douala depot.
\item \textbf{Estimate the expected growth.} A network is designed for 3 to 5 years. CamAgro plans to double the Douala workforce within 3 years: the design must reserve addressing space and switch capacity for this.
\item \textbf{Fix the budget.} The budget constrains every technical choice: redundant links, high-availability equipment, licences for supervision tools. A good engineer always proposes at least two scenarios (essential / recommended).
\item \textbf{Write it all down.} The requirements document (\emph{cahier des charges}) is signed by the client; it is the reference against which the final network will be validated.
\end{enumerate}
\end{methode}

\begin{attention}[The forgotten requirement]
The most expensive sentence in networking is: ``We did not think of that.'' Forgetting the growth forecast leads to renumbering the whole IP plan two years later; forgetting a critical application (e.g.\ the CCTV cameras) leads to congestion that no budget line can fix quickly. Always validate the requirements document with \emph{every} department, not only the IT team.
\end{attention}

\courssub{Choosing the architecture}

Once the needs are known, we choose a \cle{topology model}. Three architectures dominate modern enterprise design. The choice depends on the traffic pattern and the scale of each location.

\begin{center}
\begin{tabularx}{\linewidth}{|l|X|X|X|}
\hline
\rowcolor{popblueL} \textbf{Architecture} & \textbf{Typical use case} & \textbf{Traffic pattern} & \textbf{Scaling method} \\
\hline
Three-tier (Core / Distribution / Access) & Campus, headquarters & North–South (user $\leftrightarrow$ server) & Larger chassis or more distribution blocks \\
\hline
Spine-leaf & Data centre, server farm & East–West (server $\leftrightarrow$ server) & Add leaf switches, connect to all spines \\
\hline
SD-WAN & Inter-site WAN (branch offices) & Mixed, policy-driven & Add transport links (fibre, 4G, satellite) \\
\hline
\end{tabularx}
\end{center}

\textbf{The three-tier model.} The classic campus design organises switches into three layers: the \cle{access layer} (where users, phones and printers connect), the \cle{distribution layer} (which aggregates access switches and applies policies such as filtering, QoS and routing between VLANs), and the \cle{core layer} (a fast, simple backbone whose only job is to move packets quickly between distribution blocks). Its great virtue is \emph{predictability}: each layer has one role, so a failure is easy to locate.

\textbf{Spine-leaf (data centres).} In a data centre, most traffic flows \emph{server to server} (east-west), not server to user. The \cle{spine-leaf} architecture connects every \emph{leaf} switch (attached to servers) to every \emph{spine} switch, so any server reaches any other in exactly two hops, with constant latency and massive horizontal scalability: to grow, you add a leaf, not a bigger chassis.

\textbf{SD-WAN (inter-site links).} Connecting Yaoundé, Douala and Bafoussam with traditional MPLS lines leased from a single operator is expensive and rigid. \cle{SD-WAN} (Software-Defined Wide Area Network) solves this.

\begin{definition}[SD-WAN]
SD-WAN is a wide-area networking approach that \textbf{separates the control plane from the data plane}. The \emph{data plane} (the routers at each site) simply forwards encrypted traffic over \emph{any} available transport — MPLS, fibre, 4G/LTE from MTN or Orange. The \emph{control plane} (routing decisions, security policies, path selection) is centralised in an \textbf{orchestrator} managed from a single console. A policy written once — ``VoIP always takes the lowest-latency link'' — is pushed automatically to all sites.
\end{definition}

\begin{savaistu}[SD-WAN in Cameroon]
Because leased MPLS lines remain costly relative to 4G/fibre Internet access in Cameroon, several banks and microfinance institutions now combine a CAMTEL fibre link with an MTN or Orange 4G backup under an SD-WAN overlay: if the fibre is cut (roadworks are a classic cause), traffic switches to 4G in seconds, often without users noticing.
\end{savaistu}

For CamAgro we therefore retain: a \textbf{three-tier} design at the headquarters, a \textbf{spine-leaf} fabric in the small Yaoundé server room, and \textbf{SD-WAN} to interconnect the three sites.

\begin{popfigure}
\begin{center}
\begin{tikzpicture}[scale=0.92, every node/.style={transform shape},
 site/.style={draw=popblue, thick, rounded corners, fill=popblueL, minimum width=2.6cm, minimum height=1.1cm, align=center, font=\small\bfseries},
 dev/.style={draw=popdark, rounded corners, fill=white, minimum width=1.5cm, minimum height=0.55cm, font=\scriptsize, align=center},
 cloud/.style={draw=poporange, very thick, dashed, ellipse, fill=orange!8, minimum width=3.4cm, minimum height=1.5cm, align=center, font=\small\bfseries},
 link/.style={thick, popdark},
 mpls/.style={very thick, poppurple},
 sdwan/.style={very thick, popgreen}]
% Cloud
\node[cloud, align=center] (wan) at (5.4,3.4){WAN cloud\\ \scriptsize MPLS + 4G/SD-WAN};
% Sites
\node[site, align=center] (yde) at (0,0){Yaoundé HQ\\ \scriptsize 180 users};
\node[site, align=center] (dla) at (5.4,0){Douala\\ \scriptsize 60 users};
\node[site, align=center] (baf) at (10.8,0){Bafoussam\\ \scriptsize 35 users};
% Edge devices
\node[dev] (ey) at (0,1.6) {Edge router + NGFW};
\node[dev] (ed) at (5.4,1.6) {Edge router + NGFW};
\node[dev] (eb) at (10.8,1.6) {Edge router + NGFW};
% Links
\draw[link] (yde) -- (ey);
\draw[link] (dla) -- (ed);
\draw[link] (baf) -- (eb);
\draw[mpls] (ey) -- (wan);
\draw[mpls] (ed) -- (wan);
\draw[sdwan] (eb) -- (wan);
% Inside HQ
\node[dev] (core) at (0,-1.3) {Core (2 switches)};
\node[dev] (dist) at (0,-2.4) {Distribution};
\node[dev] (acc1) at (-1.9,-3.5) {Access A};
\node[dev] (acc2) at (0,-3.5) {Access B};
\node[dev] (acc3) at (1.9,-3.5) {Access C};
\draw[link] (yde) -- (core);
\draw[link] (core) -- (dist);
\draw[link] (dist) -- (acc1);
\draw[link] (dist) -- (acc2);
\draw[link] (dist) -- (acc3);
% Legend
\node[font=\scriptsize, anchor=west] at (7.6,-2.6) {\textcolor{poppurple}{\rule{0.5cm}{2pt}} MPLS primary};
\node[font=\scriptsize, anchor=west] at (7.6,-3.2) {\textcolor{popgreen}{\rule{0.5cm}{2pt}} SD-WAN / 4G backup};
\end{tikzpicture}
\end{center}
\legende{CamAgro enterprise network: three-tier campus at Yaoundé HQ, SD-WAN overlay interconnecting the three sites over MPLS with 4G backup.}
\end{popfigure}

\courssub{IP addressing plan, VLAN plan and routing policy}

A coherent addressing plan is the \emph{identity card} of the network: a well-chosen address tells you immediately which site, which service and which VLAN a machine belongs to.

\textbf{VLSM in practice.} With \cle{VLSM} (Variable Length Subnet Masking) we carve the private block \texttt{10.0.0.0/8} hierarchically: one /16 per site, then subnets sized to real needs. CamAgro's plan:

\begin{center}
\begin{tabularx}{\linewidth}{|l|l|X|l|}
\hline
\rowcolor{popblueL} \textbf{Site / VLAN} & \textbf{Network} & \textbf{Purpose} & \textbf{Hosts} \\
\hline
Yaoundé site block & 10.10.0.0/16 & All HQ subnets & 65\,534 \\
\hline
\quad VLAN 10 & 10.10.10.0/24 & Staff workstations & 254 \\
\hline
\quad VLAN 20 & 10.10.20.0/24 & VoIP phones & 254 \\
\hline
\quad VLAN 30 & 10.10.30.0/27 & Servers (ERP, files) & 30 \\
\hline
\quad VLAN 99 & 10.10.99.0/29 & Management (OOB) & 6 \\
\hline
Douala site block & 10.20.0.0/16 & All Douala subnets & 65\,534 \\
\hline
Bafoussam site block & 10.30.0.0/16 & All Bafoussam subnets & 65\,534 \\
\hline
WAN links & 10.255.0.0/24 & Point-to-point /30 links & 2 each \\
\hline
\end{tabularx}
\end{center}

Notice the logic: the \textbf{second octet identifies the site} (10, 20, 30), the \textbf{third octet identifies the VLAN}. Reading \texttt{10.20.20.15} instantly tells the administrator: Douala, VoIP VLAN.

\textbf{IPv6.} The plan also reserves IPv6: \cle{GUA} (Global Unicast Addresses, obtained from the provider) for Internet-facing services, and \cle{ULA} (Unique Local Addresses, prefix \texttt{fd00::/8}, e.g.\ \texttt{fd10:10::/48} for Yaoundé) for internal addressing that must never be routed on the Internet — the IPv6 equivalent of RFC 1918 private space.

\textbf{Routing policy and summarisation.} Because each site owns a full /16, the WAN routers advertise a \emph{single summary route} per site: Douala announces \texttt{10.20.0.0/16}, not dozens of /24s. \cle{Route summarisation} keeps routing tables small, speeds up convergence and hides internal topology changes from other sites.

\begin{exoresolu}[Sizing a subnet with VLSM]
The CamAgro server room in Yaoundé must host at most 25 servers. Choose the smallest suitable subnet inside the block 10.10.0.0/16 and give its network address, mask, broadcast address and usable range.
\tcblower
We need 25 host addresses. With a prefix /$n$, the number of usable addresses is $2^{32-n} - 2$.
\begin{align*}
n = 27 &\Rightarrow 2^{5} - 2 = 30 \text{ usable addresses} \geq 25 \quad \checkmark \\
n = 28 &\Rightarrow 2^{4} - 2 = 14 \text{ usable addresses} < 25 \quad \text{(too small)}
\end{align*}
The smallest suitable subnet is therefore a \textbf{/27} (mask 255.255.255.224), with block size $256 - 224 = 32$. Choosing the third octet 30 for the server VLAN:
\begin{itemize}
\item Network address: \textbf{10.10.30.0/27}
\item Usable range: \textbf{10.10.30.1 -- 10.10.30.30}
\item Broadcast: \textbf{10.10.30.31}
\end{itemize}
30 usable addresses cover the 25 servers with 5 spare for growth.
\end{exoresolu}

\courssub{Selecting the equipment}

Equipment selection follows the architecture, never the reverse.

\begin{itemize}
\item \textbf{Core switches:} high-throughput, redundant (a pair, stacked or clustered), with 10/25 GbE uplinks. The core does \emph{no} filtering: speed only.
\item \textbf{Distribution switches:} Layer-3 capable (they route between VLANs), with ACL and QoS support.
\item \textbf{Access switches:} PoE+ to power phones and cameras, port security, 1 GbE to the desktop.
\item \textbf{Edge routers:} one per site, SD-WAN capable, with two WAN interfaces (fibre + 4G).
\item \textbf{Next-generation firewall (NGFW):} unlike a classic firewall that filters only ports and addresses, an \cle{NGFW} inspects the \emph{application} layer (it can block Facebook but allow WhatsApp Business), integrates intrusion prevention (IPS), and decrypts TLS traffic for inspection under a controlled policy.
\end{itemize}

\begin{propriete}[Least privilege on management interfaces (proof not examinable)]
Every management interface (SSH, HTTPS console, SNMP) must obey the \textbf{principle of least privilege}: it is reachable \emph{only} from the hosts that genuinely need it. In practice: (1) management traffic travels on a dedicated \textbf{out-of-band} VLAN (VLAN 99, 10.10.99.0/29), physically or logically separated from user traffic; (2) an \textbf{ACL} on each device accepts management connections only from the two administrator workstations; (3) default accounts are disabled and authentication is centralised (AAA). Consequence: even if an attacker compromises a user PC in VLAN 10, the management plane remains unreachable.
\end{propriete}

\courssub{Automating the deployment with Ansible}

Configuring 20 switches and 3 routers by hand, one SSH session at a time, is slow and — worse — \emph{irreproducible}: two engineers will never type exactly the same configuration. \cle{Automation} solves this. \textbf{Ansible} is an agentless automation tool: the engineer describes the desired configuration once, in a text file called a \emph{playbook} (YAML format), and Ansible pushes it over SSH to all devices, identically, in minutes.

\begin{exemplebox}[An Ansible playbook for base configuration]
A simplified playbook that configures hostnames, VLANs 10 and 20, an OSPF routing process and AAA authentication on all access switches:
\begin{itemize}
\item \texttt{- name: Set hostname} $\rightarrow$ \texttt{ios\_config: lines: hostname ACC-YDE-01}
\item \texttt{- name: Create VLANs} $\rightarrow$ \texttt{ios\_vlans: vlan\_id: 10, name: STAFF}
\item \texttt{- name: Enable OSPF} $\rightarrow$ \texttt{network 10.10.0.0 0.0.255.255 area 0}
\item \texttt{- name: Configure AAA} $\rightarrow$ \texttt{aaa authentication login default group tacacs+ local}
\end{itemize}
Running \texttt{ansible-playbook base-config.yml} applies these four tasks to the whole inventory. If a switch is re-installed after a failure, the \emph{same} playbook restores its exact configuration — this is called \textbf{idempotence}.
\end{exemplebox}

\begin{attention}[Automation amplifies errors]
Automation applies a mistake as faithfully as a correction — to \emph{all} devices at once. A wrong VLAN number in a playbook can cut off an entire site in seconds. Rule: always test a playbook on \textbf{one} pilot device, verify, then deploy to the fleet, and keep playbooks under version control (Git) so any change can be rolled back.
\end{attention}

\courssub{Supervision: dashboards, thresholds and event correlation}

A network that is not supervised is a network whose failures are discovered by angry phone calls. Supervision rests on three pillars.

\textbf{1. Metrics collection and dashboards.} Agents and SNMP polling feed a time-series database (e.g.\ Prometheus), and \cle{Grafana} turns the data into live dashboards: CPU and memory of each device, bandwidth per link, interface error counters, availability per site.

\begin{popfigure}
\begin{center}
\begin{tikzpicture}[scale=0.95, every node/.style={transform shape}]
% Window frame
\draw[fill=popdark, rounded corners] (0,0) rectangle (12.6,7.2);
\draw[fill=white] (0.15,0.15) rectangle (12.45,6.55);
\node[font=\small\bfseries, text=white, anchor=west] at (0.3,6.9) {Grafana --- CamAgro Network Overview};
\node[font=\scriptsize, text=popgreenL, anchor=east] at (12.3,6.9) {All systems operational};
% Panel 1: CPU gauge
\draw[fill=popblueL, rounded corners] (0.35,4.35) rectangle (4.35,6.35);
\node[font=\scriptsize\bfseries, anchor=west] at (0.5,6.1) {Core CPU load};
\draw[thick, popline] (2.35,5.0) arc[start angle=180, end angle=0, radius=1.1];
\draw[very thick, popgreen] (2.35,5.0) arc[start angle=180, end angle=115, radius=1.1];
\draw[very thick, poporange] (1.25,5.0) -- (2.05,5.55);
\node[font=\small\bfseries] at (2.35,4.6) {38\%};
% Panel 2: bandwidth graph
\draw[fill=white, draw=popline, rounded corners] (4.6,4.35) rectangle (12.25,6.35);
\node[font=\scriptsize\bfseries, anchor=west] at (4.75,6.1) {WAN bandwidth Yaoundé--Douala (Mbit/s)};
\draw[popline] (4.9,4.6) -- (12.0,4.6);
\draw[popline] (4.9,4.6) -- (4.9,5.95);
\draw[very thick, popblue] (4.9,4.8) -- (6.0,5.0) -- (7.0,4.9) -- (8.0,5.4) -- (9.0,5.3) -- (10.0,5.7) -- (11.0,5.5) -- (12.0,5.8);
\draw[very thick, popgreen, dashed] (4.9,4.7) -- (6.0,4.75) -- (7.0,4.72) -- (8.0,4.8) -- (9.0,4.78) -- (10.0,4.85) -- (11.0,4.82) -- (12.0,4.9);
\node[font=\tiny, anchor=west, text=popblue] at (9.6,5.95) {MPLS in};
\node[font=\tiny, anchor=west, text=popgreen] at (9.6,4.45) {4G backup};
% Panel 3: interface errors
\draw[fill=white, draw=popline, rounded corners] (0.35,0.35) rectangle (6.3,4.1);
\node[font=\scriptsize\bfseries, anchor=west] at (0.5,3.85) {Interface errors (per hour)};
\foreach \x/\h/\c in {0.9/0.3/popgreen, 1.7/0.5/popgreen, 2.5/0.4/popgreen, 3.3/2.6/poporange, 4.1/0.6/popgreen, 4.9/0.4/popgreen, 5.7/0.5/popgreen}
  \draw[fill=\c] (\x,0.6) rectangle (\x+0.55,0.6+\h);
\draw[popline] (0.7,0.6) -- (6.1,0.6);
\node[font=\tiny, text=poporange] at (3.6,3.45) {spike 14:00 --- Gi0/2 Douala};
% Panel 4: site availability
\draw[fill=white, draw=popline, rounded corners] (6.55,0.35) rectangle (12.25,4.1);
\node[font=\scriptsize\bfseries, anchor=west] at (6.7,3.85) {Site availability (30 days)};
\node[font=\scriptsize, anchor=west] at (6.8,3.2) {Yaoundé};
\draw[fill=popgreen] (8.6,3.1) rectangle (12.0,3.4);
\node[font=\tiny, text=white, anchor=east] at (11.9,3.25) {99.98\%};
\node[font=\scriptsize, anchor=west] at (6.8,2.4) {Douala};
\draw[fill=popgreen] (8.6,2.3) rectangle (11.85,2.6);
\node[font=\tiny, anchor=west] at (11.9,2.45) {99.71\%};
\node[font=\scriptsize, anchor=west] at (6.8,1.6) {Bafoussam};
\draw[fill=popgreen] (8.6,1.5) rectangle (11.6,1.8);
\draw[fill=poporange] (11.6,1.5) rectangle (11.75,1.8);
\node[font=\tiny, anchor=west] at (11.85,1.65) {98.90\%};
\node[font=\tiny, anchor=west] at (6.8,0.8) {Alert: Bafoussam 4G failover 3 times this week};
\end{tikzpicture}
\end{center}
\legende{Simulated Grafana dashboard: core CPU gauge, WAN bandwidth with MPLS and 4G backup curves, interface error histogram (a spike reveals a faulty port), and per-site availability bars.}
\end{popfigure}

\textbf{2. Threshold alerts.} A dashboard nobody watches is useless. The administrator defines \cle{thresholds}: ``warn if link utilisation exceeds 70\% for 10 minutes'', ``page the on-call engineer if CPU exceeds 90\% or if a site goes silent''. Alerts are delivered by e-mail, SMS or messaging tools.

\textbf{3. Event correlation.} A single fibre cut can generate 200 alarms (every device behind the cut ``disappears''). \cle{Event correlation} groups related alarms and identifies the \emph{root cause}: ``these 200 alerts all depend on the Douala edge router — investigate that link first.'' Correlation turns noise into diagnosis.

\begin{methode}[Pre-production validation checklist]
Before declaring the network operational, validate each item and record the results:
\begin{enumerate}
\item \textbf{Connectivity:} \texttt{ping} between every pair of VLANs and every site; verify expected replies and latency.
\item \textbf{Path verification:} \texttt{traceroute} from Bafoussam to the Yaoundé ERP server — the path must cross the MPLS link, not the 4G backup.
\item \textbf{Bandwidth:} run \texttt{iPerf} between sites; measured throughput must match the contracted capacity (within a reasonable margin).
\item \textbf{High availability:} physically unplug the primary MPLS link and confirm that SD-WAN fails over to 4G within the target time, with VoIP calls surviving; then restore and check failback.
\item \textbf{Security:} from a user PC, attempt SSH to a switch — it must be refused by the management ACL; from the admin workstation (VLAN 99), it must succeed.
\item \textbf{Supervision:} confirm every device appears in Grafana and that a simulated failure (shutdown of one interface) triggers the expected alert.
\end{enumerate}
\end{methode}

\begin{attention}[The missing as-built documentation]
After deployment, many teams move on without writing the \textbf{as-built documentation}: the final IP plan, VLAN table, cabling diagram, device inventory, credentials procedure and playbook repository. Months later, when the original engineer has left, every change becomes an archaeological expedition — slow, risky and expensive. The as-built file is a deliverable of the project, as important as the routers themselves. Update it after \emph{every} modification.
\end{attention}

\begin{aretenir}[Key points of this section]
\begin{itemize}
\item Design starts with \textbf{requirements}: sites, critical applications, growth, budget — validated in writing.
\item Architecture choices: \textbf{three-tier} for campuses, \textbf{spine-leaf} for data centres, \textbf{SD-WAN} (separated control/data planes, central orchestration) for inter-site connectivity.
\item A hierarchical \textbf{IP plan} (VLSM, one block per site, IPv6 GUA + ULA) enables \textbf{route summarisation} and instant readability.
\item Equipment follows the architecture: redundant core, Layer-3 distribution, PoE access, SD-WAN edge routers, NGFW.
\item \textbf{Ansible} automates base configuration, VLANs, routing and AAA — idempotently, but errors scale too: test on one device first.
\item \textbf{Supervision} = Grafana dashboards + threshold alerts + event correlation; validate everything with the pre-production checklist and finish with complete \textbf{as-built documentation}.
\end{itemize}
\end{aretenir}

\begin{exercice}[Design exercise]
CamAgro opens a new agency in Garoua with 20 users, 10 VoIP phones and 4 servers.
\begin{enumerate}
\item Propose a site block inside 10.0.0.0/8, coherent with the existing plan, and subnet it with VLSM for the three services.
\item Which single summary route will the Garoua edge router advertise to the other sites?
\item State two checks from the pre-production checklist that are specific to this new site, and justify them.
\end{enumerate}
\end{exercice}

\begin{corrige}[Exercise 1]
\begin{enumerate}
\item Following the site-numbering logic (10 = Yaoundé, 20 = Douala, 30 = Bafoussam), assign \textbf{10.40.0.0/16} to Garoua. VLSM subnets: users VLAN 10: \texttt{10.40.10.0/27} (30 usable $\geq$ 20); VoIP VLAN 20: \texttt{10.40.20.0/28} (14 usable $\geq$ 10); servers VLAN 30: \texttt{10.40.30.0/29} (6 usable $\geq$ 4). (Other coherent choices are acceptable if sizes are justified.)
\item The edge router advertises the single summary route \textbf{10.40.0.0/16}, hiding all internal subnets.
\item Two site-specific checks: (a) \texttt{traceroute} from Garoua to the Yaoundé ERP server to confirm traffic uses the MPLS path and not the 4G backup; (b) an \texttt{iPerf} test between Garoua and Yaoundé to verify the contracted bandwidth, since Garoua's long-distance link is the most likely bottleneck. The HA failover test (unplug MPLS, check 4G takeover) is also essential for this remote site.
\end{enumerate}
\end{corrige}

\coursec{Activités d'intégration : étude de cas, simulation et rapport}

Having mastered the design, deployment, and monitoring of an enterprise network in the previous section, we now bring everything together. Real-world networking is not a linear sequence of isolated tasks; it is an iterative cycle of \textbf{analysis}, \textbf{design}, \textbf{implementation}, \textbf{validation}, and \textbf{documentation}. This section simulates a complete project lifecycle, mirroring the expectations of the GCE Advanced Level practical examination and the demands of a professional ICT engineer in Cameroon.

\begin{definition}[Network Integration]
\textbf{Network Integration} is the disciplined process of assembling heterogeneous hardware (routers, switches, firewalls, servers), software (OS, monitoring agents, applications), policies (QoS, security, routing), and operational processes (backup, change management, disaster recovery) into a single, coherent, manageable, and resilient infrastructure that meets defined business requirements.
\end{definition}

\begin{propriete}[The Quality Assurance Loop: Design $\rightarrow$ Test $\rightarrow$ Validate $\rightarrow$ Document]
Every successful network project follows a closed feedback loop. A design is \textbf{not} complete until it has been tested under realistic conditions (load, failure, attack), the test results have been \textbf{validated} against the original requirements (SLAs), and the entire lifecycle has been \textbf{documented} for handover and future audits. Skipping any arrow in this loop introduces technical debt that surfaces as outages or security breaches.
\end{propriete}

\courssub{Case Study: Interconnecting Five University Campuses}

\textbf{Scenario:} The Ministry of Higher Education (MINESUP) wishes to interconnect the main campuses of five state universities: \textbf{Yaoundé I}, \textbf{Yaoundé II}, \textbf{Douala}, \textbf{Buea}, and \textbf{Dschang}. The network must support:
\begin{itemize}
    \item \textbf{VoIP telephony} (Cisco CUCM / Asterisk) between campus switchboards.
    \item \textbf{High-definition videoconferencing} for Senate meetings and remote lectures.
    \item A centralized \textbf{Learning Management System (LMS)} (Moodle) hosted in a Data Centre at Yaoundé I, accessed by all campuses.
    \item \textbf{Internet access} for students and staff, with content filtering.
    \item \textbf{Management traffic} (SNMP, Syslog, NetFlow, SSH) on a dedicated out-of-band network.
\end{itemize}

\textbf{Constraints:}
\begin{itemize}
    \item Budget favors \textbf{IPsec VPN over public Internet} (CAMTEL/MTN/Orange links) for inter-campus WAN, with a primary MPLS link only between Yaoundé I and Douala.
    \item Each campus has a /22 public IPv4 block and an IPv6 /48 allocation from ANTIC.
    \item High availability: no single point of failure for the LMS and VoIP core.
    \item QoS: Voice (EF), Video (AF41), LMS/Database (AF21), Best Effort (BE), Scavenger (CS1).
\end{itemize}

\begin{experience}[Group Design Workshop (3 Hours)]
\textbf{Objective:} Produce a \textbf{Logical Topology Diagram}, a \textbf{Comprehensive Addressing Plan} (VLSM + IPv6), and a \textbf{Technology Decision Matrix}.

\textbf{Deliverables per group (4-5 students):}
\begin{enumerate}
    \item \textbf{Logical Diagram:} Core, Distribution, Access layers per campus; WAN edge; DMZ for LMS; Security zones (Inside, DMZ, Management, Guest).
    \item \textbf{Addressing Plan (IPv4):} Use VLSM on the private 10.0.0.0/8 block.
        \begin{itemize}
            \item Yaoundé I (HQ + DC): 10.1.0.0/16 $\rightarrow$ Subnet for DC (10.1.0.0/19), Campus LANs (10.1.32.0/19), Management (10.1.254.0/24), VPN Pool (10.1.255.0/24).
            \item Yaoundé II: 10.2.0.0/16, Douala: 10.3.0.0/16, Buea: 10.4.0.0/16, Dschang: 10.5.0.0/16. (Same subnetting pattern).
            \item Inter-router WAN links (MPLS/IPsec): 10.255.255.0/30, 10.255.255.4/30, ... (Point-to-point /30s; modern practice prefers /31s per RFC 3021 to conserve address space).
        \end{itemize}
    \item \textbf{IPv6 Plan:} 2001:db8:1000::/48 (Yaoundé I), 2001:db8:2000::/48 (Yaoundé II), etc. Subnet /64 per VLAN (using the documentation prefix 2001:db8::/32 per RFC 3849).
    \item \textbf{Technology Choices:} Routing Protocol (OSPF Area 0 backbone + Area per campus / BGP for MPLS PE-CE), VPN (DMVPN Phase 3 with IPsec), QoS (MQC with LLQ), Security (Zone-Based Firewall / ASA), HA (HSRP/VRRP for gateways, StackWise/VSS for core switches).
\end{enumerate}
\end{experience}

\courssub{Simulation Lab: Packet Tracer / GNS3 / EVE-NG}

The design is implemented in a network emulator. We use \textbf{GNS3} (or EVE-NG) for realism (real IOSv/IOS-XEv images) but \textbf{Cisco Packet Tracer} is acceptable for the GCE exam context.

\begin{methode}[Step-by-Step Implementation in the Simulator]
\textbf{Phase 1: Campus LAN Core (Yaoundé I Example)}
\begin{enumerate}
    \item Deploy two Core Switches (e.g., Catalyst 9500 / IOSvL2) in \textbf{VSS} (Virtual Switching System) or \textbf{StackWise}.
    \item Configure \textbf{VLANs}: 10 (Mgmt), 20 (Servers/DC), 30 (VoIP), 40 (Video), 50 (Staff), 60 (Students), 70 (Guest), 99 (Native/Unused).
    \item Configure \textbf{SVIs} (Layer 3 interfaces) on the Core VSS for each VLAN (Gateway IPs).
    \item Enable \textbf{HSRP} (or VRRP) on SVIs for gateway redundancy (if not using VSS).
    \item Configure \textbf{Access Switches} (IOSvL2): Trunk uplinks, \texttt{spanning-tree portfast} on access ports, \texttt{voice vlan} for IP Phones.
    \item \textbf{Port Security} / \textbf{802.1X} (simulated) on access ports.
\end{enumerate}

\textbf{Phase 2: WAN Edge \& Routing}
\begin{enumerate}
    \item Deploy Edge Routers (ISR 4000 / IOSv) at each campus.
    \item \textbf{MPLS Link (Yaoundé I $\leftrightarrow$ Douala):} Configure OSPF Area 0 on the MPLS interfaces. Redistribute connected/internal routes.
    \item \textbf{Internet Links (All campuses):} Configure default routes towards ISP (static or BGP). Configure \textbf{NAT} (PAT) for Internet access.
    \item \textbf{DMVPN Phase 3 Hub (Yaoundé I) \& Spokes:}
        \begin{itemize}
            \item Tunnel interface: \texttt{ip nhrp network-id 1}, \texttt{ip nhrp nhs <hub-tunnel-ip>}, \texttt{ip nhrp shortcut}, \texttt{ip nhrp redirect}.
            \item IPsec Profile: \texttt{crypto ipsec profile DMVPN-PROF}, \texttt{set transform-set AES256-SHA}.
            \item Run \textbf{OSPF} (or EIGRP) over the tunnel interfaces (Network type broadcast/point-to-multipoint).
        \end{itemize}
\end{enumerate}

\textbf{Phase 3: Services \& Security}
\begin{enumerate}
    \item \textbf{QoS (MQC):} Class-maps matching DSCP (EF, AF41, AF21, CS1). Policy-map with \texttt{priority} for Voice, \texttt{bandwidth remaining percent} for others. Apply \texttt{service-policy output} on WAN interfaces.
    \item \textbf{Firewall (Zone-Based Policy Firewall - ZPFW):} Zones: \texttt{INSIDE}, \texttt{DMZ}, \texttt{MGMT}, \texttt{WAN}, \texttt{VPN}. Inter-zone policies: Inspect (Inside$\rightarrow$WAN), Pass (Mgmt$\rightarrow$All), Drop/Log (WAN$\rightarrow$Inside except VPN).
    \item \textbf{VPN Remote Access (AnyConnect / SSL VPN):} For staff working from home (optional but impressive).
    \item \textbf{Services:} DHCP (on Core SVIs or dedicated server), DNS (internal views), NTP (master at Yaoundé I, peers at others), Syslog/SNMP/NetFlow exporters to Mgmt Server.
\end{enumerate}
\end{methode}

\begin{popfigure}
\begin{tikzpicture}[
    node distance=1.5cm and 2cm,
    campus/.style={draw=popblue, thick, rounded corners, fill=popblue!5, minimum height=2.5cm, minimum width=3.5cm, align=center},
    core/.style={draw=popdark, thick, circle, fill=poporange!20, minimum size=1.2cm},
    dist/.style={draw=popgreen, thick, rectangle, fill=popgreen!10, minimum size=1cm},
    access/.style={draw=poppurple, thick, rectangle, fill=poppurple!10, minimum size=0.8cm},
    wan/.style={draw=poppink, thick, diamond, fill=poppink!10, aspect=2},
    link/.style={very thick, popdark},
    testlabel/.style={font=\tiny, text=popdark, fill=white, fill opacity=0.8, inner sep=1pt, rounded corners, align=left, text width=5cm}
]
% Yaoundé I Campus (HQ + DC)
\node[campus, align=center] (ya1) at (0,0){Yaoundé I (HQ/DC)\\\textbf{10.1.0.0/16}\\\textit{OSPF Area 0}};
\node[core, align=center] (ya1_core) at (-1.5,0.5){Core\\VSS};
\node[dist, align=center] (ya1_dmz) at (1.5,0.5){DMZ\\LMS/FW};
\node[dist, align=center] (ya1_mgmt) at (1.5,-0.5){Mgmt\\Server};
\node[wan, align=center] (ya1_wan) at (-3,0){MPLS/\\Internet};

% Douala Campus
\node[campus, align=center] (dla) at (8,2){Douala\\\textbf{10.3.0.0/16}\\\textit{OSPF Area 3}};
\node[core] (dla_core) at (6.5,2.5) {Core};
\node[wan] (dla_wan) at (9.5,2) {Internet};

% Yaoundé II Campus
\node[campus, align=center] (ya2) at (8,-2){Yaoundé II\\\textbf{10.2.0.0/16}\\\textit{OSPF Area 2}};
\node[core] (ya2_core) at (6.5,-1.5) {Core};
\node[wan] (ya2_wan) at (9.5,-2) {Internet};

% Buea Campus
\node[campus, align=center] (buea) at (14,3){Buea\\\textbf{10.4.0.0/16}\\\textit{OSPF Area 4}};
\node[core] (buea_core) at (12.5,3.5) {Core};
\node[wan] (buea_wan) at (15.5,3) {Internet};

% Dschang Campus
\node[campus, align=center] (dsch) at (14,-3){Dschang\\\textbf{10.5.0.0/16}\\\textit{OSPF Area 5}};
\node[core] (dsch_core) at (12.5,-2.5) {Core};
\node[wan] (dsch_wan) at (15.5,-3) {Internet};

% Connections
\draw[link] (ya1_core) -- node[testlabel, above, align=center]{MPLS: 1Gbps\\OSPF Adj: FULL} (dla_core);
\draw[link, densely dashed] (ya1_wan) -- node[testlabel, above, sloped, align=center]{DMVPN Tunnel\\IPsec: UP} (dla_wan);
\draw[link, densely dashed] (ya1_wan) -- (ya2_wan);
\draw[link, densely dashed] (ya1_wan) -- (buea_wan);
\draw[link, densely dashed] (ya1_wan) -- (dsch_wan);
\draw[link] (dla_core) -- (ya2_core);
\draw[link] (dla_core) -- (buea_core);
\draw[link] (ya2_core) -- (dsch_core);

% Internal Yaoundé I
\draw[link] (ya1_core) -- (ya1_dmz);
\draw[link] (ya1_core) -- (ya1_mgmt);

% Test Annotations
\node[testlabel, align=center] at (-1.5, -1.5){
    \textbf{Test Results (Yaoundé I Core):}\\
    iPerf TCP (VoIP): 64 kbps, Jitter 1.2ms, Loss 0\%\\
    iPerf TCP (Video): 4 Mbps, Jitter 5ms, Loss 0.01\%\\
    Failover (MPLS$\rightarrow$DMVPN): 1.8s convergence\\
    Nmap Scan (WAN): Only 500/4500 UDP (IKE) open\\
    ACL Test: Student VLAN $\nrightarrow$ Mgmt VLAN: \textcolor{popgreen}{DROP}
};
\node[testlabel, align=center] at (15.5, 0){
    \textbf{DMVPN Hub Status:}\\
    4 Spokes Registered\\
    NHRP Cache: Complete\\
    OSPF Neighbors: 4/4 Full\\
    IPv6 Tunnels: UP
};
\end{tikzpicture}
\legende{Final Deployed Topology in Simulator with Key Test Result Annotations}
\end{popfigure}

\courssub{Validation: Load, Failover, and Security Testing}

A network is not "done" when the configuration is saved. It is done when it proves it meets SLAs under stress.

\begin{experience}[Test Plan Execution (2 Hours)]
\textbf{Tools:} \texttt{iPerf3} (installed on Linux VMs in DC and each campus), \texttt{Nmap}, \texttt{Wireshark} (SPAN port), Simulator link failure injection.

\begin{enumerate}
    \item \textbf{Baseline Performance (iPerf3):}
\begin{itemize}
    \item \texttt{iperf3 -c \textless DC\_IP\textgreater -p 5001 -t 60 -P 4} (4 parallel streams, 60s) for LMS/Database traffic (AF21). Target: $> 900$ Mbps on 1G links.
    \item \texttt{iperf3 -c \textless DC\_IP\textgreater -u -b 64k -t 60} (UDP 64kbps) for VoIP (EF). Target: Jitter $< 30$ms, Loss $< 1$\%.
    \item \texttt{iperf3 -c \textless DC\_IP\textgreater -u -b 4M -t 60} for Video (AF41). Target: Jitter $< 50$ms.
\end{itemize}
    \item \textbf{Failover / Resilience Testing:}
        \begin{itemize}
            \item \textbf{Scenario A:} Shut MPLS interface on Yaoundé I Core. Measure OSPF/DMVPN convergence time (ping continuous from Douala to DC). Target: $< 2$s (sub-second with BFD + tuned timers).
            \item \textbf{Scenario B:} Shut Core Switch 1 in Yaoundé I (VSS split). Verify VSS dual-active detection and recovery (or HSRP failover if not VSS). Target: $< 1$s.
            \item \textbf{Scenario C:} Simulate ISP outage on Douala Internet link. Verify DMVPN re-establishes via backup ISP (if dual-homed) or traffic shifts to MPLS.
        \end{itemize}
    \item \textbf{Security Validation:}
\begin{itemize}
            \item \textbf{Nmap Scan from Internet:} \texttt{nmap -sS -p- \textless Yaoundé\_I\_Public\_IP\textgreater}. Verify \textbf{only} UDP 500/4500 (IKE/IPsec NAT-T) and TCP 443 (SSL VPN) are open. All others filtered/closed.
            \item \textbf{Internal Segmentation Test:} From Student VLAN PC, attempt SSH to Mgmt Server (10.1.254.10). Verify \texttt{Connection refused / Timeout} (ACL Drop). Verify \texttt{syslog} shows deny.
            \item \textbf{VPN Penetration:} Attempt to connect DMVPN spoke with wrong pre-shared key. Verify ISAKMP failure logs.
        \end{itemize}
\end{enumerate}
\end{experience}

\begin{exoresolu}[Worked Example: Analyzing iPerf3 Output for QoS Validation]
\textbf{Statement:} During the VoIP test (UDP 64kbps, EF marking) from Douala to Yaoundé I DC over the DMVPN tunnel, the following output is observed on the receiver (DC):
\texttt{[  5]   0.00-60.00 sec  480 KBytes  65.5 Kbits/sec  2.15 ms  0/342 (0\%)}\\
\texttt{[  5] Sent 342 datagrams}

During a concurrent "Student Download" test (TCP, Best Effort, 500 Mbps) saturating the 100 Mbps WAN link, the VoIP test is repeated:
\texttt{[  5]   0.00-60.00 sec  475 KBytes  64.8 Kbits/sec  48.3 ms  12/340 (3.5\%)}\\
\texttt{[  5] Sent 340 datagrams}

\textbf{Question:} Is the QoS policy working correctly? Justify your answer with reference to the markings and queue behavior.

\textbf{Solution:}
\begin{enumerate}
    \item \textbf{Baseline (No Load):} Jitter = 2.15 ms, Loss = 0\%. \textbf{Excellent.} The LLQ (Priority Queue) for EF traffic is serviced immediately.
    \item \textbf{Under Load (Congestion):} Jitter jumps to 48.3 ms, Loss = 3.5\%.
    \item \textbf{Analysis:} For G.711/G.729 VoIP, the ITU-T G.114 recommendation states \textbf{one-way delay $< 150$ms} (jitter contributes to delay variation). 48ms jitter is \textbf{acceptable but borderline} for high quality. However, \textbf{3.5\% packet loss is unacceptable} (target $< 1\text{\%}$). This indicates the \textbf{Priority Queue (PQ) is being policed} (likely \texttt{priority percent 10} or similar) and the VoIP traffic is exceeding the reserved bandwidth, causing tail-drop in the PQ, OR the shaper/policer on the WAN interface is dropping EF packets because the \textbf{parent shaper} rate is exceeded and the child priority queue is not guaranteed its rate under the shaper.
    \item \textbf{Remediation:} Increase the \texttt{priority} bandwidth allocation for EF class (e.g., to 20\% or fixed kbps matching max concurrent calls $\times$ 80kbps). Ensure \texttt{police} is not dropping conforming packets. Verify \texttt{shape average} on parent policy allows child priority queue to burst.
\end{enumerate}
\end{exoresolu}

\courssub{Technical Report Writing}

The report is the permanent record of the project. In the GCE exam, the report \textbf{is} a major part of the mark. It must be structured, evidence-based, and reflective.

\begin{popfigure}
\begin{tikzpicture}[
    node distance=0.2cm,
    section/.style={draw=popblue, thick, rounded corners, fill=popblue!10, minimum height=1.2cm, text width=5cm, align=center, font=\bfseries},
    detail/.style={draw=popgreen, rounded corners, fill=popgreen!5, minimum height=0.8cm, text width=5cm, align=left, font=\small},
    arrow/.style={->, thick, popdark}
]
% Title
\node[section, fill=popdark, text=white, minimum width=12cm] (title) {TECHNICAL PROJECT REPORT: MINESUP Inter-Campus Network};
% Sections
\node[section] (s1) [below=of title] {1. Executive Summary};
\node[detail] (d1) [right=of s1] {Business need, scope, key results (SLA met?), recommendation.};

\node[section] (s2) [below=of s1] {2. Requirements Analysis};
\node[detail] (d2) [right=of s2] {Functional (VoIP, Video, LMS), Non-functional (Availability 99.9\%, Security, Scalability), Constraints (Budget, IPv4/IPv6).};

\node[section] (s3) [below=of s2] {3. Logical \& Physical Design};
\node[detail] (d3) [right=of s3] {Topology diagram (logical + physical rack), Addressing Plan (VLSM Table), VLAN Plan, Routing Design (OSPF Areas, DMVPN), QoS Model (8-class), Security Zones.};

\node[section] (s4) [below=of s3] {4. Implementation Details};
\node[detail] (d4) [right=of s4] {Key Config Snippets (Core VSS, DMVPN Hub/Spoke, ZPFW Policy, QoS Policy-map, HSRP). \emph{Not full configs!} Version numbers (IOS 17.3.x).};

\node[section] (s5) [below=of s4] {5. Test Plan \& Results};
\node[detail] (d5) [right=of s5] {Test Case ID, Description, Tool/Command, Expected Result, Actual Result (Screenshot/Log), Pass/Fail. \textbf{Include iPerf tables, Failover convergence times, Nmap output.}};

\node[section] (s6) [below=of s5] {6. Lessons Learned \& Future Improvements};
\node[detail] (d6) [right=of s6] {What failed? (e.g., MTU issue on DMVPN $\rightarrow$ \texttt{ip tcp adjust-mss 1360}). What would you change? (BGP for MPLS, SD-WAN migration, IPv6 native). Disaster Recovery Plan status.};

\node[section] (s7) [below=of s6] {7. Appendices};
\node[detail] (d7) [right=of s7] {Full Configs (zipped), Cable Schedule, License Inventory, Gantt Chart, Risk Register.};

% Arrows
\foreach \s/\d in {s1/d1, s2/d2, s3/d3, s4/d4, s5/d5, s6/d6, s7/d7}
    \draw[arrow] (\s.east) -- (\d.west);
\end{tikzpicture}
\legende{Standard Technical Report Structure (Mandatory for GCE Project Submission)}
\end{popfigure}

\begin{methode}[Writing the "Test Results" Section (Section 5) – The Evidence]
Never write "The network works." Write:
\begin{enumerate}
    \item \textbf{Test Case ID:} TC-WAN-03.
    \item \textbf{Objective:} Verify VoIP QoS (EF) under WAN congestion.
    \item \textbf{Setup:} iPerf3 Server (DC, 10.1.10.5), Client (Douala, 10.3.50.10). Background load: 5 TCP streams (BE) saturating 100Mbps WAN.
    \item \textbf{Command:} \texttt{iperf3 -c 10.1.10.5 -u -b 64k -t 60 -S 0xB8} (DSCP 46 = EF).
    \item \textbf{Expected:} Jitter $< 30$ms, Loss $< 1$\%.
    \item \textbf{Actual:} Jitter 12.4ms, Loss 0.2\%. \textbf{PASS}.
    \item \textbf{Evidence:} \textit{(Insert screenshot of terminal output + router \texttt{show policy-map interface} showing priority queue counters incrementing)}.
\end{enumerate}
\end{methode}

\courssub{Oral Defence: Simulated GCE Jury}

The final step is a 10-15 minute presentation followed by 10 minutes of questions. Treat it like a project defence.

\begin{aretenir}[GCE Oral Defence: Marking Grid Criteria (Typical)]
The examiner (jury) evaluates five competencies. Prepare one slide / talking point per competency.
\begin{center}
\begin{tabularx}{\linewidth}{|l|X|}\hline
\textbf{Competency} & \textbf{What to Demonstrate} \\ \hline
\textbf{1. Analysis} & You understood the Ministry's needs (VoIP, Video, LMS) and translated them into technical specs (Bandwidth, Latency, Security Zones). \\ \hline
\textbf{2. Design} & Your diagram is logical (hierarchical), addressing is scalable (VLSM), protocols are justified (OSPF vs BGP, DMVPN vs MPLS), QoS classes match apps. \\ \hline
\textbf{3. Implementation} & You can navigate the simulator CLI: \texttt{show run | section}, \texttt{debug}, \texttt{ping/traceroute} sourcing from correct VRF/interface. You fix a deliberate error injected by examiner. \\ \hline
\textbf{4. Documentation} & Your report is professional: TOC, numbered figures/tables, config snippets explained, test tables with PASS/FAIL, lessons learned are honest. \\ \hline
\textbf{5. Communication} & Clear speech, eye contact, correct terminology (not "cable" but "link/interface", not "box" but "router/switch"), admits limitations ("We didn't test IPv6 multicast"). \\ \hline
\end{tabularx}
\end{center}
\end{aretenir}

\begin{attention}[Common "Jury Trap" Questions – Prepare Answers]
\begin{itemize}
    \item \textbf{"Why OSPF and not BGP everywhere?"} $\rightarrow$ OSPF for fast convergence inside AS (campuses), BGP only at MPLS PE-CE or Internet edge. DMVPN uses OSPF/EIGRP over tunnel for simplicity.
    \item \textbf{"Your DMVPN spokes use the same tunnel subnet (10.255.255.0/24). How does Hub distinguish them?"} $\rightarrow$ NHRP mapping (Public IP $\leftrightarrow$ Tunnel IP). Hub uses \texttt{ip nhrp map multicast dynamic}.
    \item \textbf{"What happens to VoIP if the Priority Queue fills up?"} $\rightarrow$ Packets are dropped (tail drop) $\rightarrow$ Voice quality degrades. Solution: Admission Control (Call Admission Control - CAC) on CUCM or \texttt{priority} with \texttt{police} conform-action transmit exceed-action drop (but better to size queue correctly).
    \item \textbf{"How do you recover the LMS database if the Yaoundé I DC burns down?"} $\rightarrow$ \textbf{This is the DR question.} Answer: Async replication to Douala DC (standby VM), RPO $< 1$hr, RTO $< 4$hr. \textbf{If you haven't tested the failover script, you fail this question.}
    \item \textbf{"Show me the command to verify IPsec SA on the Hub."} $\rightarrow$ \texttt{show crypto ipsec sa} / \texttt{show crypto session detail}.
\end{itemize}
\end{attention}

\begin{savaistu}[Cameroon Context: ANTIC \& Cybersecurity Law]
Law No. 2010/012 on Cybersecurity and Cybercrime mandates that critical infrastructure (like a Ministry network) must implement \textbf{security audits}, \textbf{incident response plans}, and \textbf{data protection} measures. In your report's "Security" and "Lessons Learned" sections, explicitly reference:
\begin{itemize}
    \item Compliance with \textbf{ANTIC} guidelines for public administration networks.
    \item Logging retention (1 year minimum) on the Syslog server.
    \item Encryption of data in transit (IPsec, TLS 1.2+ for LMS) and at rest (LUKS/BitLocker on DB servers).
    \item Regular vulnerability scanning (OpenVAS/Nessus) scheduled via cron.
\end{itemize}
This shows professional awareness beyond the classroom.
\end{savaistu}

\courssub{Consolidation Exercises}

\begin{exercice}[Integration Challenge: The "Friday Afternoon" Scenario]
It is Friday 4:30 PM. The Douala Campus network admin calls: "Students cannot access the LMS (Moodle) since 2 PM. VoIP calls to Yaoundé sound robotic. Internet is fine."
You have remote SSH access to Douala Edge Router and Core Switch.
\begin{enumerate}
    \item List the \textbf{first 5 verification commands} you would run (in order) to isolate the fault domain (L1/L2/L3/Service/Security).
    \item The \texttt{show ip ospf neighbor} on Douala Core shows the neighbor to Yaoundé I Hub is \textbf{DOWN/INIT} (stuck). The DMVPN tunnel interface is \textbf{Up/Up}. What are the \textbf{top 3 causes}? (Hint: MTU, ACL, OSPF Network Type).
    \item You find an ACL on the Douala WAN interface blocking \texttt{protocol 89} (OSPF). You fix it. Neighbors come \textbf{FULL}. But VoIP is still robotic. \texttt{show policy-map interface Tunnel0} shows \textbf{0 packets} matched in the \texttt{class VOICE}. Why? (Hint: DSCP marking / Trust boundary / Tunnel header copy).
    \item Write the \textbf{exact configuration commands} to fix the QoS marking trust on the Douala Access Switch uplinks and ensure DSCP is copied to the IPsec outer header (or preserved inside).
\end{enumerate}
\end{exercice}

\begin{corrige}[Integration Challenge]
\begin{enumerate}
    \item \textbf{Verification Sequence:}
        1. \texttt{ping 10.1.10.5 source vlan 50} (L3 reachability to LMS from Staff VLAN).
        2. \texttt{show ip route 10.1.10.5} (Routing table check).
        3. \texttt{show crypto session detail} / \texttt{show ip nhrp} (DMVPN tunnel health).
        4. \texttt{show policy-map interface Tunnel0} (QoS policy applied? counters moving?).
        5. \texttt{show logging | include \%SECURITY|\%ACL|\%OSPF} (Recent security/routing events).
    \item \textbf{OSPF Stuck in INIT over DMVPN (Tunnel Up):}
        1. \textbf{MTU Mismatch:} IPsec overhead (approx 60-80 bytes) reduces effective MTU. OSPF Database Description (DBD) packets are large and have DF bit set. \textbf{Fix:} \texttt{ip mtu 1400} on Tunnel interface + \texttt{ip tcp adjust-mss 1360}.
        2. \textbf{ACL Blocking OSPF (Protocol 89) or Multicast 224.0.0.5/6} on physical WAN interface or Tunnel interface.
        3. \textbf{OSPF Network Type Mismatch:} Hub \texttt{broadcast}, Spoke \texttt{point-to-multipoint} (or vice versa) $\rightarrow$ Hello parameters mismatch.
    \item \textbf{QoS Counters Zero on Tunnel:}
        The traffic entering the tunnel is \textbf{not marked} (or marking is stripped). On the Access Switch, the uplink to Core/Router must \textbf{trust DSCP} (\texttt{mls qos trust dscp} or \texttt{trust device cisco-phone} then \texttt{trust dscp}). If the phone marks EF, but the switch port is \texttt{untrusted} (default), the switch resets to BE (0). Also, verify \texttt{qos pre-classify} on the Tunnel interface so the policy-map sees the \textbf{inner packet's DSCP} before encryption.
    \item \textbf{Fix Commands (Douala Access Switch Uplink \& Router Tunnel):}
\texttt{! On Access Switch Uplink (GigabitEthernet1/0/1)}\\
\texttt{interface GigabitEthernet1/0/1}\\
\texttt{ switchport trunk encapsulation dot1q}\\
\texttt{ switchport mode trunk}\\
\texttt{ mls qos trust dscp          ! Trust the marking from IP Phone / PC}\\
\texttt{ spanning-tree portfast trunk}\\
\\
\texttt{! On Edge Router Tunnel Interface}\\
\texttt{interface Tunnel0}\\
\texttt{ ip mtu 1400}\\
\texttt{ ip tcp adjust-mss 1360}\\
\texttt{ qos pre-classify            ! Classify based on inner header (pre-encryption)}\\
\texttt{ service-policy output WAN\_QOS\_OUT}
\end{enumerate}
\end{corrige}

\begin{exercice}[GCE Style: Report Critique]
You are given a student's "Lessons Learned" section (excerpt below). Identify \textbf{three major weaknesses} and rewrite the section professionally.

\textit{Student Excerpt:} "We learned that networking is hard. The simulator crashed a lot. We couldn't get BGP to work so we used OSPF. The firewall blocked some traffic but we disabled it for testing. Next time we will buy better computers. We didn't test the backup link because we ran out of time. The report was written the night before."
\end{exercice}

\begin{corrige}[Report Critique]
\textbf{Weaknesses:}
1. \textbf{Unprofessional tone / Blaming tools:} "Simulator crashed", "better computers". Engineers solve tool constraints.
2. \textbf{Lack of technical depth:} "Couldn't get BGP to work" $\rightarrow$ No root cause analysis (ASN mismatch? Missing network command?).
3. \textbf{Security violation / Bad practice:} "Disabled firewall for testing". Never disable security controls; modify rules with \texttt{log} to see drops.
4. \textbf{Incomplete Validation:} "Didn't test backup link" $\rightarrow$ Failed the Design $\rightarrow$ Test loop. Major mark deduction.
5. \textbf{Poor time management:} "Written night before".

\textbf{Professional Rewrite:}
\begin{quote}
\textbf{6. Lessons Learned \& Future Improvements}
\begin{itemize}
    \item \textbf{Routing Protocol Selection:} Initial design proposed eBGP for WAN. During implementation, BGP peering failed due to \texttt{ebgp-multihop} requirement over DMVPN and missing \texttt{network} statements for connected subnets. Due to time constraints, we migrated to OSPF (network type \texttt{point-to-multipoint}) which established adjacencies immediately. \textbf{Action:} In production, use BGP for MPLS PE-CE and Internet edge; retain OSPF for DMVPN overlay.
    \item \textbf{Firewall Policy Validation:} Early testing showed LMS access failures. Investigation revealed the Zone-Based Firewall policy between \texttt{VPN} and \texttt{DMZ} zones lacked an \texttt{inspect} action for the custom TCP port 8443 (Moodle HTTPS). The rule was corrected; \textbf{the firewall was never disabled}. All deny actions log to Syslog.
    \item \textbf{Resilience Testing Gap:} The MPLS-to-DMVPN failover test (TC-HA-02) was \textbf{not executed} due to simulator license expiry on the final day. This is a \textbf{critical gap}. \textbf{Mitigation:} Schedule failover tests 48h before handover. Documented DR runbook (Appendix D) remains unvalidated.
    \item \textbf{Tooling Stability:} GNS3 VM crashed twice during high-throughput iPerf tests (10Gbps simulated). Migrated heavy load tests to physical lab switches (Catalyst 2960) for final validation.
\end{itemize}
\end{quote}
\end{corrige}

\begin{aretenir}[Key Takeaways for Section 4]
\begin{itemize}
    \item \textbf{Integration} is the art of making disparate technologies (Routing, Switching, Security, QoS, VPN, IPv6, Services) work \textbf{together} to serve business apps (VoIP, Video, LMS).
    \item The \textbf{Design $\rightarrow$ Test $\rightarrow$ Validate $\rightarrow$ Document} loop is non-negotiable. Untested redundancy is a myth.
    \item \textbf{Simulation} (Packet Tracer/GNS3) is where you prove the design. Capture \textbf{evidence} (screenshots, logs, \texttt{show} outputs) for the report.
    \item The \textbf{Technical Report} is your deliverable. Structure it rigorously: Requirements $\rightarrow$ Design $\rightarrow$ Config Highlights $\rightarrow$ Test Evidence $\rightarrow$ Lessons Learned.
    \item The \textbf{Oral Defence} tests your ownership. Know \textbf{why} you chose every protocol, every IP address, every QoS class. Be ready to troubleshoot live.
    \item \textbf{Cameroon Reality:} Reference ANTIC guidelines, CAMTEL/MTN/Orange as ISPs, Mobile Money API integration security, and the frequent need for \textbf{cost-effective} solutions (DMVPN over MPLS).
\end{itemize}
\end{aretenir}

\newpage

\coursec{Integration activity}

\coursec{Integration Activities: Designing, Deploying and Supervising an Enterprise Network}

\courssub{Aims of the activity}
\begin{aretenir}[Learning Outcomes]
By the end of this integration activity, you will be able to:
\begin{itemize}
    \item Analyse the \cle{complexity factors} (scale, heterogeneity, dynamics, administrative domains) in a multi-campus enterprise network.
    \item Identify and mitigate \cle{operational challenges}: performance bottlenecks, high-availability requirements, security threats, and management overhead.
    \item Design a complete network architecture: logical/physical topology, IP addressing plan, VLAN segmentation, routing protocols (OSPF/BGP), redundancy mechanisms (VRRP, link aggregation).
    \item Implement the design in a network emulator (GNS3) with realistic device configurations (Cisco IOS / FRR).
    \item Automate initial deployment and ongoing configuration changes using \cle{Ansible} playbooks and Jinja2 templates.
    \item Deploy a modern \cle{observability stack} (Prometheus node-exporter, SNMP exporter, Alertmanager, Grafana) for real-time monitoring and alerting.
    \item Conduct structured \cle{resilience testing} (link failure, device reboot, congestion) and document results.
    \item Produce a professional \cle{technical report} (15 pages) and deliver a concise \cle{oral presentation} (10 minutes) aligned with GCE Advanced Level marking criteria.
\end{itemize}
\end{aretenir}

\courssub{Situation}
\begin{experience}[CampusNet — Inter-Campus Network for Cameroon Public Universities]
The Ministry of Higher Education (MINESUP) has launched the \textbf{CampusNet} project to interconnect three major university campuses:
\begin{itemize}
    \item \textbf{Campus A — Yaoundé (Main Campus)}: 12\,000 students, 800 staff, Data Centre hosting LMS (Moodle), library system, ERP, research storage (2\,PB), VoIP PBX, video-conferencing MCU.
    \item \textbf{Campus B — Douala (Economic Campus)}: 8\,500 students, 500 staff, local LMS mirror, IoT sensor network for smart-building (500 devices), student Wi-Fi (3\,000 concurrent).
    \item \textbf{Campus C — Buea (Technology Campus)}: 6\,000 students, 400 staff, High-Performance Computing (HPC) cluster (500 cores), research data replication, eduroam RADIUS.
\end{itemize}

\medskip
\textbf{Current infrastructure:}
\begin{itemize}
    \item Each campus has a collapsed core (two Catalyst 9500 stacks) and access switches (Catalyst 2960X).
    \item WAN links: Yaoundé–Douala 2\,Gbps dark fibre (CAMTEL), Yaoundé–Buea 1\,Gbps MPLS (Orange Cameroon), Douala–Buea 500\,Mbps backup VPN over Internet (MTN Business).
    \item Internet exit: Yaoundé 2\,Gbps (CAMTEL), Douala 1\,Gbps (Orange), Buea 500\,Mbps (MTN). Each campus runs its own BGP AS (65001, 65002, 65003) peering with ISPs.
    \item Management: SSH/CLI only, no automation, no centralised logging, SNMP v2c read-only.
    \item Security: Perimeter ASA firewalls, no internal segmentation, flat VLAN 1 everywhere.
\end{itemize}

\medskip
\textbf{Requirements from MINESUP (Cahier des charges):}
\begin{enumerate}
    \item \textbf{Performance}: Latency $\le$ 15\,ms between any two campuses; jitter $\le$ 5\,ms for VoIP/video; throughput $\ge$ 800\,Mbps sustained for HPC replication.
    \item \textbf{Availability}: 99.99\% uptime for core services (LMS, ERP, VoIP); automatic failover $<$ 500\,ms for WAN link failure.
    \item \textbf{Security}: Zero-Trust segmentation (micro-perimeters per service), 802.1X for wired/wireless, centralised logging (Syslog/NetFlow), intrusion detection (Suricata), encrypted management (SSH keys, SNMPv3).
    \item \textbf{Management}: Infrastructure-as-Code (GitLab repo), automated provisioning (Ansible), configuration drift detection, change-management workflow (GitOps).
    \item \textbf{Observability}: Unified dashboards (Grafana) for latency, bandwidth, errors, CPU/memory, BGP state, certificate expiry; alerting via Email/Telegram/Slack with escalation.
    \item \textbf{Documentation}: Living diagrams (NetBox), IPAM (NetBox), runbooks for top 10 incidents, disaster-recovery plan tested quarterly.
\end{enumerate}

\medskip
\textbf{Constraints:}
\begin{itemize}
    \item Budget: 150 million XAF CAPEX (new hardware), 30 million XAF/year OPEX.
    \item Timeline: Design (2 weeks), Implementation in GNS3 (3 weeks), Automation \& Monitoring (2 weeks), Testing \& Documentation (2 weeks).
    \item Team: 4 students per group (roles: Network Architect, Automation Engineer, Security Analyst, Documentation Lead).
    \item Assessment: GCE Practical Project grid (Design 30\%, Implementation 30\%, Automation/Monitoring 20\%, Report/Presentation 20\%).
\end{itemize}
\end{experience}

\courssub{Tasks}
\begin{enumerate}
    \item \textbf{Complexity Analysis} — Produce a table mapping each \cle{complexity factor} (scale, heterogeneity, dynamics, administrative boundaries, policy diversity) to concrete manifestations in CampusNet. Rate impact (High/Medium/Low) and propose a mitigation strategy.
    \item \textbf{Network Design} —
    \begin{enumerate}
        \item Draw the \textbf{logical topology} (VLANs, routing domains, security zones) and \textbf{physical topology} (device placement, link types, redundancy).
        \item Define an \textbf{IP addressing plan} (IPv4 / IPv6) using VLSM: allocate /20 per campus, /24 per building, /30 for point-to-point, /64 per VLAN for IPv6. Document in a table.
        \item Select routing protocols: \textbf{OSPF} (Area 0 backbone, Area 1/2/3 per campus) for internal routing; \textbf{iBGP} between campus border routers for policy control; \textbf{eBGP} to ISPs. Justify.
        \item Design \textbf{high availability}: VRRP/HSRP for default gateways, LACP (802.3ad) for all inter-switch links, BFD for OSPF/BGP fast failure detection.
    \end{enumerate}
    \item \textbf{GNS3 Implementation} —
    \begin{enumerate}
        \item Build the topology in GNS3 using Cisco IOSv (routers), IOU/IOL (switches), FRR containers (Linux BGP), and Ubuntu servers (LMS, RADIUS, Syslog).
        \item Configure: VLANs, SVIs, OSPF (with area authentication), iBGP (route-reflectors), eBGP (with route-maps for prefix filtering), VRRP, 802.1X (RADIUS), ACLs for inter-zone traffic.
        \item Verify end-to-end connectivity, VoIP DSCP marking (EF), QoS policies (priority queue for voice, shaped queue for HPC replication).
    \end{enumerate}
    \item \textbf{Automation with Ansible} —
    \begin{enumerate}
        \item Create an Ansible inventory (YAML) grouped by campus, role (core, access, border, server).
        \item Write playbooks for: initial bootstrap (hostname, management IP, SSH keys, SNMPv3, NTP, syslog), VLAN/SVI deployment (Jinja2 template), OSPF/BGP configuration (templates with variables), VRRP/LACP, firewall ACLs.
        \item Implement \textbf{idempotency} checks and \textbf{configuration drift detection} (run playbook in check mode daily via GitLab CI).
    \end{enumerate}
    \item \textbf{Observability Stack Deployment} —
    \begin{enumerate}
        \item Deploy Prometheus (scrape configs for node-exporter, snmp-exporter, blackbox-exporter), Alertmanager (routes, inhibition rules, receivers), Grafana (dashboards: Campus Overview, Link Health, BGP Peers, Device Resources, Security Events).
        \item Define 10 critical alerts (e.g., \texttt{BGPPeerDown}, \texttt{HighLatency}, \texttt{LinkUtilization>80\%}, \texttt{CPU>90\%}, \texttt{CertificateExpiry<30d}) with severity labels and runbook links.
    \end{enumerate}
    \item \textbf{Resilience Testing} —
    \begin{enumerate}
        \item Simulate: (a) Primary WAN link failure (Yaoundé–Douala fibre cut), (b) Core switch reboot (Yaoundé), (c) BGP session flap (Douala–Orange), (d) HPC replication traffic surge (1\,Gbps for 10\,min).
        \item Measure: failover time, packet loss, latency, application impact (VoIP MOS, LMS page load). Capture pcaps and Grafana screenshots.
    \end{enumerate}
    \item \textbf{Technical Report \& Presentation} —
    \begin{enumerate}
        \item Write a 15-page report: Executive Summary, Requirements Analysis, Design Decisions (with diagrams/tables), Implementation Details (config snippets, playbook excerpts), Test Results (tables/graphs), Lessons Learned, Appendices (full configs, Ansible repo structure).
        \item Prepare a 10-minute slide deck (12 slides max) highlighting problem, solution, key innovations, demo of automation/monitoring, and Q\&A readiness.
    \end{enumerate}
\end{enumerate}

\courssub{Model answer}
\begin{exoresolu}[Task 1 — Complexity Analysis]
\textbf{Statement:} Produce a table mapping each complexity factor to concrete manifestations in CampusNet, rate impact, and propose mitigation.

\tcblower
\begin{center}
\begin{tabularx}{\linewidth}{|l|X|l|X|}
\hline
\rowcolor{popblueL}
\textbf{Factor} & \textbf{Manifestation in CampusNet} & \textbf{Impact} & \textbf{Mitigation Strategy} \\
\hline
Scale & 26\,500 users, 5\,000+ devices, 3 sites, 2\,PB storage & High & Hierarchical design (core/distribution/access), route summarisation, IPv6 /48 per campus \\
\hline
Heterogeneity & Cisco, FRR, Ubuntu, Windows, IoT (CoAP/MQTT), HPC (InfiniBand gateway) & High & Standardised OS images, NETCONF/YANG for config, abstraction via Ansible roles \\
\hline
Dynamics & Student BYOD churn (3\,000/day), VM migration, HPC job bursts, research data flows & High & BFD + fast-hello OSPF, ECMP, dynamic buffer allocation, telemetry streaming (gNMI) \\
\hline
Administrative Boundaries & MINESUP (policy), ANTIC (security), CAMTEL/Orange/MTN (SLA), Campus IT teams (autonomy) & Medium & Federation agreements, shared NetBox/IPAM, RBAC per campus, audit logs \\
\hline
Policy Diversity & QoS per app (VoIP EF, HPC AF31, Web BE), 802.1X vs PSK, encryption mandates (IPsec for inter-campus) & High & Class-based QoS with MQC, Cisco ISE / FreeRADIUS centralised, IPsec VTI templates \\
\hline
\end{tabularx}
\end{center}
\end{exoresolu}

\begin{exoresolu}[Task 2 — Network Design]
\textbf{Statement:} Logical/physical topology, IP addressing plan, routing protocol selection, HA design.

\tcblower
\textbf{2.1 Logical Topology (security zones and routing domains)}
\begin{popfigure}
\begin{tikzpicture}[node distance=2.5cm, every node/.style={font=\small}]
\node[draw, rounded corners, fill=popgreenL, thick] (core) {Core Zone (Transit)};
\node[draw, rounded corners, fill=popblueL, right=of core] (dmz) {DMZ / Public Services};
\node[draw, rounded corners, fill=poporange, below=of core] (admin) {Admin / Mgmt};
\node[draw, rounded corners, fill=poppink, below=of dmz] (user) {User Access (802.1X)};
\node[draw, rounded corners, fill=popgold, left=of user] (iot) {IoT / OT};
\node[draw, rounded corners, fill=poppurple, left=of iot] (hpc) {HPC / Research};
\draw[thick, ->] (core) -- node[midway, above] {OSPF Area 0} (dmz);
\draw[thick, ->] (core) -- (admin);
\draw[thick, ->] (core) -- (user);
\draw[thick, ->] (core) -- (iot);
\draw[thick, ->] (core) -- (hpc);
\node[draw, rounded corners, fill=poppink, above=of core] (wan) {WAN / ISP Peering};
\draw[thick, <->] (wan) -- node[midway, right] {iBGP / IPsec} (core);
\end{tikzpicture}
\legende{Logical security zones and routing domains}
\end{popfigure}

\textbf{2.2 Physical Topology}
\begin{itemize}
    \item Each campus: 2$\times$ Core (9500 stack) — 2$\times$ Distribution (9500) — Access (2960X) in redundant star.
    \item Inter-campus: Dual-homed border routers (ISR 4451) with LACP 2$\times$10G (Yaoundé–Douala), 2$\times$1G (Yaoundé–Buea), 1$\times$1G VPN backup (Douala–Buea).
    \item Out-of-band management: Dedicated VRF \texttt{MGMT} on separate 1G switch, console server.
\end{itemize}

\textbf{2.3 IP Addressing Plan (excerpt)}
\begin{center}
\begin{tabularx}{\linewidth}{|l|l|l|X|}
\hline
\rowcolor{popblueL}
\textbf{Segment} & \textbf{IPv4 (VLSM)} & \textbf{IPv6} & \textbf{Description} \\
\hline
Yaoundé Core Loopbacks & 10.0.0.0/22 & 2001:db8:1::/48 & OSPF Router-IDs, BGP source \\
\hline
Douala Core Loopbacks & 10.0.4.0/22 & 2001:db8:2::/48 & OSPF Router-IDs, BGP source \\
\hline
Buea Core Loopbacks & 10.0.8.0/22 & 2001:db8:3::/48 & OSPF Router-IDs, BGP source \\
\hline
PtP Links (Yaoundé–Douala) & 10.1.0.0/30, 10.1.0.4/30 & 2001:db8:0:100::/64, 2001:db8:0:101::/64 & LACP member links \\
\hline
PtP Links (Yaoundé–Buea) & 10.1.1.0/30, 10.1.1.4/30 & 2001:db8:0:200::/64, 2001:db8:0:201::/64 & LACP member links \\
\hline
Yaoundé User VLANs & 10.10.0.0/20 (/24 per building) & 2001:db8:1:1000::/56 (/64 per VLAN) & 802.1X dynamic VLAN assignment \\
\hline
Douala User VLANs & 10.20.0.0/20 & 2001:db8:2:1000::/56 & 802.1X dynamic VLAN assignment \\
\hline
Buea User VLANs & 10.30.0.0/20 & 2001:db8:3:1000::/56 & 802.1X dynamic VLAN assignment \\
\hline
Server Farm (Yaoundé) & 10.100.0.0/22 & 2001:db8:1:2000::/56 & LMS, ERP, DB, Syslog \\
\hline
DMZ (all campuses) & 10.200.0.0/22 & 2001:db8:0:3000::/56 & Public DNS, Web, VPN gateway \\
\hline
Mgmt VRF & 10.250.0.0/24 & 2001:db8:0:4000::/64 & OOB, NetBox, Ansible, Prometheus \\
\hline
\end{tabularx}
\end{center}

\textbf{2.4 Routing Protocol Justification}
\begin{itemize}
    \item \textbf{OSPF} (Area 0 = inter-campus core links; Area 1/2/3 = campus internal): fast convergence, hierarchical summarisation, vendor-agnostic. Authentication: SHA-256.
    \item \textbf{iBGP} (full mesh via Route Reflectors at each campus core): carries customer prefixes (campus subnets), applies local-pref for primary/backup path selection, enables traffic engineering (MED, communities).
    \item \textbf{eBGP} to ISPs: standard upstream peering, prefix filtering (max-prefix 100), RPKI validation.
    \item \textbf{Redistribution}: OSPF $\leftrightarrow$ iBGP only at border routers with route-maps (tagging, filtering).
\end{itemize}

\textbf{2.5 High Availability Mechanisms}
\begin{itemize}
    \item \textbf{Gateway redundancy}: VRRPv3 (IPv4/IPv6) on all SVIs, preempt delay 30\,s, tracking uplink interfaces.
    \item \textbf{Link aggregation}: LACP (active) on all inter-switch and server uplinks; min-links 2.
    \item \textbf{Fast failure detection}: BFD (interval 300\,ms, multiplier 3) for OSPF and BGP neighbours.
    \item \textbf{Control-plane policing}: CoPP on core routers to protect BGP/OSPF/SSH.
\end{itemize}
\end{exoresolu}

\begin{exoresolu}[Task 3 — GNS3 Implementation (Key Config Snippets)]
\textbf{Statement:} Build topology, configure VLANs, OSPF, BGP, VRRP, QoS, 802.1X.

\tcblower
\textbf{3.1 Core Switch VLAN/SVI (Catalyst 9500)}\par
\texttt{vlan 10,20,30,100,200}\\
\texttt{\ name MGMT, USER\_YAO, USER\_DLA, SERVER, DMZ}\\
\texttt{interface Vlan10}\\
\texttt{\ description MGMT\_VRF}\\
\texttt{\ vrf forwarding MGMT}\\
\texttt{\ ip address 10.250.0.1 255.255.255.0}\\
\texttt{\ ipv6 address 2001:db8:0:4000::1/64}\\
\texttt{\ no shutdown}\\
\texttt{interface Vlan100}\\
\texttt{\ description SERVER\_FARM}\\
\texttt{\ ip address 10.100.0.1 255.255.252.0}\\
\texttt{\ ipv6 address 2001:db8:1:2000::1/56}\\
\texttt{\ vrrp 1 address-family ipv4}\\
\texttt{\ \ address 10.100.0.1}\\
\texttt{\ \ priority 110}\\
\texttt{\ \ preempt delay minimum 30}\\
\texttt{\ \ track 1 decrement 20}\\
\texttt{\ vrrp 1 address-family ipv6}\\
\texttt{\ \ address 2001:db8:1:2000::1}\\
\texttt{\ \ priority 110}\\
\texttt{\ \ preempt delay minimum 30}\\
\texttt{\ no shutdown}

\medskip
\textbf{3.2 OSPF with BFD (Router Yaoundé-Core1)}\par
\texttt{router ospf 1}\\
\texttt{\ router-id 10.0.0.1}\\
\texttt{\ area 0 authentication message-digest}\\
\texttt{\ passive-interface default}\\
\texttt{\ no passive-interface TenGigabitEthernet1/0/1}\\
\texttt{\ no passive-interface TenGigabitEthernet1/0/2}\\
\texttt{\ network 10.0.0.0 0.0.3.255 area 0}\\
\texttt{\ network 10.1.0.0 0.0.0.3 area 0}\\
\texttt{\ network 10.1.1.0 0.0.0.3 area 0}\\
\texttt{\ bfd all-interfaces}\\
\texttt{interface TenGigabitEthernet1/0/1}\\
\texttt{\ description Link\_to\_Douala\_Core1}\\
\texttt{\ ip address 10.1.0.1 255.255.255.252}\\
\texttt{\ ip ospf network point-to-point}\\
\texttt{\ ip ospf authentication message-digest}\\
\texttt{\ ip ospf message-digest-key 1 sha256 0 <strong-key>}\\
\texttt{\ bfd interval 300 min\_rx 300 multiplier 3}\\
\texttt{\ no shutdown}

\medskip
\textbf{3.3 iBGP with Route Reflector (Yaoundé Core1)}\par
\texttt{router bgp 65001}\\
\texttt{\ bgp router-id 10.0.0.1}\\
\texttt{\ bgp log-neighbor-changes}\\
\texttt{\ no bgp default ipv4-unicast}\\
\texttt{\ neighbor RR\_CLIENT peer-group}\\
\texttt{\ neighbor RR\_CLIENT remote-as 65001}\\
\texttt{\ neighbor RR\_CLIENT update-source Loopback0}\\
\texttt{\ neighbor RR\_CLIENT bfd}\\
\texttt{\ neighbor 10.0.0.2 peer-group RR\_CLIENT}\\
\texttt{\ neighbor 10.0.4.1 peer-group RR\_CLIENT\ \ ! Douala RR}\\
\texttt{\ neighbor 10.0.8.1 peer-group RR\_CLIENT\ \ ! Buea RR}\\
\texttt{\ address-family ipv4 unicast}\\
\texttt{\ \ neighbor RR\_CLIENT activate}\\
\texttt{\ \ neighbor RR\_CLIENT route-reflector-client}\\
\texttt{\ \ neighbor RR\_CLIENT send-community both}\\
\texttt{\ \ network 10.10.0.0 mask 255.255.240.0}\\
\texttt{\ \ network 10.100.0.0 mask 255.255.252.0}\\
\texttt{\ \ maximum-paths ibgp 4}\\
\texttt{\ exit-address-family}\\
\texttt{\ address-family ipv6 unicast}\\
\texttt{\ \ neighbor RR\_CLIENT activate}\\
\texttt{\ \ neighbor RR\_CLIENT route-reflector-client}\\
\texttt{\ \ network 2001:db8:1::/48}\\
\texttt{\ exit-address-family}

\medskip
\textbf{3.4 QoS for VoIP \& HPC (MQC)}\par
\texttt{class-map match-any VOICE}\\
\texttt{\ match dscp ef}\\
\texttt{class-map match-any HPC\_REPL}\\
\texttt{\ match dscp af31}\\
\texttt{policy-map CAMPUS\_QOS}\\
\texttt{\ class VOICE}\\
\texttt{\ \ priority percent 20}\\
\texttt{\ \ police cir 200000000 conform-action transmit exceed-action drop}\\
\texttt{\ class HPC\_REPL}\\
\texttt{\ \ bandwidth remaining percent 50}\\
\texttt{\ \ random-detect dscp-based}\\
\texttt{\ class class-default}\\
\texttt{\ \ fair-queue}\\
\texttt{\ \ random-detect}\\
\texttt{interface TenGigabitEthernet1/0/1}\\
\texttt{\ service-policy output CAMPUS\_QOS}

\medskip
\textbf{3.5 802.1X (Access Switch Port Template)}\par
\texttt{interface range GigabitEthernet1/0/1 - 48}\\
\texttt{\ description USER\_ACCESS}\\
\texttt{\ switchport mode access}\\
\texttt{\ authentication port-control auto}\\
\texttt{\ dot1x pae authenticator}\\
\texttt{\ mab}\\
\texttt{\ dot1x timeout tx-period 10}\\
\texttt{\ spanning-tree portfast}\\
\texttt{\ spanning-tree bpduguard enable}\\
\texttt{radius server ISE\_YAO}\\
\texttt{\ address ipv4 10.100.10.10 auth-port 1812 acct-port 1813}\\
\texttt{\ key 7 <radius-shared-key>}\\
\texttt{aaa group server radius RADIUS\_GRP}\\
\texttt{\ server name ISE\_YAO}\\
\texttt{\ server name ISE\_DLA}\\
\texttt{\ server name ISE\_BUE}\\
\texttt{aaa authentication dot1x default group RADIUS\_GRP}\\
\texttt{aaa authorization network default group RADIUS\_GRP}\\
\texttt{aaa accounting dot1x default start-stop group RADIUS\_GRP}
\end{exoresolu}

\begin{exoresolu}[Task 4 — Ansible Automation]
\textbf{Statement:} Inventory, playbooks for bootstrap, VLAN/SVI, routing, drift detection.

\tcblower
\textbf{4.1 Inventory (\texttt{inventory/campusnet.yml})}\par
\texttt{all:}\\
\texttt{\ \ vars:}\\
\texttt{\ \ \ \ ansible\_network\_os: ios}\\
\texttt{\ \ \ \ ansible\_connection: network\_cli}\\
\texttt{\ \ \ \ ansible\_user: "\{\{ vault\_ansible\_user \}\}"}\\
\texttt{\ \ \ \ ansible\_become: yes}\\
\texttt{\ \ \ \ ansible\_become\_method: enable}\\
\texttt{\ \ \ \ ntp\_servers: [10.250.0.10, 10.250.0.11]}\\
\texttt{\ \ \ \ snmp\_v3\_user: campusnet\_mon}\\
\texttt{\ \ children:}\\
\texttt{\ \ \ \ yaounde:}\\
\texttt{\ \ \ \ \ \ vars: \{ campus: yaounde, asn: 65001, ipv4\_base: 10.0.0.0/22,}\\
\texttt{\ \ \ \ \ \ \ \ \ \ \ \ \ \ ipv6\_base: 2001:db8:1::/48 \}}\\
\texttt{\ \ \ \ \ \ children:}\\
\texttt{\ \ \ \ \ \ \ \ core:\ \ hosts: \{ yao-core1: \{ansible\_host: 10.250.0.1\},}\\
\texttt{\ \ \ \ \ \ \ \ \ \ \ \ \ \ \ \ \ \ \ yao-core2: \{ansible\_host: 10.250.0.2\} \}}\\
\texttt{\ \ \ \ \ \ \ \ border: hosts: \{ yao-br1: \{ansible\_host: 10.250.0.3\} \}}\\
\texttt{\ \ \ \ douala:}\\
\texttt{\ \ \ \ \ \ vars: \{ campus: douala, asn: 65002, ipv4\_base: 10.0.4.0/22,}\\
\texttt{\ \ \ \ \ \ \ \ \ \ \ \ \ \ ipv6\_base: 2001:db8:2::/48 \}}\\
\texttt{\ \ \ \ \ \ children:}\\
\texttt{\ \ \ \ \ \ \ \ core:\ \ hosts: \{ dla-core1: \{ansible\_host: 10.250.0.21\} \}}\\
\texttt{\ \ \ \ \ \ \ \ border: hosts: \{ dla-br1: \{ansible\_host: 10.250.0.23\} \}}\\
\texttt{\ \ \ \ buea:}\\
\texttt{\ \ \ \ \ \ vars: \{ campus: buea, asn: 65003, ipv4\_base: 10.0.8.0/22,}\\
\texttt{\ \ \ \ \ \ \ \ \ \ \ \ \ \ ipv6\_base: 2001:db8:3::/48 \}}\\
\texttt{\ \ \ \ \ \ children:}\\
\texttt{\ \ \ \ \ \ \ \ core:\ \ hosts: \{ bue-core1: \{ansible\_host: 10.250.0.41\} \}}\\
\texttt{\ \ \ \ \ \ \ \ border: hosts: \{ bue-br1: \{ansible\_host: 10.250.0.43\} \}}

\medskip
\textbf{4.2 Bootstrap Playbook (\texttt{playbooks/bootstrap.yml})}\par
\texttt{- name: Bootstrap network devices}\\
\texttt{\ \ hosts: all}\\
\texttt{\ \ gather\_facts: no}\\
\texttt{\ \ tasks:}\\
\texttt{\ \ \ \ - name: Set hostname}\\
\texttt{\ \ \ \ \ \ ios\_config:}\\
\texttt{\ \ \ \ \ \ \ \ lines: "hostname \{\{ inventory\_hostname \}\}"}\\
\texttt{\ \ \ \ \ \ \ \ save\_when: changed}\\
\texttt{\ \ \ \ - name: Enable SSH v2, disable Telnet}\\
\texttt{\ \ \ \ \ \ ios\_config:}\\
\texttt{\ \ \ \ \ \ \ \ lines:}\\
\texttt{\ \ \ \ \ \ \ \ \ \ - "ip ssh version 2"}\\
\texttt{\ \ \ \ \ \ \ \ \ \ - "line vty 0 15"}\\
\texttt{\ \ \ \ \ \ \ \ \ \ - " transport input ssh"}\\
\texttt{\ \ \ \ \ \ \ \ \ \ - " login local"}\\
\texttt{\ \ \ \ \ \ \ \ save\_when: changed}\\
\texttt{\ \ \ \ - name: Configure NTP}\\
\texttt{\ \ \ \ \ \ ios\_config:}\\
\texttt{\ \ \ \ \ \ \ \ lines: "ntp server \{\{ item \}\} vrf MGMT"}\\
\texttt{\ \ \ \ \ \ \ \ loop: "\{\{ ntp\_servers \}\}"}\\
\texttt{\ \ \ \ \ \ \ \ save\_when: changed}\\
\texttt{\ \ \ \ - name: Configure SNMPv3}\\
\texttt{\ \ \ \ \ \ ios\_config:}\\
\texttt{\ \ \ \ \ \ \ \ lines:}\\
\texttt{\ \ \ \ \ \ \ \ \ \ - "snmp-server user \{\{ snmp\_v3\_user \}\} netadmin v3 auth sha}\\
\texttt{\ \ \ \ \ \ \ \ \ \ \ \ \{\{ vault\_snmp\_auth \}\} priv aes 256 \{\{ vault\_snmp\_priv \}\}"}\\
\texttt{\ \ \ \ \ \ \ \ \ \ - "snmp-server host 10.250.0.100 version 3 auth \{\{ snmp\_v3\_user \}\}"}\\
\texttt{\ \ \ \ \ \ \ \ save\_when: changed}\\
\texttt{\ \ \ \ - name: Configure Syslog}\\
\texttt{\ \ \ \ \ \ ios\_config:}\\
\texttt{\ \ \ \ \ \ \ \ lines:}\\
\texttt{\ \ \ \ \ \ \ \ \ \ - "logging host 10.250.0.100 transport tcp port 6514"}\\
\texttt{\ \ \ \ \ \ \ \ \ \ - "logging trap notifications"}\\
\texttt{\ \ \ \ \ \ \ \ save\_when: changed}

\medskip
\textbf{4.3 VLAN/SVI Deployment (Jinja2 Template — \texttt{templates/vlans.j2})}\par
\texttt{\{\% for vlan in vlans \%\}}\\
\texttt{vlan \{\{ vlan.id \}\}}\\
\texttt{\ name \{\{ vlan.name \}\}}\\
\texttt{interface Vlan\{\{ vlan.id \}\}}\\
\texttt{\ description \{\{ vlan.description \}\}}\\
\texttt{\ ip address \{\{ vlan.ipv4 \}\} \{\{ vlan.mask \}\}}\\
\texttt{\ \{\% if vlan.ipv6 \%\}ipv6 address \{\{ vlan.ipv6 \}\}\{\% endif \%\}}\\
\texttt{\ \{\% if vlan.vrrp \%\}}\\
\texttt{\ vrrp \{\{ vlan.vrrp.group \}\} address-family ipv4}\\
\texttt{\ \ address \{\{ vlan.vrrp.vip \}\}}\\
\texttt{\ \ priority \{\{ vlan.vrrp.priority \}\}}\\
\texttt{\ \ preempt delay minimum 30}\\
\texttt{\ \{\% endif \%\}}\\
\texttt{\ no shutdown}\\
\texttt{\{\% endfor \%\}}

\medskip
\textbf{4.4 Drift Detection (GitLab CI — \texttt{.gitlab-ci.yml} excerpt)}\par
\texttt{stages: [validate, deploy, drift-check]}\\
\texttt{validate\_config:}\\
\texttt{\ \ stage: validate}\\
\texttt{\ \ script:}\\
\texttt{\ \ \ \ - pip install ansible ansible-lint}\\
\texttt{\ \ \ \ - ansible-lint playbooks/}\\
\texttt{\ \ \ \ - ansible-playbook -i inventory/campusnet.yml playbooks/site.yml {-}{-}syntax-check}\\
\texttt{drift\_detection:}\\
\texttt{\ \ stage: drift-check}\\
\texttt{\ \ only: [schedules]}\\
\texttt{\ \ script:}\\
\texttt{\ \ \ \ - ansible-playbook -i inventory/campusnet.yml playbooks/site.yml {-}{-}check {-}{-}diff}\\
\texttt{\ \ \ \ - \ \# fail the pipeline if any task reports "changed"}
\end{exoresolu}

\begin{exoresolu}[Task 5 — Observability Stack]
\textbf{Statement:} Deploy Prometheus, Alertmanager, Grafana; define 10 critical alerts.

\tcblower
\textbf{5.1 Docker Compose (\texttt{monitoring/docker-compose.yml})}\par
\texttt{version: '3.8'}\\
\texttt{services:}\\
\texttt{\ \ prometheus:}\\
\texttt{\ \ \ \ image: prom/prometheus:v2.48.0}\\
\texttt{\ \ \ \ volumes: [./prometheus.yml:/etc/prometheus/prometheus.yml,}\\
\texttt{\ \ \ \ \ \ \ \ \ \ \ \ \ \ ./rules:/etc/prometheus/rules, prom\_data:/prometheus]}\\
\texttt{\ \ \ \ ports: ["9090:9090"]}\\
\texttt{\ \ alertmanager:}\\
\texttt{\ \ \ \ image: prom/alertmanager:v0.26.0}\\
\texttt{\ \ \ \ volumes: [./alertmanager.yml:/etc/alertmanager/alertmanager.yml]}\\
\texttt{\ \ \ \ ports: ["9093:9093"]}\\
\texttt{\ \ grafana:}\\
\texttt{\ \ \ \ image: grafana/grafana:10.2.0}\\
\texttt{\ \ \ \ volumes: [./grafana:/etc/grafana/provisioning, grafana\_data:/var/lib/grafana]}\\
\texttt{\ \ \ \ ports: ["3000:3000"]}\\
\texttt{\ \ snmp-exporter:}\\
\texttt{\ \ \ \ image: prom/snmp-exporter:v0.25.0}\\
\texttt{\ \ \ \ ports: ["9116:9116"]}\\
\texttt{\ \ blackbox-exporter:}\\
\texttt{\ \ \ \ image: prom/blackbox-exporter:v0.24.0}\\
\texttt{\ \ \ \ ports: ["9115:9115"]}\\
\texttt{volumes: \{ prom\_data: \{\}, grafana\_data: \{\} \}}

\medskip
\textbf{5.2 Prometheus Scrape Config (\texttt{prometheus.yml} excerpt)}\par
\texttt{global: \{ scrape\_interval: 15s, evaluation\_interval: 15s \}}\\
\texttt{alerting:}\\
\texttt{\ \ alertmanagers: [ static\_configs: [ targets: ['alertmanager:9093'] ] ]}\\
\texttt{rule\_files: ['rules/*.yml']}\\
\texttt{scrape\_configs:}\\
\texttt{\ \ - job\_name: 'network-devices'}\\
\texttt{\ \ \ \ metrics\_path: /snmp}\\
\texttt{\ \ \ \ params: \{ module: [if\_mib, bgp4\_mib] \}}\\
\texttt{\ \ \ \ static\_configs:}\\
\texttt{\ \ \ \ \ \ - targets: [10.250.0.1, 10.250.0.2, 10.250.0.21, 10.250.0.41]}\\
\texttt{\ \ \ \ relabel\_configs:}\\
\texttt{\ \ \ \ \ \ - \{ source\_labels: [\_\_address\_\_], target\_label: \_\_param\_target \}}\\
\texttt{\ \ \ \ \ \ - \{ source\_labels: [\_\_param\_target], target\_label: instance \}}\\
\texttt{\ \ \ \ \ \ - \{ target\_label: \_\_address\_\_, replacement: snmp-exporter:9116 \}}\\
\texttt{\ \ - job\_name: 'blackbox-icmp'}\\
\texttt{\ \ \ \ metrics\_path: /probe}\\
\texttt{\ \ \ \ params: \{ module: [icmp] \}}\\
\texttt{\ \ \ \ static\_configs:}\\
\texttt{\ \ \ \ \ \ - targets: [10.1.0.2, 10.1.1.2, 2001:db8:2::1, 2001:db8:3::1]}\\
\texttt{\ \ \ \ relabel\_configs:}\\
\texttt{\ \ \ \ \ \ - \{ source\_labels: [\_\_address\_\_], target\_label: \_\_param\_target \}}\\
\texttt{\ \ \ \ \ \ - \{ target\_label: \_\_address\_\_, replacement: blackbox-exporter:9115 \}}

\medskip
\textbf{5.3 Alert Rules (\texttt{rules/campusnet.yml})}\par
\texttt{groups:}\\
\texttt{- name: campusnet\_critical}\\
\texttt{\ \ interval: 30s}\\
\texttt{\ \ rules:}\\
\texttt{\ \ - alert: BGPPeerDown}\\
\texttt{\ \ \ \ expr: bgp\_peer\_state\{state="established"\} == 0}\\
\texttt{\ \ \ \ for: 1m}\\
\texttt{\ \ \ \ labels: \{ severity: critical, team: network \}}\\
\texttt{\ \ - alert: HighLatencyInterCampus}\\
\texttt{\ \ \ \ expr: probe\_duration\_seconds > 0.015}\\
\texttt{\ \ \ \ for: 2m}\\
\texttt{\ \ \ \ labels: \{ severity: warning, team: network \}}\\
\texttt{\ \ - alert: LinkUtilizationHigh}\\
\texttt{\ \ \ \ expr: (rate(ifHCInOctets[5m]) + rate(ifHCOutOctets[5m])) * 8}\\
\texttt{\ \ \ \ \ \ \ \ \ \ / ifHighSpeed / 1e6 > 0.8}\\
\texttt{\ \ \ \ for: 5m}\\
\texttt{\ \ \ \ labels: \{ severity: warning, team: network \}}\\
\texttt{\ \ - alert: CPUHigh}\\
\texttt{\ \ \ \ expr: cisco\_cpu\_5sec > 90}\\
\texttt{\ \ \ \ for: 5m}\\
\texttt{\ \ \ \ labels: \{ severity: warning, team: network \}}\\
\texttt{\ \ - alert: MemoryHigh}\\
\texttt{\ \ \ \ expr: (cisco\_memory\_used / cisco\_memory\_total) > 0.85}\\
\texttt{\ \ \ \ for: 5m}\\
\texttt{\ \ \ \ labels: \{ severity: warning, team: network \}}\\
\texttt{\ \ - alert: VRRPStateChange}\\
\texttt{\ \ \ \ expr: changes(vrrp\_state[10m]) > 0}\\
\texttt{\ \ \ \ labels: \{ severity: critical, team: network \}}\\
\texttt{\ \ - alert: CertificateExpirySoon}\\
\texttt{\ \ \ \ expr: ssl\_cert\_not\_after - time() < 30 * 86400}\\
\texttt{\ \ \ \ labels: \{ severity: warning, team: security \}}\\
\texttt{\ \ - alert: HPCReplicationLag}\\
\texttt{\ \ \ \ expr: hpc\_replication\_lag\_seconds > 300}\\
\texttt{\ \ \ \ for: 10m}\\
\texttt{\ \ \ \ labels: \{ severity: critical, team: research \}}\\
\texttt{\ \ - alert: RadiusAuthFailures}\\
\texttt{\ \ \ \ expr: rate(radius\_auth\_failures\_total[5m]) > 10}\\
\texttt{\ \ \ \ for: 2m}\\
\texttt{\ \ \ \ labels: \{ severity: warning, team: security \}}\\
\texttt{\ \ - alert: ConfigDriftDetected}\\
\texttt{\ \ \ \ expr: ansible\_drift\_detected == 1}\\
\texttt{\ \ \ \ labels: \{ severity: warning, team: automation \}}

\medskip
\textbf{5.4 Alertmanager Config (\texttt{alertmanager.yml})}\par
\texttt{route:}\\
\texttt{\ \ group\_by: ['alertname', 'campus', 'severity']}\\
\texttt{\ \ group\_wait: 30s}\\
\texttt{\ \ group\_interval: 5m}\\
\texttt{\ \ repeat\_interval: 4h}\\
\texttt{\ \ receiver: 'default'}\\
\texttt{\ \ routes:}\\
\texttt{\ \ \ \ - \{ match: \{ severity: critical \}, receiver: 'critical-pager', continue: true \}}\\
\texttt{\ \ \ \ - \{ match: \{ team: security \}, receiver: 'security-team' \}}\\
\texttt{receivers:}\\
\texttt{\ \ - name: 'default'}\\
\texttt{\ \ \ \ email\_configs: [ \{ to: 'noc@campusnet.cm', send\_resolved: true \} ]}\\
\texttt{\ \ - name: 'critical-pager'}\\
\texttt{\ \ \ \ email\_configs: [ \{ to: 'oncall-net@campusnet.cm', send\_resolved: true \} ]}\\
\texttt{\ \ - name: 'security-team'}\\
\texttt{\ \ \ \ email\_configs: [ \{ to: 'secops@campusnet.cm', send\_resolved: true \} ]}\\
\texttt{inhibit\_rules:}\\
\texttt{\ \ - source\_match: \{ alertname: 'BGPPeerDown' \}}\\
\texttt{\ \ \ \ target\_match\_re: \{ alertname: 'LinkUtilizationHigh|HighLatencyInterCampus' \}}\\
\texttt{\ \ \ \ equal: ['instance']}

\medskip
\textbf{5.5 Grafana Dashboard — key panels}
\begin{itemize}
    \item \textbf{Stat panel}: BGP Peers Up/Down (query: \texttt{sum(bgp\_peer\_state) by (state)}).
    \item \textbf{Time series}: Inter-campus latency P50/P95/P99 (Blackbox ICMP).
    \item \textbf{Heatmap}: Interface utilisation \% (SNMP ifHCIn/OutOctets).
    \item \textbf{Table}: Top 10 CPU/Memory devices.
    \item \textbf{Geomap}: Campus sites with link health (green/yellow/red).
    \item \textbf{Logs panel}: Loki query for \texttt{authentication failure} on RADIUS.
\end{itemize}
\end{exoresolu}

\begin{exoresolu}[Task 6 — Resilience Testing]
\textbf{Statement:} Simulate failures, measure failover, packet loss, application impact.

\tcblower
\textbf{Test Plan \& Results Summary}
\begin{center}
\begin{tabularx}{\linewidth}{|l|c|c|c|c|X|}
\hline
\rowcolor{popblueL}
\textbf{Scenario} & \textbf{Failover Time} & \textbf{Packet Loss} & \textbf{Latency Spike} & \textbf{VoIP MOS} & \textbf{Observations} \\
\hline
Yaoundé–Douala fibre cut (primary 2G) & 320 ms (BFD) & 0.1\% (2 pkts) & 8 ms & 4.3 $\to$ 4.1 & Traffic shifts to Orange MPLS (1G) via iBGP local-pref; ECMP rebalances in 1.2\,s \\
\hline
Yaoundé Core1 reload & 410 ms (VRRP + BFD) & 0.3\% & 12 ms & 4.3 $\to$ 3.9 & Core2 becomes VRRP master; OSPF SPF recompute 180\,ms; BGP graceful restart preserves sessions \\
\hline
Douala–Orange BGP flap (simulated) & 1.8 s (BGP hold timer) & 1.2\% & 45 ms & 4.3 $\to$ 3.5 & Route reflector withdraws prefixes; backup VPN (MTN) takes over after BGP convergence \\
\hline
HPC replication surge (1\,Gbps, 10\,min) & N/A & 0\% (QoS) & 2 ms & 4.3 (stable) & AF31 queue absorbs burst; voice EF queue untouched; buffer usage peaks at 68\% \\
\hline
\end{tabularx}
\end{center}

\textbf{Key Measurements (Grafana screenshots referenced in report):}
\begin{itemize}
    \item \textbf{Failover time}: measured by ICMP ping stream (100 pps) from Yaoundé server to Douala server; first successful reply after failure.
    \item \textbf{Packet loss}: counted via \texttt{ping -c 1000} during a 30\,s window around failure.
    \item \textbf{VoIP MOS}: simulated with a SIPp scenario (G.711, 20\,ms ptime); MOS computed from the R-factor.
    \item \textbf{HPC replication}: \texttt{rsync} of a 50\,GB dataset; throughput monitored via \texttt{ifstat} and Prometheus \texttt{rate(ifHCOutOctets)}.
\end{itemize}

\textbf{Remediation Actions Identified:}
\begin{enumerate}
    \item Reduce BGP hold timer to 9\,s (keepalive 3\,s) for faster ISP failover.
    \item Enable BGP Graceful Restart on all border routers.
    \item Tune VRRP advertisement interval (sub-second timers) for faster gateway failover.
    \item Add \texttt{bfd all-interfaces} to OSPF on access uplinks.
\end{enumerate}
\end{exoresolu}

\begin{exoresolu}[Task 7 — Technical Report \& Presentation Outline]
\textbf{Statement:} 15-page report structure and 12-slide presentation.

\tcblower
\textbf{7.1 Report Structure (15 pages)}
\begin{enumerate}
    \item \textbf{Title Page} — Project name, team members, matricule numbers, supervisor, date.
    \item \textbf{Executive Summary} (1 page) — Problem, solution, key results, compliance with MINESUP requirements.
    \item \textbf{Requirements Analysis} (2 pages) — Complexity factor table (Task 1), operational challenges mapping, constraint matrix.
    \item \textbf{Network Design} (4 pages) —
          \begin{itemize}
              \item Logical topology diagram + security zone description.
              \item Physical topology diagram + bill of materials (BoM) with costs (XAF).
              \item IP addressing plan (full table, IPv4/IPv6).
              \item Routing design: OSPF areas, iBGP/eBGP policies, route maps, communities.
              \item HA design: VRRP, LACP, BFD, CoPP.
              \item QoS design: class maps, policy maps, DSCP mapping table.
              \item Security design: 802.1X, RADIUS, ACLs, IPsec VTI, segmentation.
          \end{itemize}
    \item \textbf{Implementation Details} (3 pages) —
          \begin{itemize}
              \item GNS3 topology screenshot, device list, versions.
              \item Representative configuration snippets (core, border, access) with annotations.
              \item Ansible repo structure (tree), key playbooks, Jinja2 templates, vault usage.
              \item Monitoring stack architecture, scrape configs, alert rules, dashboard screenshots.
          \end{itemize}
    \item \textbf{Testing \& Validation} (2 pages) — Test plan table, results table (Task 6), pcap analysis highlights, Grafana alert screenshots.
    \item \textbf{Operational Procedures} (1.5 pages) — Runbooks for top 5 alerts, change management workflow (GitOps), backup/restore, disaster recovery (RTO/RPO).
    \item \textbf{Lessons Learned \& Future Work} (1 page) — Automation challenges, GNS3 limitations, IPv6 deployment gaps, proposed SD-WAN migration.
    \item \textbf{Appendices} (0.5 pages) — Full device configs (GitLab link), Ansible playbook index, Prometheus rule files, glossary.
\end{enumerate}

\textbf{7.2 Presentation Slides (12 slides, 10 minutes)}
\begin{enumerate}
    \item \textbf{Title} — CampusNet: Inter-Campus Network for Cameroon Universities.
    \item \textbf{Context \& Challenges} — 3 campuses, 26k users, MINESUP requirements, complexity factors.
    \item \textbf{Design Philosophy} — Hierarchical, Zero-Trust, GitOps, Observability-first.
    \item \textbf{Logical \& Physical Topology} — Diagram with security zones, link types, redundancy.
    \item \textbf{IP Plan \& Routing} — VLSM table, OSPF areas, iBGP RR, eBGP policies.
    \item \textbf{High Availability \& QoS} — VRRP/BFD/LACP, MQC for VoIP/HPC.
    \item \textbf{Automation with Ansible} — Inventory, bootstrap, VLAN/SVI templates, drift detection CI.
    \item \textbf{Security Implementation} — 802.1X, RADIUS, ACLs, IPsec, SNMPv3, Syslog.
    \item \textbf{Observability Stack} — Prometheus/Alertmanager/Grafana, 10 critical alerts, dashboard demo.
    \item \textbf{Resilience Test Results} — Table + Grafana screenshots, failover $<$ 500\,ms achieved.
    \item \textbf{Compliance \& Cost} — Requirements traceability matrix, BoM within 150M XAF.
    \item \textbf{Conclusion \& Q\&A} — Key achievements, next steps (SD-WAN, Intent-Based Networking).
\end{enumerate}

\textbf{7.3 GCE Marking Grid Alignment}
\begin{center}
\begin{tabularx}{\linewidth}{|l|X|c|}
\hline
\rowcolor{popblueL}
\textbf{Criterion} & \textbf{Evidence in Deliverables} & \textbf{Weight} \\
\hline
Design & Topology diagrams, IP plan, routing/HA/QoS/security rationale, BoM & 30\% \\
Implementation & GNS3 running config, connectivity proof, QoS/VoIP demo, 802.1X capture & 30\% \\
Automation \& Monitoring & Ansible playbooks (idempotent, vault), Prometheus/Grafana/Alertmanager, drift CI & 20\% \\
Report \& Presentation & 15-page report structure, 10-min slides, traceability, lessons learned & 20\% \\
\hline
\end{tabularx}
\end{center}
\end{exoresolu}

\begin{attention}[Common Pitfalls in CampusNet Project]
\begin{itemize}
    \item \textbf{Forgetting IPv6} — MINESUP requires dual-stack; omitting IPv6 addressing/OSPFv3/BGP address-family loses marks.
    \item \textbf{Static routes instead of dynamic routing} — Not scalable; you must justify OSPF/iBGP/eBGP.
    \item \textbf{No authentication on routing protocols} — Plain OSPF/BGP is a security fail; use SHA-256 authentication.
    \item \textbf{Ansible playbooks not idempotent} — Using \texttt{ios\_command} instead of \texttt{ios\_config} modules; no \texttt{{-}{-}check} support.
    \item \textbf{Monitoring only up/down} — Must include latency, utilisation, BGP state, certificate expiry, application metrics.
    \item \textbf{Alert fatigue} — Too many alerts without grouping/inhibition; no runbook links.
    \item \textbf{Incomplete resilience tests} — Only testing link failover, not device reload, congestion, or security events.
    \item \textbf{Report exceeds 15 pages or lacks executive summary} — GCE examiners penalise non-compliance.
    \item \textbf{Presentation reads slides} — You must demonstrate a live dashboard, an Ansible run, or a failover video.
\end{itemize}
\end{attention}

\begin{savaistu}[Cameroon Context: CampusNet in Real Life]
The interconnection of Cameroonian public universities (Yaoundé I and II, Douala, Buea, Dschang, Ngaoundéré, Maroua) through the national research and education network, supported by CAMTEL fibre infrastructure and supervised by ANTIC for security, follows a very similar architecture: a national backbone, eduroam federation, shared computing resources, and centralised monitoring. Your CampusNet project mirrors the actual technical and governance challenges faced by Cameroonian network engineers — making this activity a genuine pre-career simulation.
\end{savaistu}

\newpage

\coursec{Methods and worked examples}

\coursec{Factors of Complexity in Modern Networks}

\courssub{Analysing Network Complexity Factors}

\begin{methode}[Analysing Network Complexity Factors]
\textbf{Objective:} Systematically identify and categorise the factors that increase complexity in a modern network infrastructure.
\begin{enumerate}
    \item \textbf{Inventory the Physical \& Logical Topology:} Map all sites (HQ, branches, data centres, cloud), links (MPLS, VPN, SD-WAN, 4G/5G backup), and devices (routers, switches, firewalls, APs, controllers).
    \item \textbf{Assess Heterogeneity:} Count distinct OS types (IOS, NX-OS, Linux, proprietary), vendor platforms, and hardware generations. High heterogeneity $\rightarrow$ higher ops overhead.
    \item \textbf{Evaluate Virtualisation \& Overlay Layers:} Identify VRFs, VXLAN/EVPN fabrics, SDN controllers (Cisco DNA Center, Meraki, FortiManager), and network functions virtualisation (NFV) instances.
    \item \textbf{Quantify Scale \& Dynamics:} Number of endpoints, VLANs, routing prefixes, frequency of changes (adds/moves/changes per week), and IoT/OT device count.
    \item \textbf{Classify Complexity Level:} Use a simple scoring matrix (Low/Medium/High) for each factor; aggregate to justify automation/orchestration investment.
\end{enumerate}
\textbf{Key Reflex:} Complexity is not just ``many devices''; it is \emph{state explosion} (routing tables, policy rules, MAC tables) multiplied by \emph{rate of change}.
\end{methode}

\begin{exoresolu}[Complexity Assessment for a Cameroonian Bank]
\textbf{Scenario:} \textit{Banque Atlantique Cameroon (BAC) operates a headquarters in Douala, 15 branches nationwide, and a disaster recovery (DR) site in Yaound\'e. The WAN uses MPLS (CAMTEL) with IPsec VPN backup over MTN/Orange 4G. The core runs Cisco Nexus 9K (NX-OS) with VXLAN/EVPN; branches run ISR 4000 (IOS-XE). A new SD-WAN overlay (Viptela) is being piloted at 5 branches. The network carries 120 VLANs, 5 VRFs, and 3,500 endpoints. Weekly change requests average 40. IoT includes 200 IP cameras and 50 smart ATMs per branch.}

\textbf{Questions:}
\begin{enumerate}
    \item Using the method above, produce a complexity scorecard (Low/Medium/High) for each of the five factors. Justify each rating.
    \item Calculate an overall complexity index (simple average of the five scores, mapping Low=1, Medium=2, High=3). Interpret the result for the CIO.
    \item Recommend \textbf{two} specific automation strategies to reduce operational friction.
\end{enumerate}
\tcblower
\textbf{Detailed Solution}

\noindent\textbf{1. Complexity Scorecard}

\begin{center}
\begin{tabularx}{\linewidth}{|l|X|X|c|}
\hline
\rowcolor{popblueL}
\textbf{Factor} & \textbf{Observation} & \textbf{Justification} & \textbf{Rating} \\
\hline
Physical/Logical Topology & Multi-site (17 sites), Hybrid WAN (MPLS + 4G VPN), Dual DC & Multiple transport types, asymmetric routing risk, DR sync requirements & \textbf{High} \\
\hline
Heterogeneity & 2 OS families (NX-OS, IOS-XE), 2 Hardware platforms, SD-WAN controller (vManage) & Distinct CLI/config models, separate upgrade cycles, skill fragmentation & \textbf{High} \\
\hline
Virtualisation/Overlay & VXLAN/EVPN (DC), VRFs (5), SD-WAN overlay (pilot), NFV (future) & Overlays on overlays; policy mapping between VXLAN and SD-WAN segments complex & \textbf{High} \\
\hline
Scale \& Dynamics & 120 VLANs, 3.5k endpoints, 40 CRs/week, 250 IoT/branch & High change velocity; IoT adds MAC churn and micro-segmentation needs & \textbf{Medium} \\
\hline
Management/Tooling & Likely fragmented (Prime, DNA Center, vManage, CLI) & No single pane of glass; correlation of events across domains difficult & \textbf{High} \\
\hline
\end{tabularx}
\end{center}

\noindent\textbf{2. Overall Complexity Index}
\[
\text{Scores: } 3 + 3 + 3 + 2 + 3 = 14 \quad\Rightarrow\quad \text{Average} = \frac{14}{5} = 2.8
\]
\textbf{Interpretation:} An index of \textbf{2.8/3.0 (Very High)} indicates the network has exceeded the capacity of manual, CLI-driven operations. The CIO must prioritise \textbf{Infrastructure-as-Code (IaC)} and a \textbf{unified assurance platform} to avoid configuration drift and prolonged MTTR (Mean Time To Repair).

\noindent\textbf{3. Automation Strategies}
\begin{enumerate}
    \item \textbf{NetDevOps Pipeline with Ansible/Terraform:} Store all device configs (Nexus, ISR, vEdge) as version-controlled YAML/Jinja2 templates. CI/CD pipeline validates syntax, runs \texttt{ansible-lint} and \texttt{batfish} analysis, then pushes via API/NETCONF. Eliminates manual CLI errors on 40+ weekly changes.
    \item \textbf{Intent-Based Networking (IBN) Assurance:} Deploy Cisco DNA Center Assurance (or equivalent) to continuously verify \emph{intent} (e.g., ``All ATMs in VRF-Banking must reach Core Banking App with $<50$ ms latency''). Auto-generates remediation playbooks for drift detection.
\end{enumerate}

\noindent\textbf{Examiner's Tip:} Always link complexity metrics to \emph{business risk} (downtime, compliance fines, audit findings). GCE marks are awarded for justification, not just listing.
\end{exoresolu}

\coursec{Operational Challenges: Performance, Availability, Security, Management}

\courssub{Evaluating Operational Challenges}

\begin{methode}[Evaluating Operational Challenges: Performance, Availability, Security, Management]
\textbf{Objective:} Translate business requirements into measurable network KPIs and identify gaps.
\begin{enumerate}
    \item \textbf{Define Service-Level Objectives (SLOs) per Traffic Class:}
    \begin{itemize}
        \item \textbf{Voice/Video (EF):} One-way delay $< 150$ ms, Jitter $< 30$ ms, Loss $< 1\%$.
        \item \textbf{Transactional (AF31/AF21):} RTT $< 200$ ms, Availability $99.99\%$.
        \item \textbf{Bulk/Scavenger (BE/CS1):} Best effort, shaped to $\le 10\%$ link.
    \end{itemize}
    \item \textbf{Map Availability Architecture:} Identify Single Points of Failure (SPOF) at L1 (fibre cut), L2 (STP root), L3 (HSRP/VRRP), L4-L7 (Load Balancer, Firewall HA). Target: \textbf{Five 9s ($99.999\%$)} for core banking $\approx 5.26$ min downtime/year.
    \item \textbf{Security Posture Assessment:} Verify Zero Trust segmentation (micro-segmentation, 802.1X, MACsec), Encryption (IPsec, TLS 1.3), Visibility (SSL inspection, NetFlow), and Compliance (PCI-DSS, Cameroon Law 2010/012 on Cybersecurity).
    \item \textbf{Management Maturity Model:} Rate on 5 levels: 1-Reactive (CLI), 2-Scripted, 3-Monitored (SNMP/Syslog), 4-Automated (Ansible/NETCONF), 5-Intent-Based (AI/ML assurance). Target Level 4+.
    \item \textbf{Gap Analysis \& Remediation Plan:} Prioritise gaps by \emph{Risk $\times$ Effort} matrix.
\end{enumerate}
\textbf{Key Formula:} \textbf{Availability (\%)} $= \frac{\text{MTBF}}{\text{MTBF} + \text{MTTR}} \times 100$. Redundancy improves MTBF; Automation/Monitoring reduces MTTR.
\end{methode}

\begin{exoresolu}[Operational KPI Gap Analysis for an E-Gov Platform]
\textbf{Scenario:} \textit{The Ministry of Posts and Telecommunications (MINPOSTEL) hosts the ``e-Citizen'' portal (tax, ID, land registry) in a private cloud at the National Data Centre (Yaound\'e). Access: 500 concurrent civil servants via IPsec VPN from regional delegations; 50,000 citizens/day via public internet (HTTPS). Current metrics (30-day window):}
\begin{itemize}
    \item Portal uptime: $99.2\%$ (target $99.95\%$).
    \item Average HTTPS response time: $4.2$ s (target $< 2$ s).
    \item VPN jitter for VoIP softphones: $65$ ms (target $< 30$ ms).
    \item Firewall rule base: $1,200$ rules, last review $18$ months ago.
    \item Change success rate: $78\%$ (target $> 95\%$).
\end{itemize}

\textbf{Questions:}
\begin{enumerate}
    \item Classify each metric into Performance, Availability, Security, or Management. State the gap magnitude.
    \item Calculate the current annual downtime in hours. How many minutes of downtime does the target allow?
    \item Propose \textbf{three} concrete technical actions (one per domain: Perf, Avail, Sec) to close the largest gaps.
\end{enumerate}
\tcblower
\textbf{Detailed Solution}

\noindent\textbf{1. Classification \& Gap Magnitude}

\begin{center}
\begin{tabularx}{\linewidth}{|l|l|c|c|X|}
\hline
\rowcolor{popblueL}
\textbf{Metric} & \textbf{Domain} & \textbf{Current} & \textbf{Target} & \textbf{Gap} \\
\hline
Portal Uptime & Availability & $99.2\%$ & $99.95\%$ & $0.75\%$ deficit ($\approx 65.7$ hrs/yr excess) \\
\hline
HTTPS Response Time & Performance & $4.2$ s & $< 2$ s & $+110\%$ over target \\
\hline
VPN VoIP Jitter & Performance (Real-time) & $65$ ms & $< 30$ ms & $+117\%$ over target \\
\hline
Firewall Rule Review & Security & $18$ months & $6$ months (PCI) & $12$ months overdue \\
\hline
Change Success Rate & Management & $78\%$ & $> 95\%$ & $17$ pp deficit \\
\hline
\end{tabularx}
\end{center}

\noindent\textbf{2. Downtime Calculation}
\[
\text{Current Downtime} = (1 - 0.992) \times 365 \times 24 = \mathbf{70.08\ hours/year}
\]
\[
\text{Target Downtime} = (1 - 0.9995) \times 365 \times 24 \times 60 = \mathbf{262.8\ minutes/year} \approx \mathbf{4.38\ hours}
\]
\textbf{Gap:} Current exceeds target by $\approx 65.7$ hours/year. This is a \textbf{critical availability breach} for a government service.

\noindent\textbf{3. Remediation Actions}
\begin{enumerate}
    \item \textbf{Performance (HTTPS/VoIP):} Deploy \textbf{SD-WAN with Application-Aware Routing (AAR)} at regional delegations. Configure SLA probes (loss, latency, jitter) on both MPLS and 4G/5G links; steer VoIP (DSCP EF) to best path dynamically. Enable TCP optimisation (WAN acceleration) and TLS 1.3 + HTTP/2 on web servers to cut response time.
    \item \textbf{Availability (Portal Uptime):} Implement \textbf{Active-Active Geo-Clustering} between Yaound\'e DC and a secondary DR site (e.g., Douala). Use Global Server Load Balancing (GSLB / DNS failover) with health checks $< 10$ s. Automate failover via Ansible playbooks triggered by monitoring (Prometheus/Alertmanager).
    \item \textbf{Security (Firewall Hygiene):} Execute \textbf{Firewall Rule Recertification}: Use tooling (Tufin, FireMon, or custom script) to identify unused rules (hit count $= 0$ in 90 days), shadowed rules, and overly broad `any-any' rules. Enforce \textbf{Change Advisory Board (CAB)} workflow with automated pre-deployment simulation (Batfish) to raise change success rate.
\end{enumerate}

\noindent\textbf{Examiner's Tip:} In GCE, always convert percentages to absolute time (hours/minutes) for availability --- it demonstrates quantitative rigour. Cite specific technologies (AAR, GSLB, Batfish) not just generic terms.
\end{exoresolu}

\coursec{Design, Deployment and Supervision of an Enterprise Network}

\courssub{Hierarchical Model and Addressing Design}

\begin{methode}[Designing an Enterprise Network Architecture (Hierarchical Model)]
\textbf{Objective:} Produce a scalable, resilient, and manageable three-tier (or collapsed core) design with addressing, VLAN, and routing plan.
\begin{enumerate}
    \item \textbf{Requirements Gathering:} Sites, user counts, server farms, Wi-Fi density, WAN circuits, security zones, growth forecast (3-5 years).
    \item \textbf{Select Tier Model:}
    \begin{itemize}
        \item \textbf{Three-Tier (Core / Distribution / Access):} Large campus ($> 3$ buildings, $> 2000$ ports). Core = high-speed switching only; Distribution = policy, routing boundaries, aggregation; Access = PoE+, 802.1X.
        \item \textbf{Collapsed Core (Distribution + Core merged):} Medium campus ($< 2000$ ports). Simpler, cheaper, fewer hops.
    \end{itemize}
    \item \textbf{IP Addressing \& VLAN Plan:} Use private IPv4 (/16 or /20 per site) + IPv6 (/48 per site). VLANs per function (Mgmt, Voice, User, IoT, Guest, Server, DMZ). Map VLAN $\leftrightarrow$ Subnet $\leftrightarrow$ Security Zone (Firewall context).
    \item \textbf{Routing Design:} IGP (OSPF Area 0 at Core, stub/totally-stub areas at Distribution). Summarise routes at Distribution $\rightarrow$ Core. BGP at WAN edge / DCI. Default route only at Access.
    \item \textbf{Redundancy \& Loop Prevention:} MLAG/vPC at Access-Distribution; HSRP/VRRP for gateway; STP (Rapid PVST+) or better: \textbf{No L2 loops} via Routed Access (L3 to Access) or VXLAN EVPN.
    \item \textbf{Wireless \& Policy:} Centralised WLC (or embedded), FlexConnect for branches. 802.1X + RADIUS (NPS/ISE) for dynamic VLAN/ACL assignment.
    \item \textbf{Documentation Deliverables:} Logical Diagram (L3), Physical Diagram (Rack/Elevation), IP/ VLAN Table, Routing Table, Config Templates.
\end{enumerate}
\textbf{Key Reflex:} \emph{Routed Access} (L3 links to access switches) eliminates STP, sub-second convergence, and simplifies troubleshooting --- preferred in modern designs.
\end{methode}

\begin{exoresolu}[Campus Network Design for a University]
\textbf{Scenario:} \textit{University of Buea (UB) is building a new Faculty of Engineering campus. Requirements:}
\begin{itemize}
    \item $4$ buildings (A, B, C, D), each $4$ floors, $40$ access ports/floor ($640$ ports total).
    \item $2$ Server rooms (Building A \& C), each $20$ servers ($1$ Gbps), future $10$ Gbps.
    \item Wireless: $2$ APs/floor, high density in labs.
    \item WAN: $2 \times 1$ Gbps fibre to main campus (MPLS), $1 \times 100$ Mbps 4G backup.
    \item Security zones: Admin, Student, Lab, Server, Mgmt, Guest.
    \item Budget allows: $2$ Core switches, $8$ Distribution switches, $32$ Access switches (PoE+), $2$ WAN routers, $1$ Firewall pair (HA), $1$ WLC pair (HA).
\end{itemize}

\textbf{Questions:}
\begin{enumerate}
    \item Choose the tier model (Three-Tier or Collapsed Core) and justify.
    \item Propose an IPv4 addressing scheme (VLSM) for the $6$ security zones across the campus, using $10.10.0.0/16$. Show subnet, mask, usable hosts, and VLAN ID for each zone per building.
    \item Draw a \textbf{logical L3 diagram} (text description acceptable) showing routing adjacencies, HSRP/VRRP groups, and summarisation points.
    \item Explain why \textbf{Routed Access} is preferred over L2 Access here.
\end{enumerate}
\tcblower
\textbf{Detailed Solution}

\noindent\textbf{1. Tier Model Selection}
\textbf{Choice: Collapsed Core (Distribution + Core combined).}
\textbf{Justification:}
\begin{itemize}
    \item Port count ($640$ access + $40$ server + $64$ uplink) $< 2000$; fits in two high-density $48$-port $10$G/40G core switches (e.g., Catalyst 9500/9600).
    \item Only $4$ buildings $\rightarrow$ short fibre runs; no need for separate Core layer.
    \item Reduces latency (one less hop), cost (2 fewer chassis), and config complexity.
    \item Distribution layer handles all L3 routing, policy, aggregation; Access is pure L2 (or Routed Access).
\end{itemize}

\noindent\textbf{2. IPv4 VLSM Addressing Plan ($10.10.0.0/16$)}
\textbf{Strategy:} Allocate /20 per building ($4094$ hosts) $\rightarrow$ plenty for growth. Within each building /20, carve /24 per zone per floor or aggregate per zone.
\textbf{Simplified Allocation (per Building):}

\begin{center}
\begin{tabularx}{\linewidth}{|c|l|c|c|c|X|}
\hline
\rowcolor{popblueL}
\textbf{Zone} & \textbf{Subnet (Building A)} & \textbf{Mask} & \textbf{VLAN} & \textbf{Usable Hosts} & \textbf{Notes} \\
\hline
Mgmt (OOB) & $10.10.0.0/24$ & /24 & 10 & 254 & Switch mgmt, iDRAC, WLC \\
Admin & $10.10.1.0/24$ & /24 & 20 & 254 & Staff wired \\
Student & $10.10.2.0/23$ & /23 & 30 & 510 & High density \\
Lab & $10.10.4.0/23$ & /23 & 40 & 510 & Engineering labs \\
Server (A) & $10.10.8.0/24$ & /24 & 50 & 254 & Building A server room \\
Guest & $10.10.9.0/24$ & /24 & 60 & 254 & Captive portal \\
Wireless (AP Mgmt) & $10.10.10.0/24$ & /24 & 99 & 254 & CAPWAP tunnel \\
\hline
\end{tabularx}
\end{center}
\textit{Repeat for Building B ($10.10.16.0/20$), C ($10.10.32.0/20$), D ($10.10.48.0/20$). Server Room C uses $10.10.40.0/24$ (VLAN 50). WAN links: $10.10.255.0/30$, $10.10.255.4/30$, $10.10.255.8/30$. Loopbacks: $10.10.254.x/32$.}

\noindent\textbf{3. Logical L3 Diagram (Text Description)}
\begin{flushleft}
\texttt{
[WAN Router 1] --(OSPF Area 0)--> [Core/Dist Switch Pair (vPC/MLAG)] <--(OSPF Area 0)-- [WAN Router 2]\newline
       |                                                              |\newline
       | (BGP / Default Route)                                        | (BGP / Default Route)\newline
       v                                                              v\newline
[MPLS Cloud]                                                    [4G Backup]\newline
\newline
Core/Dist Pair (VSS/StackWise Virtual):\newline
  - L3 Interfaces (SVIs) for ALL VLANs (10,20,30,40,50,60,99) -> Gateway IP = .1\newline
  - HSRP Group per VLAN (Active/Standby) -> Virtual IP = .1\newline
  - OSPF Area 0: Advertises connected subnets + Loopbacks.\newline
  - Route Summarisation: \newline
      * To WAN: advertise 10.10.0.0/16 summary (suppress specifics).\newline
      * From WAN: receive default route only.\newline
\newline
Access Switches (per Building, per Floor):\newline
  - Option A (L2 Access): Trunk VLANs 10,20,30,40,60,99 to Core/Dist. STP Root = Core/Dist.\newline
  - Option B (ROUTED ACCESS - RECOMMENDED): \newline
      * L3 Point-to-Point links (/31 or /30) to Core/Dist.\newline
      * OSPF Area 1 (Stub) per Building.\newline
      * Core/Dist summarises Building A 10.10.0.0/20, B 10.10.16.0/20, etc. into Area 0.\newline
      * NO SVIs on Access, NO STP, NO HSRP on Access.
}
\end{flushleft}

\noindent\textbf{4. Why Routed Access?}
\begin{enumerate}
    \item \textbf{Sub-second Convergence:} OSPF L3 convergence ($\approx 50-200$ ms) vs STP ($\approx 30-50$ s for PVST+, $< 1$ s for RPVST+ but still L2 risk).
    \item \textbf{No L2 Loops / Broadcast Storms:} L3 boundary stops broadcast/multicast/unknown unicast flooding across buildings.
    \item \textbf{Equal-Cost Multi-Path (ECMP):} Access switch can load-balance across two uplinks to Core/Dist pair actively (no blocking port).
    \item \textbf{Simpler Troubleshooting:} \texttt{show ip route} on access switch shows exact path; no \texttt{show spanning-tree} ambiguity.
    \item \textbf{Security:} L3 boundary allows ACLs/firewalling at every hop; easier micro-segmentation.
\end{enumerate}

\noindent\textbf{Examiner's Tip:} Always specify \textbf{summarisation boundaries} (Area 0 $\leftrightarrow$ Stub Areas) and \textbf{HSRP/VRRP} on the \emph{Distribution/Core} layer, never on Access. Draw the diagram clearly showing L3 links for Routed Access.
\end{exoresolu}

\courssub{High Availability: Redundancy Protocols and Link Aggregation}

\begin{methode}[Configuring High Availability: Redundancy Protocols \& Link Aggregation]
\textbf{Objective:} Implement gateway redundancy and link-level redundancy without loops.
\begin{enumerate}
    \item \textbf{First-Hop Redundancy (FHRP):} Choose \textbf{HSRPv2} (Cisco) or \textbf{VRRPv3} (Multi-vendor).
    \begin{itemize}
        \item Configure \texttt{priority} (Active $= 110$, Standby $= 100$), \texttt{preempt}, \texttt{timers} (hello $100$ ms, hold $300$ ms for sub-second failover).
        \item Track uplink interfaces: \texttt{track 1 interface Gig1/0/1 line-protocol} $\rightarrow$ \texttt{standby 10 track 1 decrement 20}.
    \end{itemize}
    \item \textbf{Link Aggregation (LACP - 802.3ad):} Bundle $2+$ physical links into one logical EtherChannel.
    \begin{itemize}
        \item Mode: \texttt{active} (both sides) for negotiation.
        \item Load-balancing hash: \texttt{port-channel load-balance src-dst-ip} (L3/L4) or \texttt{src-dst-mac-ip} (default).
        \item \textbf{MLAG/vPC/StackWise Virtual:} Dual-homed Access to two Distribution switches $\rightarrow$ single logical control plane, active-active forwarding, no STP blocking.
    \end{itemize}
    \item \textbf{Spanning Tree (if L2 Access unavoidable):} \textbf{Rapid PVST+} (Cisco) or \textbf{MSTP} (Multi-vendor).
    \begin{itemize}
        \item Set Core/Dist as Root Primary/Secondary for all VLANs: \texttt{spanning-tree vlan 1-4094 root primary/secondary}.
        \item Enable \texttt{spanning-tree portfast bpduguard default} on all access ports.
        \item Enable \texttt{spanning-tree loopguard default} on uplinks.
    \end{itemize}
    \item \textbf{Verification Commands:} \texttt{show standby brief}, \texttt{show etherchannel summary}, \texttt{show spanning-tree vlan 10 detail}, \texttt{show lacp neighbor}.
\end{enumerate}
\textbf{Key Reflex:} \emph{Always track uplinks in HSRP/VRRP.} Without tracking, the Active router stays Active even if its WAN link fails $\rightarrow$ blackhole traffic.
\end{methode}

\begin{exoresolu}[HSRP + EtherChannel Configuration \& Troubleshooting]
\textbf{Scenario:} \textit{A network engineer configures HSRPv2 and LACP EtherChannel between two Distribution switches (DSW1, DSW2) and an Access switch (ASW1).}
\textbf{DSW1 Config (relevant):}
\begin{flushleft}
\texttt{
interface Port-channel10\newline
 description Link to ASW1\newline
 switchport trunk allowed vlan 10,20,30,99\newline
 spanning-tree portfast trunk\newline
!\newline
interface GigabitEthernet1/0/1\newline
 channel-group 10 mode active\newline
!\newline
interface Vlan10\newline
 ip address 10.10.10.2 255.255.255.0\newline
 standby version 2\newline
 standby 10 ip 10.10.10.1\newline
 standby 10 priority 110\newline
 standby 10 preempt\newline
 standby 10 timers 100 300\newline
 standby 10 track 1 decrement 20\newline
!\newline
track 1 interface GigabitEthernet1/0/2 line-protocol
}
\end{flushleft}
\textbf{DSW2 Config:} Identical but \texttt{priority 100}, \texttt{track 1 interface GigabitEthernet2/0/2}.
\textbf{ASW1 Config:}
\begin{flushleft}
\texttt{
interface Port-channel10\newline
 switchport trunk allowed vlan 10,20,30,99\newline
 spanning-tree portfast trunk\newline
!\newline
interface range GigabitEthernet1/0/1 - 2\newline
 channel-group 10 mode active
}
\end{flushleft}
\textbf{Observation:} After a link failure on DSW1 \texttt{Gig1/0/2} (WAN uplink), HSRP fails over correctly. But when \texttt{Gig1/0/1} (EtherChannel member) fails, HSRP \textbf{does not} fail over, and users on ASW1 lose connectivity for $30$ seconds.

\textbf{Questions:}
\begin{enumerate}
    \item Identify \textbf{two} configuration errors causing the 30-second outage.
    \item Provide the corrected configuration snippets for DSW1 and ASW1.
    \item Explain why \texttt{spanning-tree portfast trunk} on the Port-channel is dangerous if ASW1 is a single switch (not MLAG).
\end{enumerate}
\tcblower
\textbf{Detailed Solution}

\noindent\textbf{1. Configuration Errors}
\begin{enumerate}
    \item \textbf{Missing LACP \texttt{min-links} / \texttt{port-channel min-links}:} The EtherChannel (Po10) has 2 members. When one fails (\texttt{Gig1/0/1}), the channel stays \texttt{Up} (1 link left) but bandwidth halves. HSRP tracks \texttt{Gig1/0/2} (WAN), not the EtherChannel. \textbf{Critical:} HSRP tracks the \emph{WAN interface}, not the \emph{downlink to Access}. If the downlink fails partially, HSRP stays Active $\rightarrow$ traffic blackholed at DSW1 because ASW1 still sends traffic to DSW1 via remaining link, but DSW1 cannot forward to WAN? Wait, scenario says ``users on ASW1 lose connectivity''. If DSW1 WAN link is up, but DSW1-ASW1 link loses one member, traffic should flow on remaining member. \textbf{Re-read:} ``When \texttt{Gig1/0/1} (EtherChannel member) fails, HSRP does not fail over, and users lose connectivity for 30 seconds.'' $\rightarrow$ This implies \textbf{STP reconvergence} (30 sec = STP Max Age + Forward Delay). \texttt{spanning-tree portfast trunk} on a \textbf{trunk} between switches \textbf{does not} prevent STP on that trunk; it only skips Listening/Learning on \emph{access} ports. On a trunk, STP runs normally. If the EtherChannel is not negotiated correctly (see error 2), STP sees a topology change and blocks.
    \item \textbf{LACP Mode Mismatch / Non-negotiation:} DSW1 uses \texttt{channel-group 10 mode active}. ASW1 uses \texttt{channel-group 10 mode active}. This is correct. \textbf{But:} If the Port-channel interface on DSW1/DSW2 is Layer 2 (\texttt{switchport}), and ASW1 is single-homed, STP runs. The 30s outage is classic STP Topology Change Notification (TCN) + Max Age (20s) + Forward Delay (15s) $\approx$ 30-50s. \textbf{Root Cause:} \texttt{spanning-tree portfast trunk} is \textbf{invalid} on a trunk port-channel connecting two switches. It only applies to access ports. The port-channel goes through Listening/Learning/Forwarding on link failure/recovery.
\end{enumerate}

\noindent\textbf{2. Corrected Configuration}
\textbf{DSW1 / DSW2 (Best Practice: Routed Access or MLAG):} \textit{Assuming L2 Access with MLAG/vPC is not available, use Routed Access (L3). If forced L2, fix STP.}
\textbf{Option A: Routed Access (Eliminates STP entirely) - RECOMMENDED}
\begin{flushleft}
\texttt{
! DSW1 \& DSW2\newline
interface Port-channel10\newline
 no switchport\newline
 ip address 10.10.255.1/31  ! /31 for p2p (RFC 3021) or /30\newline
 ip ospf network point-to-point\newline
 ip ospf 1 area 1\newline
!\newline
interface GigabitEthernet1/0/1\newline
 no switchport\newline
 channel-group 10 mode active\newline
\newline
! ASW1\newline
interface Port-channel10\newline
 no switchport\newline
 ip address 10.10.255.0/31\newline
 ip ospf network point-to-point\newline
 ip ospf 1 area 1\newline
!\newline
interface range GigabitEthernet1/0/1 - 2\newline
 no switchport\newline
 channel-group 10 mode active
}
\end{flushleft}
\textbf{Option B: L2 Access with MLAG/vPC (Active-Active, No STP Blocking)}
\begin{flushleft}
\texttt{
! DSW1 \& DSW2 (vPC Peer-Link established)\newline
interface Port-channel10\newline
 description vPC to ASW1\newline
 switchport trunk allowed vlan 10,20,30,99\newline
 spanning-tree port type network  ! vPC requires 'network' type, not portfast\newline
 vpc 10\newline
!\newline
interface GigabitEthernet1/0/1\newline
 channel-group 10 mode active\newline
\newline
! ASW1 (Single switch, dual-homed)\newline
interface Port-channel10\newline
 switchport trunk allowed vlan 10,20,30,99\newline
 spanning-tree port type network  ! Correct type for switch-to-switch trunk
}
\end{flushleft}
\textbf{Option C: L2 Access Single-Homed (Fix STP Timers)}
\begin{flushleft}
\texttt{
! DSW1 \& DSW2 (HSRP Peers)\newline
interface Port-channel10\newline
 switchport trunk allowed vlan 10,20,30,99\newline
 spanning-tree port type network  ! Correct type for switch-to-switch trunk\newline
! REMOVE 'spanning-tree portfast trunk'\newline
!\newline
! Global: Rapid PVST+ is default. Ensure no blocking.\newline
! HSRP Tracking: Track the Port-channel LINE PROTOCOL, not physical member.\newline
track 10 interface Port-channel10 line-protocol\newline
!\newline
interface Vlan10\newline
 standby 10 track 10 decrement 20
}
\end{flushleft}

\noindent\textbf{3. Why \texttt{spanning-tree portfast trunk} is Dangerous}
\begin{itemize}
    \item \texttt{portfast} transitions a port immediately to \textbf{Forwarding} state, skipping Listening/Learning.
    \item On a \textbf{trunk} connecting two switches, this creates an \textbf{instant L2 loop} if there is any other path (e.g., another uplink, or mis-cabling).
    \item BPDU Guard (\texttt{bpduguard}) would shut the port if a BPDU is received, but if the peer also has portfast, BPDUs might not be processed fast enough.
    \item \textbf{Correct usage:} \texttt{spanning-tree port type edge} (access ports only) or \texttt{spanning-tree port type network} (switch-to-switch links, enables Bridge Assurance).
\end{itemize}

\noindent\textbf{Examiner's Tip:} Know the difference between \texttt{portfast} (legacy CLI), \texttt{port type edge} (access), \texttt{port type network} (uplinks). GCE practical questions often test STP port types and HSRP tracking objects.
\end{exoresolu}

\courssub{Monitoring, Baselining and Alerting Strategy}

\begin{methode}[Network Monitoring, Baselining \& Alerting Strategy]
\textbf{Objective:} Establish proactive visibility to detect anomalies before users complain.
\begin{enumerate}
    \item \textbf{Telemetry Collection Stack:}
    \begin{itemize}
        \item \textbf{Device Health:} SNMPv3 (polls: CPU, Memory, Interface errors, Temperature, Power) every $60$ s.
        \item \textbf{Flow Visibility:} NetFlow/IPFIX/sFlow from all L3 interfaces $\rightarrow$ Collector (nProbe, ElastiFlow, SolarWinds, PRTG). Retain $30$ days.
        \item \textbf{Event Logs:} Syslog (RFC 5424) + SNMP Traps $\rightarrow$ SIEM (Splunk, ELK, FortiAnalyzer). Severity 0-4 (Emergency-Warning) alert immediately.
        \item \textbf{Synthetic Probes:} IP SLA / TWAMP / Two-Way Active Measurement for Latency, Jitter, Loss, MOS on critical paths (Voice, VPN, Cloud).
    \end{itemize}
    \item \textbf{Baselining (Dynamic Thresholds):} Collect $4$ weeks of data. Compute \textbf{Mean ($\mu$)} and \textbf{Standard Deviation ($\sigma$)} per metric per hour-of-day/day-of-week.
    \item \textbf{Alerting Rules (Multi-Level):}
    \begin{itemize}
        \item \textbf{Critical (Page):} Device down, Link down, CPU $> 90\%$ for $5$ min, Interface errors $> 1\%$, BGP peer down, Firewall HA split-brain.
        \item \textbf{Warning (Ticket):} CPU $> 75\%$, Memory $> 85\%$, Latency $> \mu + 3\sigma$, Jitter $> 30$ ms, Flow volume anomaly ($> 3\sigma$).
        \item \textbf{Info (Log):} Config change, New MAC, New DHCP lease, Route flap.
    \end{itemize}
    \item \textbf{Correlation \& Root Cause:} Use topology-aware correlation (e.g., ``Core switch CPU high'' + ``All downstream interfaces error increment'' $\rightarrow$ Root cause = Core switch).
    \item \textbf{Reporting:} Daily/Weekly/Monthly dashboards: Availability \%, Top Talkers, Application Response Time, Capacity Trend (Interface util $> 70\%$ trigger upgrade plan).
\end{enumerate}
\textbf{Key Formula:} \textbf{Interface Utilization \%} $= \frac{\text{Octets} \times 8}{\text{Speed} \times \text{Interval}} \times 100$. Alert at $70\%$ (plan), $85\%$ (urgent), $95\%$ (critical).
\end{methode}

\begin{exoresolu}[Designing a Monitoring Policy for a Mobile Money Platform]
\textbf{Scenario:} \textit{MTN Cameroon MoMo platform runs on a clustered database (Active-Standby) across two data centres (Douala DC1, Yaound\'e DC2) linked by $2 \times 10$ Gbps DWDM. Application servers ($50$ VMs) in each DC. Users access via USSD ($*123\#$), Mobile App (HTTPS/REST), and Agent POS (TCP 8443). Current monitoring: Basic SNMP (up/down) + Ping. No flow data. Recent outage: DB replication lag $> 5$ min caused failover failure; detected $20$ min late by manual check.}

\textbf{Questions:}
\begin{enumerate}
    \item List the \textbf{minimum 5 critical metrics} (with source: SNMP/Flow/Syslog/IP SLA) that must be monitored to prevent recurrence.
    \item Define a \textbf{baselining strategy} for the DCI link utilisation. How do you distinguish ``normal backup traffic'' from ``replication lag congestion''?
    \item Write a \textbf{pseudo-alert rule} (in plain logic) for ``Database Replication Lag Critical'' using available telemetry.
\end{enumerate}
\tcblower
\textbf{Detailed Solution}

\noindent\textbf{1. Minimum 5 Critical Metrics}

\begin{center}
\begin{tabularx}{\linewidth}{|l|l|X|c|}
\hline
\rowcolor{popblueL}
\textbf{Metric} & \textbf{Source} & \textbf{Why Critical for MoMo} & \textbf{Threshold Example} \\
\hline
DCI Link Utilisation (Tx/Rx) & SNMP (IF-MIB) / NetFlow & Replication traffic saturates link $\rightarrow$ lag & $> 70\%$ (Warn), $> 85\%$ (Crit) \\
\hline
DB Replication Lag (Seconds) & Custom Script / DB Agent (Prometheus exporter) & Direct measure of RPO risk & $> 60$ s (Warn), $> 300$ s (Crit) \\
\hline
Application Response Time (USSD/REST) & IP SLA (HTTP/TCP) / APM (AppDynamics) & User-facing SLA & $> 2$ s (Warn), $> 5$ s (Crit) \\
\hline
Cluster Heartbeat / Split-Brain Status & SNMP (Vendor MIB) / Syslog & HA integrity & Any ``Split-Brain'' or ``Heartbeat Lost'' = Crit \\
\hline
Interface Errors / Discards (DCI, Server NICs) & SNMP (IF-MIB) & Physical layer degradation $\rightarrow$ silent corruption & Errors $> 0.1\%$ (Warn), $> 1\%$ (Crit) \\
\hline
\end{tabularx}
\end{center}
\textit{Additional: BGP Peer State (WAN), Firewall Session Count, Certificate Expiry (TLS).}

\noindent\textbf{2. Baselining Strategy for DCI Link}
\begin{enumerate}
    \item \textbf{Collection:} Poll \texttt{ifHCInOctets} / \texttt{ifHCOutOctets} (64-bit counters) every $60$ s via SNMPv3. Calculate $5$-min average utilisation \%.
    \item \textbf{Time-Series Database:} Store in InfluxDB/Prometheus/Graphite with tags: \texttt{interface=DCI-1}, \texttt{direction=tx}, \texttt{site=Douala}.
    \item \textbf{Dynamic Baseline:} Compute \textbf{Hourly Percentiles (P50, P95, P99)} over rolling $4$ weeks.
    \item \textbf{Anomaly Detection:} 
    \begin{itemize}
        \item \textbf{Normal Backup:} Predictable spike $02{:}00$--$04{:}00$ (P95 $\approx 60\%$). Baseline learns this pattern.
        \item \textbf{Replication Lag Congestion:} Sustained high utilisation \textbf{outside} backup window, or utilisation $> \text{P99} + 10\%$ for $> 10$ min.
        \item \textbf{Correlation:} Join with \textbf{NetFlow Top Talkers} on DCI. If Top Talker = DB Server IP (port 1521/5432/3306) $\rightarrow$ Replication traffic. If Top Talker = Backup Server IP (port 22/443/NDMP) $\rightarrow$ Backup.
    \end{itemize}
\end{enumerate}

\noindent\textbf{3. Pseudo-Alert Rule: ``DB Replication Lag Critical''}
\begin{flushleft}
\texttt{
ALERT "DB\_Replication\_Lag\_Critical"\newline
IF  (metric:db\_replication\_lag\_seconds > 300) \newline
    AND (metric:dci\_utilization\_tx > 85\% FOR 10m)\newline
    AND (netflow:top\_talker\_dci.dst\_port IN (1521, 5432, 3306) \newline
         AND netflow:top\_talker\_dci.bytes > 50\%\_of\_link)\newline
    AND (state:db\_cluster\_role == "Primary" OR "Standby")\newline
THEN\newline
    SEVERITY = CRITICAL (PAGE)\newline
    RUNBOOK = "Check DB alert log; Verify network path (ping/traceroute DCI); \newline
               Check for long-running transactions; Consider manual failover if lag > 600s."\newline
    NOTIFY = [NOC\_OnCall, DBA\_OnCall, Platform\_Lead]\newline
    AUTO\_TICKET = JIRA/ServiceNow (Priority 1)\newline
END
}
\end{flushleft}

\noindent\textbf{Examiner's Tip:} Baselining is not static thresholds. Always mention \textbf{time-of-day/week seasonality} and \textbf{correlation across data sources} (SNMP + Flow + Logs). This distinguishes a ``monitoring config'' from a ``monitoring strategy''.
\end{exoresolu}

\courssub{Structured Troubleshooting Methodology}

\begin{methode}[Structured Troubleshooting Methodology (GCE Exam Approach)]
\textbf{Objective:} Solve any network fault systematically under exam pressure.
\begin{enumerate}
    \item \textbf{Define the Problem Precisely:} ``User A cannot reach Server B via App Port X'' (not ``Internet down''). Note: Who, What, Where, When, Scope (1 user / VLAN / Site).
    \item \textbf{Gather Facts (No Assumptions):} 
    \begin{itemize}
        \item \texttt{ping} / \texttt{traceroute} (source $\rightarrow$ dest, dest $\rightarrow$ source).
        \item \texttt{show ip route <dest>} (LFIB/FIB check).
        \item \texttt{show interface <id> | include errors|drops|util}.
        \item \texttt{show logging | include <keywords>}.
        \item Check recent changes (CAB log, Git diff).
    \end{itemize}
    \item \textbf{Choose Troubleshooting Approach:}
    \begin{itemize}
        \item \textbf{Bottom-Up (OSI L1 $\rightarrow$ L7):} Physical $\rightarrow$ Data Link $\rightarrow$ Network $\rightarrow$ Transport $\rightarrow$ App. Best for ``Link Down'', ``No Adjacency''.
        \item \textbf{Top-Down (L7 $\rightarrow$ L1):} App logs $\rightarrow$ Port open? $\rightarrow$ Firewall? $\rightarrow$ Route? Best for ``Slow App'', ``Timeout''.
        \item \textbf{Divide \& Conquer (Binary Search):} Test midpoint (e.g., Firewall inside interface). If OK, problem is downstream; else upstream. Fastest for path issues.
        \item \textbf{Follow-the-Path:} Trace packet hop-by-hop (ingress $\rightarrow$ egress) checking ACL, NAT, PBR, QoS, Routing at each hop.
    \end{itemize}
    \item \textbf{Hypothesise \& Test:} One change at a time. Document command + output. Rollback if no effect.
    \item \textbf{Verify Resolution \& Close:} Confirm with user/app owner. Update documentation/Knowledge Base. Root Cause Analysis (RCA) if major.
\end{enumerate}
\textbf{Key Reflex:} \textbf{``Show, Don't Guess.''} Every claim in the exam must be backed by a \texttt{show} command output or logical deduction from provided outputs.
\end{methode}

\begin{exoresolu}[Troubleshooting Intermittent VoIP Quality]
\textbf{Scenario:} \textit{Users at the Bamenda branch report ``robotic voice'' and ``call drops'' on the Cisco IP Phones (SCCP) between 08:00--10:00 and 14:00--16:00. Branch connects via $2 \times 100$ Mbps Metro Ethernet (Primary/Backup) to HQ (Douala). HQ CallManager cluster. QoS: Voice (DSCP EF) priority queue $30\%$. Branch Router: ISR 4331. Switch: Catalyst 2960X (PoE+).}

\textbf{Provided Outputs (Simulated):}
\begin{enumerate}
    \item \texttt{show policy-map interface Gig0/0/0 (Branch WAN Primary)} shows: \texttt{Class VOICE: packets 12,450, bytes 1.2M, offered rate 850 kbps, \textbf{drop rate 12\%}}.
    \item \texttt{show interface Gig0/0/0}: \texttt{output queue 0/40, 125 drops; no buffer failures}.
    \item \texttt{show platform hardware qfp active feature qos interface Gig0/0/0}: \texttt{Queue 0 (Priority) depth 512/2048 packets}.
    \item \texttt{show ip sla statistics 10 (VoIP Jitter Probe to HQ)}: \texttt{RTT Avg 45ms, Jitter Avg 28ms, \textbf{One-way Jitter Src->Dst 42ms, Dst->Src 15ms}}.
    \item \texttt{show logging}: \texttt{\%LINK-3-UPDOWN: Interface Gig0/0/1 (Backup), changed state to Up/Down} flapping every $3$ min during peak hours.
    \item \texttt{show run | section eigrp}: \texttt{network 10.0.0.0} on both Primary and Backup. No \texttt{bandwidth} command on Backup interface.
\end{enumerate}

\textbf{Questions:}
\begin{enumerate}
    \item Using Divide \& Conquer, identify the \textbf{root cause} of the voice quality issues.
    \item Explain the mechanism linking the Backup link flap to the Primary link voice drops.
    \item Provide the \textbf{exact configuration commands} to fix the root cause on the Branch Router.
\end{enumerate}
\tcblower
\textbf{Detailed Solution}

\noindent\textbf{1. Root Cause Identification (Divide \& Conquer)}
\textbf{Path:} Phone $\rightarrow$ Access Switch $\rightarrow$ Branch Router (WAN Primary) $\rightarrow$ Metro-E $\rightarrow$ HQ Router $\rightarrow$ CallManager.
\begin{itemize}
    \item \textbf{L1/L2 (Access):} No errors on switch ports (implied). Phones register $\rightarrow$ PoE/CDP/LLDP OK.
    \item \textbf{L3 (Routing):} EIGRP adjacency stable on Primary? \textbf{Critical Clue:} Backup link (Gig0/0/1) flaps \textbf{every 3 min}. EIGRP runs on both. No \texttt{bandwidth} on Backup $\rightarrow$ default $100$ Mbps (or $1.544$ Mbps if serial? Assume Metro-E $100$M). \textbf{Metric Calculation:} EIGRP metric uses \texttt{bandwidth} + \texttt{delay}. If Backup has same BW/Delay as Primary, ECMP installs \textbf{two equal-cost paths}.
    \item \textbf{L4/L7 (QoS/Voice):} Policy-map shows \textbf{12\% drops in Priority Queue} on Primary. Queue depth $512/2048$ (25\% full) $\rightarrow$ not tail drop. \textbf{Drop cause:} \textbf{Policer/Shaper} or \textbf{Queue limit exceeded due to burst}.
    \item \textbf{The Link:} When Backup flaps \texttt{Up}, EIGRP installs route via Backup. Traffic \textbf{hashes} across both links (per-flow load balancing). Voice packets (single flow $\rightarrow$ single hash) \textbf{sometimes take Backup link}. Backup link may have \textbf{no QoS policy} (or different policy) $\rightarrow$ Voice treated as Best Effort $\rightarrow$ Jitter/Loss. When Backup flaps \texttt{Down}, route withdrawn $\rightarrow$ reconvergence $\rightarrow$ packet loss.
    \item \textbf{IP SLA Asymmetric Jitter:} Src->Dst (Branch $\rightarrow$ HQ) $42$ ms vs Dst->Src $15$ ms. Confirms \textbf{egress congestion at Branch} (Primary or Backup).
\end{itemize}
\textbf{Root Cause:} \textbf{EIGRP Equal-Cost Multi-Path (ECMP) over unequal links (Primary Metro-E + Flapping Backup) without QoS on Backup, causing voice packets to be load-balanced onto an unmanaged, unstable link.}

\noindent\textbf{2. Mechanism}
\begin{enumerate}
    \item Backup interface comes \texttt{Up} $\rightarrow$ EIGRP forms adjacency $\rightarrow$ learns same prefix via Backup.
    \item Metric equal (default bandwidth/delay) $\rightarrow$ \texttt{variance 1} (default) $\rightarrow$ ECMP installed in RIB/FIB.
    \item CEF Per-Flow Load Balancing: Voice flow (SRC IP:Port, DST IP:Port, Protocol) hashes to \textbf{one} link.
    \item If hash selects Backup: Voice packets traverse Backup interface $\rightarrow$ \textbf{No Priority Queue configured} $\rightarrow$ Queued in Best Effort $\rightarrow$ High Jitter/Loss (robotic voice).
    \item Backup flaps \texttt{Down} $\rightarrow$ Route removed $\rightarrow$ Traffic shifts to Primary $\rightarrow$ Brief reconvergence loss (call drop).
    \item Cycle repeats every $3$ min $\rightarrow$ Intermittent quality exactly during flap windows (peak hours = more traffic = more noticeable).
\end{enumerate}

\noindent\textbf{3. Fix Configuration (Branch Router)}
\textbf{Goal:} Prevent Backup from being used for Voice (or any traffic) unless Primary fails. Make Backup a \textbf{True Standby} (Floating Static or EIGRP with higher metric).
\textbf{Solution: Floating Static Route + IP SLA Tracking (Most Robust) OR EIGRP Metric Tuning.}
\textbf{Here: EIGRP Metric Tuning (Simpler, fits scenario).}

\begin{flushleft}
\texttt{
! 1. Set accurate bandwidth on Backup (if 100M Metro-E)\newline
interface GigabitEthernet0/0/1\newline
 description Backup\_WAN\_MTN\newline
 bandwidth 100000      ! kbps (100 Mbps)\newline
 delay 100             ! tens of microseconds (10 ms) - HIGHER than Primary (e.g., Primary delay 10)\newline
! Primary Gig0/0/0 should have: bandwidth 100000, delay 10 (1 ms)\newline
\newline
! 2. Ensure EIGRP prefers Primary via Delay (K1=1, K3=1 default)\newline
! Primary Delay = 10 (1ms), Backup Delay = 100 (10ms) -> Metric Backup ~ 10x Primary\newline
! No ECMP. Backup only used if Primary fails.\newline
\newline
! 3. (BEST PRACTICE) Apply SAME QoS Policy to Backup Interface\newline
interface GigabitEthernet0/0/1\newline
 service-policy output WAN\_QOS\_POLICY   ! Same as Primary\newline
\newline
! 4. (ALTERNATIVE) Floating Static Route with IP SLA Tracking (If not running EIGRP on Backup)\newline
ip sla 100\newline
 icmp-echo 10.10.10.1 source-interface GigabitEthernet0/0/0  ! Primary GW\newline
 frequency 5\newline
 threshold 200\newline
 timeout 500\newline
ip sla schedule 100 life forever start-time now\newline
track 100 ip sla 100 reachability\newline
!\newline
ip route 0.0.0.0 0.0.0.0 GigabitEthernet0/0/1 10.10.20.1 254 track 100  ! AD 254 > EIGRP 90\newline
! Primary Default via EIGRP (AD 90). Backup Static (AD 254) only if Track Down.
}
\end{flushleft}

\noindent\textbf{Verification:}
\begin{flushleft}
\texttt{
show ip route | include 0.0.0.0  ! Only one default via Primary\newline
show ip eigrp topology          ! Backup route FD > Primary FD (Feasible Distance)\newline
show policy-map int Gig0/0/1    ! Voice drops = 0\%\newline
show ip sla stat 100            ! Track UP
}
\end{flushleft}

\noindent\textbf{Examiner's Tip:} In troubleshooting questions, \textbf{correlate timestamps} (flap every 3 min $\leftrightarrow$ user complaints 08-10, 14-16). Always check \texttt{show ip route} for ECMP and \texttt{show policy-map interface} for drops. The fix must address \textbf{both} routing stability \textbf{and} QoS consistency.
\end{exoresolu}

\coursec{Integration Activities: Case Study, Simulation and Reporting}

\courssub{Integration Project Lifecycle}

\begin{methode}[Integration Project Lifecycle: Case Study, Simulation \& Reporting]
\textbf{Objective:} Execute a complete network integration project from requirements to handover documentation.
\begin{enumerate}
    \item \textbf{Phase 1: Requirements \& Constraints (RACI Matrix).}
    \begin{itemize}
        \item Stakeholders: Sponsor, NetOps, SecOps, Server/App Teams, Vendor PM.
        \item Constraints: Budget, Maintenance Windows (Cameroon: often weekends 22h--06h), Regulatory (ANTIC, Telecommunications Law), Existing Contracts.
        \item Deliverable: \textbf{Project Charter} + \textbf{High-Level Design (HLD)}.
    \end{itemize}
    \item \textbf{Phase 2: Detailed Design \& Simulation (LLD).}
    \begin{itemize}
        \item Tools: Cisco Packet Tracer / GNS3 / EVE-NG / CML / NSO.
        \item Build \textbf{Digital Twin}: Exact device images (IOS-XE 17.x), topology, config templates.
        \item Test Scenarios: Normal ops, Single link failure, Device failure, Failover convergence time, Traffic path verification (traceroute), Security policy enforcement (pen test lite).
        \item Deliverable: \textbf{Low-Level Design (LLD)} + \textbf{Test Plan (ATP - Acceptance Test Plan)} + \textbf{Simulation Report (Screenshots/Logs)}.
    \end{itemize}
    \item \textbf{Phase 3: Staged Deployment.}
    \begin{itemize}
        \item \textbf{Lab Staging:} Pre-config devices, burn-in (48h), upgrade firmware to approved baseline.
        \item \textbf{Pilot:} One non-critical site / VLAN. Monitor $48$h.
        \item \textbf{Phased Rollout:} Site-by-site or Module-by-module (Core $\rightarrow$ Dist $\rightarrow$ Access). Rollback plan per step ($< 30$ min).
        \item \textbf{Change Management:} RFC per site, CAB approval, Post-Implementation Review (PIR).
    \end{itemize}
    \item \textbf{Phase 4: Validation \& Handover.}
    \begin{itemize}
        \item Execute ATP: Ping, iPerf (BW), IP SLA (Jitter), Security Scan (Nessus/OpenVAS), Config Audit (Nipper/Netmiko).
        \item \textbf{As-Built Documentation:} Physical/Logical Diagrams (Visio/Draw.io), IPAM Export (NetBox/InfoBlox), Config Backups (Git), Runbooks (``How to replace Access Switch''), Contact List.
        \item \textbf{Knowledge Transfer:} Workshop with NetOps team.
    \end{itemize}
\end{enumerate}
\textbf{Key Deliverable:} \textbf{Project Close-Out Report} containing: Executive Summary, Scope vs Delivered, Test Results (Pass/Fail), Open Issues (Risk Register), Lessons Learned, Sign-offs.
\end{methode}

\begin{exoresolu}[Integration Case Study: Regional Hospital Network Upgrade]
\textbf{Scenario:} \textit{The Ministry of Public Health (MINSANT\'E) mandates upgrading the LAN/WAN of 5 Regional Hospitals (Douala, Yaound\'e, Bafoussam, Garoua, Maroua) to support Telemedicine (HD Video Conferencing), PACS (Medical Imaging), and HMIS (Hospital Management Information System). Current: Flat L2, 100 Mbps, no redundancy, no QoS, unmanaged switches. Budget: 150M FCFA. Timeline: 3 months.}

\textbf{Questions:}
\begin{enumerate}
    \item Draft a \textbf{High-Level Design (HLD)} summary (max 150 words) covering Architecture, WAN, Security, and Management.
    \item Define \textbf{3 Acceptance Test Cases (ATP)} for the Telemedicine VLAN (VLAN 100) with Pass/Fail criteria.
    \item Outline the \textbf{Staged Rollout Plan} for one hospital (Douala Laquintinie) minimising disruption to 24/7 Emergency Dept.
    \item List \textbf{4 critical items} for the Final Handover Documentation package.
\end{enumerate}
\tcblower
\textbf{Detailed Solution}

\noindent\textbf{1. High-Level Design (HLD) Summary}
\textbf{Architecture:} Collapsed Core (Core/Distribution) per hospital using 2x Catalyst 9500 (StackWise Virtual) for $10$G uplinks/servers, 9300 Access (PoE+) for $1$G/2.5G to endpoints. \textbf{WAN:} SD-WAN (Cisco Meraki / Viptela) over Dual ISP (CAMTEL Fibre 100M + MTN 4G/5G 50M) with Application-Aware Routing steering Telemedicine (DSCP EF/AF41) to best path. \textbf{Security:} Zone-Based Firewall (HQ FortiGate HA) + Local Segmentation (TrustSec/SGT) for Medical (PACS/HMIS), Admin, Guest, IoT (Biomed). \textbf{Management:} Cloud Controller (Meraki/DNA Center) for Zero-Touch Provisioning (ZTP), Assurance, Config Compliance. \textbf{IPv4/IPv6 Dual Stack} throughout. \textbf{Resilience:} Active-Active Core, Dual-homed Access (MLAG), $99.99\%$ target.

\noindent\textbf{2. Acceptance Test Cases (Telemedicine VLAN 100)}

\begin{center}
\begin{tabularx}{\linewidth}{|c|X|X|X|}
\hline
\rowcolor{popblueL}
\textbf{ID} & \textbf{Test Case} & \textbf{Procedure} & \textbf{Pass Criteria} \\
\hline
ATP-TM-01 & \textbf{End-to-End HD Video Quality} & Initiate 1080p30 Zoom/Teams call between Doctor (Hosp) $\leftrightarrow$ Specialist (HQ). Run $30$ min. Capture IP SLA / NetFlow. & MOS $\ge 4.0$; One-way Delay $< 150$ ms; Jitter $< 30$ ms; Loss $< 0.5\%$; Bandwidth $\ge 4$ Mbps sustained. \\
\hline
ATP-TM-02 & \textbf{WAN Failover \& Session Persistence} & Simulate Primary ISP failure (shutdown WAN1). Monitor active Telemedicine call. & Call \textbf{does not drop}; Failover time $< 2$ s (SD-WAN sub-second); Jitter spike $< 50$ ms during switch; Auto-failback when Primary restored. \\
\hline
ATP-TM-03 & \textbf{Segmentation \& Security Enforcement} & From Telemedicine VLAN (10.20.100.x), attempt: (a) SSH to PACS Server (10.20.50.x), (b) Internet (google.com), (c) Admin VLAN (10.20.10.x). & (a) \textbf{Permit} (TCP 104/443 DICOM); (b) \textbf{Permit} (Proxy/NAT); (c) \textbf{Deny} (Log generated). No lateral movement. \\
\hline
\end{tabularx}
\end{center}

\noindent\textbf{3. Staged Rollout Plan: Douala Laquintinie Hospital}
\begin{enumerate}
    \item \textbf{Week 1 (Prep):} Rack/Stack new Core (9500 SVL) in Server Room A (existing). Rack Access switches in IDFs (per floor). Cable Mgmt, Labeling (TIA-606). Firmware upgrade to Gold Image. Configure Mgmt VLAN, OOB (Console Server), SNMPv3, Syslog, NTP.
    \item \textbf{Week 2 (Core \& WAN):} Configure Core SVL, Routing (OSPF/BGP), SD-WAN Edge onboarding, Firewall HA, DHCP/DNS. \textbf{Test in Lab/Isolated VLAN}. \textbf{Maintenance Window (Sat 22h--Sun 06h):} Cutover WAN links (CAMTEL/MTN) to New Core. Validate Internet/MPLS. \textbf{Rollback:} Re-patch old router ($< 15$ min).
    \item \textbf{Week 3 (Pilot - Non-Critical):} Migrate \textbf{Admin Building} (Offices, OPD Admin) only. New Access switches $\rightarrow$ Core. Enable 802.1X, QoS, Telemedicine VLAN (test with dummy endpoint). Monitor $48$h.
    \item \textbf{Week 4 (Phased Clinical - Night Windows Only):} 
    \begin{itemize}
        \item Night 1: Floor 1 (Wards) $\rightarrow$ Migrate Access uplinks to New Core. Validate Nurse stations, VoIP phones.
        \item Night 2: Floor 2 (ICU/OT) $\rightarrow$ \textbf{Critical.} Pre-stage configs. Dual-homed APs. Validate Biomed devices (IoT VLAN).
        \item Night 3: Floor 3 (Imaging/PACS) $\rightarrow$ High BW test (iPerf 10G). Validate DICOM flow.
        \item Night 4: Emergency Dept (ED) $\rightarrow$ \textbf{Last.} Run parallel old/new for $1$ hour. Cutover during lowest census (03h--05h). Immediate validation: Patient monitor network, VoIP, HMIS.
    \end{itemize}
    \item \textbf{Week 5:} Decommission old switches. Final ATP execution. As-Built docs. Handover.
\end{enumerate}

\noindent\textbf{4. Final Handover Documentation (4 Critical Items)}
\begin{enumerate}
    \item \textbf{As-Built Diagrams (L1/L2/L3):} Physical rack elevations (port mapping), Logical topology (VLAN, IP, Routing, Security Zones), WAN circuits (Circuit IDs, Provider contacts).
    \item \textbf{Configuration Repository:} Git repo with \textbf{version-controlled} configs for every device (Core, Access, FW, WLC, SD-WAN), tagged \texttt{v1.0-handover}. Includes \texttt{show tech-support} outputs.
    \item \textbf{Operational Runbooks (Playbooks):} 
    \begin{itemize}
        \item ``Replace Failed Access Switch (ZTP Procedure)''.
        \item ``SD-WAN Failover Verification / Manual Failback''.
        \item ``Add New Telemedicine Endpoint (VLAN, ACL, QoS)''.
        \item ``Disaster Recovery: Core Switch Failure Recovery''.
    \end{itemize}
    \item \textbf{Acceptance Test Report \& Sign-Off:} Completed ATP spreadsheet (Pass/Fail/Evidence), Open Issues Log (with Owner/Target Date), Capacity Baseline Report (Interface Util, CPU/Mem), \textbf{Stakeholder Signatures} (Hospital IT Dir, MINSANT\'E PM, Vendor PM).
\end{enumerate}

\noindent\textbf{Examiner's Tip:} Integration questions test \textbf{project management} as much as technical skill. Always mention: \textbf{Maintenance Windows}, \textbf{Rollback Plans}, \textbf{Parallel Run} for critical areas, and \textbf{As-Built vs Design} documentation. Cameroon context: Factor in \textbf{rainy season fibre cuts} (redundancy), \textbf{power instability} (UPS/Generator checks), and \textbf{ANTIC compliance} (logging, encryption).
\end{exoresolu}

\courssub{Network Simulation and Modelling}

\begin{methode}[Network Simulation \& Modelling with Packet Tracer / GNS3]
\textbf{Objective:} Validate designs and configurations in a virtual environment before deployment.
\begin{enumerate}
    \item \textbf{Topology Construction:} Replicate exact physical/logical topology. Use \textbf{real device images} (IOSv, IOS-XEv, CSR1000v, NX-OSv 9000, vWLC, vFMC) matching production versions.
    \item \textbf{Configuration Injection:} 
    \begin{itemize}
        \item Packet Tracer: Paste config via CLI tab or \texttt{config.text} import.
        \item GNS3/EVE-NG: Use \textbf{Cloud/NAT} nodes for Internet access; \texttt{config replace} or Netmiko/Ansible push to VMs.
    \end{itemize}
    \item \textbf{Traffic Generation:}
    \begin{itemize}
        \item \textbf{Built-in:} Packet Tracer Simulation Mode (step-by-step packet capture), PDU inspection.
        \item \textbf{External:} \texttt{iPerf3} (VMs), \texttt{TRex} / \texttt{Ostinato} (L2-L4), \texttt{Scapy} (custom packets), \texttt{IP SLA} responders on routers.
    \end{itemize}
    \item \textbf{Test Scenarios (Automated via Python/Expect):}
    \begin{enumerate}
        \item \textbf{Baseline:} Ping, Traceroute, iPerf (TCP/UDP), BGP/OSPF convergence.
        \item \textbf{Failure Injection:} \texttt{shutdown} interface, \texttt{poweroff} VM, \texttt{tc qdisc} (latency/loss/jitter emulation), Cable pull (virtual).
        \item \textbf{Security:} Nmap scan, Metasploit aux, Custom exploit (safe lab).
    \end{enumerate}
    \item \textbf{Data Collection \& Analysis:} 
    \begin{itemize}
        \item Packet Captures (PCAP) $\rightarrow$ Wireshark (IO Graphs, Expert Info, VoIP Analysis).
        \item Telemetry: SNMP/Streaming Telemetry (gNMI) $\rightarrow$ Prometheus/Grafana dashboards.
        \item Convergence Timing: \texttt{show ip route} timestamps, EIGRP/OSPF event logs, BFD flaps.
    \end{itemize}
    \item \textbf{Report:} Topology screenshot, Config snippets, Test Table (Scenario, Expected, Actual, Pass/Fail, PCAP link), Recommendations.
\end{enumerate}
\textbf{Key Reflex:} Simulation is \textbf{not} a substitute for hardware ASIC behaviour (QoS queues, TCAM limits, CPU punt). Always note \textbf{simulation limitations} in report.
\end{methode}

\begin{exoresolu}[Simulating QoS \& Failover for a Branch Office]
\textbf{Scenario:} \textit{You must validate the QoS policy for a branch router (ISR 4000) before deploying to 50 sites. The policy prioritises Voice (EF), Video (AF41), Critical Data (AF31), and scavenges Bulk (CS1). WAN: $50$ Mbps Metro-E. You use GNS3 with a CSR1000v (Router) + 2x Ubuntu VMs (Client/Server).}

\textbf{Questions:}
\begin{enumerate}
    \item Describe how to \textbf{emulate the 50 Mbps WAN bottleneck} on the CSR1000v GigabitEthernet interface in GNS3.
    \item Write the \textbf{MQC (Modular QoS CLI) configuration} for the 4-class policy (Voice 30\%, Video 20\%, Critical 20\%, Best Effort 30\%).
    \item Design a \textbf{traffic generation script} (using \texttt{iPerf3} or \texttt{Scapy}) to simultaneously load all 4 classes and verify prioritisation under congestion.
    \item What \textbf{Wireshark filters/columns} would you use to prove Voice packets experienced near-zero queuing delay while Bulk packets were delayed/dropped?
\end{enumerate}
\tcblower
\textbf{Detailed Solution}

\noindent\textbf{1. Emulating 50 Mbps Bottleneck on CSR1000v}
CSR1000v interfaces are virtual (virtio) and run at line rate (10G+). Two methods:
\begin{enumerate}
    \item \textbf{Policy-Map Shaping (Preferred, tests QoS logic):} Apply a \textbf{Parent Shaper} to the physical interface, then \textbf{Child Policy} for queuing.
    \begin{flushleft}
    \texttt{
    policy-map WAN\_PARENT\_SHAPER\newline
     class class-default\newline
      shape average 50000000  ! 50 Mbps\newline
      service-policy WAN\_CHILD\_QOS\newline
    !\newline
    interface GigabitEthernet2  ! WAN Interface\newline
     service-policy output WAN\_PARENT\_SHAPER
    }
    \end{flushleft}
    \item \textbf{GNS3 Link Emulation (Tests L1/L2 behaviour):} Edit the Link between CSR and ``WAN Cloud''/Next Hop. In GNS3 Link properties: \textbf{Bandwidth Limit = 50000 Kbps}, \textbf{Latency = 10 ms}, \textbf{Jitter = 2 ms}, \textbf{Loss = 0.1\%}. \textit{Note: This shapes at the virtual wire level, bypassing router QoS queues. Use Method 1 for QoS validation.}
\end{enumerate}

\noindent\textbf{2. MQC 4-Class Configuration}
\begin{flushleft}
\texttt{
! 1. Class Maps (Match DSCP)\newline
class-map match-any VOICE\newline
 match dscp ef\newline
class-map match-any VIDEO\newline
 match dscp af41 cs4\newline
class-map match-any CRITICAL\_DATA\newline
 match dscp af31 af32 af33\newline
class-map match-any SCAVENGER\newline
 match dscp cs1\newline
\newline
! 2. Child Policy (Queuing - LLQ + CBWFQ)\newline
policy-map WAN\_CHILD\_QOS\newline
 class VOICE\newline
  priority level 1\newline
  police cir 15000000 bc 187500  ! 30\% of 50M = 15M, burst ~12ms\newline
  conform-action transmit\newline
  exceed-action drop\newline
 class VIDEO\newline
  priority level 2\newline
  police cir 10000000 bc 250000  ! 20\% = 10M\newline
  conform-action transmit\newline
  exceed-action drop\newline
 class CRITICAL\_DATA\newline
  bandwidth remaining percent 20\newline
  random-detect dscp-based\newline
 class class-default\newline
  bandwidth remaining percent 30\newline
  random-detect dscp-based\newline
  ! Scavenger (CS1) falls here; give minimal or use 'shape' to limit\newline
  ! Better: Explicit class SCAVENGER with 'bandwidth remaining percent 5' \newline
  ! and class-default 25\% for true Best Effort.\newline
\newline
! 3. Parent Shaper (Applied to Interface)\newline
policy-map WAN\_PARENT\_SHAPER\newline
 class class-default\newline
  shape average 50000000\newline
  service-policy WAN\_CHILD\_QOS
}
\end{flushleft}

\noindent\textbf{3. Traffic Generation Script (Python + iPerf3 / Scapy)}
\textbf{Concept:} Run 4 parallel streams from Client VM to Server VM, each marking DSCP via \texttt{IP\_TOS} socket option or Scapy.
\begin{flushleft}
\texttt{
\# traffic\_gen.py (Run on Client VM)\newline
import subprocess, threading, time\newline
from scapy.all import IP, UDP, Raw, send, Ether\newline
\newline
SERVER\_IP = "10.10.20.10"\newline
DURATION = 60 \#\# seconds\newline
\newline
def iperf\_stream(dscp\_val, port, bw\_mbps, label):\newline
    \# iPerf3 supports -S (TOS) for DSCP (TOS = DSCP << 2)\newline
    tos = dscp\_val << 2\newline
    cmd = f"iperf3 -c \{SERVER\_IP\} -p \{port\} -u -b \{bw\_mbps\}M -t \{DURATION\} -S \{tos\} --get-server-output"\newline
    print(f"[\{label\}] Starting: \{cmd\}")\newline
    subprocess.run(cmd, shell=True)\newline
\newline
\# Scapy for precise control (e.g., constant bit rate Voice)\newline
def scapy\_voice\_stream():\newline
    \# G.711: 160 bytes payload + 40 bytes IP/UDP/RTP = 200 bytes @ 50 pps = 80 kbps\newline
    \# DSCP EF = 46 -> TOS 0xB8\newline
    pkt = Ether()/IP(dst=SERVER\_IP, tos=0xB8)/UDP(sport=50000, dport=50000)/Raw(b'V'*160)\newline
    send(pkt, inter=0.02, loop=1, verbose=0)  \# 50 pps\newline
\newline
\# Launch Threads\newline
threads = [\newline
    threading.Thread(target=iperf\_stream, args=(46, 5201, 20, "VOICE-EF")),      \# 20M offered > 15M policed\newline
    threading.Thread(target=iperf\_stream, args=(34, 5202, 15, "VIDEO-AF41")),    \# 15M offered > 10M policed\newline
    threading.Thread(target=iperf\_stream, args=(26, 5203, 15, "CRITICAL-AF31")), \# 15M offered\newline
    threading.Thread(target=iperf\_stream, args=(8,  5204, 30, "BULK-CS1")),      \# 30M offered\newline
]\newline
for t in threads: t.start()\newline
for t in threads: t.join()\newline
print("Test Complete.")\newline
}
\end{flushleft}
\textit{Server side:} \texttt{iperf3 -s -p 5201 -p 5202 -p 5203 -p 5204} (run multiple instances or use single port with `-P` parallel streams differentiated by DSCP not possible easily; use multiple ports).

\noindent\textbf{4. Wireshark Analysis}
\begin{itemize}
    \item \textbf{Capture Points:} SPAN/Mirror on CSR \texttt{Gig2} (WAN) egress (TX) $\rightarrow$ Best view of queuing.
    \item \textbf{Columns to Add:} \texttt{Frame Time Delta Displayed} ($\Delta t$), \texttt{IP DSCP Value}, \texttt{IP Packet Length}, \texttt{Queue Index} (if available via metadata, else infer).
    \item \textbf{Filters:}
    \begin{itemize}
        \item \texttt{ip.dscp == 46} (Voice) $\rightarrow$ Plot \texttt{IO Graph}: Y-Axis = \texttt{Frame Time Delta} (inter-packet gap). Expect \textbf{flat line $\approx 20$ ms} (50 pps). Jitter = variation in $\Delta t$.
        \item \texttt{ip.dscp == 8} (Bulk) $\rightarrow$ IO Graph: Expect \textbf{bursts, gaps, retransmissions} (TCP) or large $\Delta t$ (UDP).
        \item \texttt{tcp.analysis.retransmission} $\rightarrow$ Count for Bulk/Critical.
        \item \texttt{frame.len > 1400 \&\& ip.dscp == 46} $\rightarrow$ Verify no fragmentation.
    \end{itemize}
    \item \textbf{Key Metric:} \textbf{Queuing Delay} = (Timestamp at Router Egress) - (Timestamp at Router Ingress). Requires 2 capture points or router timestamping (PTP). In sim: Compare \texttt{IP SLA} responder timestamps.
\end{itemize}

\noindent\textbf{Examiner's Tip:} In GCE practicals, you may be asked to \textbf{interpret a provided PCAP} or \texttt{show policy-map interface} output. Know how to calculate \textbf{Drop Rate} = \texttt{packets dropped} / \texttt{packets offered}. Know that \texttt{priority} without \texttt{police} is dangerous (starves other traffic); \texttt{priority} \textbf{implies} policing in modern IOS-XE.
\end{exoresolu}

\courssub{Professional Technical Report Writing}

\begin{methode}[Professional Technical Report Writing for GCE \& Industry]
\textbf{Objective:} Structure findings, designs, or incident reports to meet professional and examination standards.
\begin{enumerate}
    \item \textbf{Title Page:} Project Title, Student Name/Index, Centre, Date, Supervisor.
    \item \textbf{Executive Summary (Abstract):} $< 250$ words. Context, Problem, Methodology, Key Results, Conclusion, Recommendation. \textbf{No jargon unexplained.}
    \item \textbf{Table of Contents, Figures, Tables, Acronyms.}
    \item \textbf{1. Introduction:} Background, Scope, Objectives, Constraints (Budget, Time, Law), Assumptions.
    \item \textbf{2. Literature Review / Technology Overview:} Relevant standards (IEEE 802.1Q, RFC 791, MEF 10.3), Vendor docs, Previous work.
    \item \textbf{3. Methodology / Design Approach:} Top-down (PPDIOO) or Agile. Tools used (Packet Tracer, GNS3, Visio, NetBox, Ansible).
    \item \textbf{4. Detailed Design / Implementation:}
    \begin{itemize}
        \item Logical/Physical Diagrams (Numbered, Captioned, Referenced in text).
        \item Addressing Tables (VLAN, Subnet, VRF, Interface, Loopback).
        \item Configuration Snippets (Key sections only: Routing, QoS, Security, HA). \textbf{Annotate} critical lines.
        \item Scripts/Automation Code (Appendix full, snippet in body).
    \end{itemize}
    \item \textbf{5. Testing, Validation \& Results:} Test Plan Table (ID, Description, Steps, Expected, Actual, Pass/Fail, Evidence Ref). Graphs (Throughput, Convergence, CPU). PCAP Analysis screenshots.
    \item \textbf{6. Discussion / Analysis:} Interpret results. Compare with targets. Explain anomalies. Link to theory (e.g., ``TCP Global Synchronisation observed due to Tail Drop; RED configured mitigatedit'').
    \item \textbf{7. Conclusion \& Recommendations:} Summary of achievement. Future work (IPv6, Segmentation, AI-Ops). Lessons Learned.
    \item \textbf{References:} IEEE, RFC, Vendor Config Guides, Cameroon Law (e.g., Law 2010/012), MINESEC Syllabus. Harvard/IEEE citation style.
    \item \textbf{Appendices:} Full Configs, Raw Data (CSV), Scripts, Glossary.
\end{enumerate}
\textbf{Golden Rules:} 
\begin{itemize}
    \item \textbf{Third Person Passive:} ``The router was configured...'' not ``I configured...''.
    \item \textbf{Figures/Tables:} Must be referenced \textbf{before} they appear (``Figure 3 shows...''). Captions \textbf{below} figures, \textbf{above} tables.
    \item \textbf{Units:} SI units (Gbps, ms, \%). No ``Megs'' or ``K''.
    \item \textbf{Version Control:} Report filename: \texttt{ProjectName\_v1.2\_Date\_Author.pdf}.
\end{itemize}
\end{methode}

\begin{exoresolu}[Critiquing a Student Project Report]
\textbf{Scenario:} \textit{An Upper Sixth student submits a project report titled ``Design of a Secure Network for a Small Business''. The examiner notes the following issues in the draft:}
\begin{enumerate}
    \item The ``Executive Summary'' is $2$ pages long and contains configuration commands.
    \item The ``Methodology'' section says ``We used Packet Tracer because it is easy''.
    \item A topology diagram (Figure 4) has no caption, no legend, uses coloured lines only (no labels), and is referenced as ``the diagram below'' on the previous page.
    \item The ``Testing'' section shows screenshots of \texttt{ping} working, but no \texttt{traceroute}, no \texttt{show ip route}, no failure testing.
    \item The ``Conclusion'' introduces a new idea: ``We should add a Wireless Controller''.
    \item References include ``Google.com'' and ``Cisco Website''.
\end{enumerate}

\textbf{Questions:}
\begin{enumerate}
    \item For each issue (1-6), state the \textbf{specific reporting standard violated} (refer to the Method above) and the \textbf{required correction}.
    \item Rewrite a \textbf{correct Executive Summary} (max 150 words) for a project deploying SD-WAN for a 10-branch enterprise.
\end{enumerate}
\tcblower
\textbf{Detailed Solution}

\noindent\textbf{1. Critique \& Corrections}

\begin{center}
\begin{tabularx}{\linewidth}{|c|X|X|}
\hline
\rowcolor{popblueL}
\textbf{Issue} & \textbf{Standard Violated} & \textbf{Required Correction} \\
\hline
1. Exec Summary too long + contains config & \textbf{Length ($<250$ words)}, \textbf{Content (no jargon, no config)} & Condense to 1 paragraph ($< 150$ words): Context, Problem, Solution, Key Result, Recommendation. Move configs to Section 4/Appendix. \\
\hline
2. Methodology vague (``easy'') & \textbf{Methodology rigour:} Must justify tool choice against requirements (scale, fidelity, cost). & Rewrite: ``Packet Tracer v8.2 selected for Layer 2/3 protocol support (OSPF, EIGRP, QoS, AAA), multi-device simulation capacity ($> 20$ nodes), and zero cost. GNS3 considered but rejected due to host resource limits and lack of native switching ASIC simulation.'' \\
\hline
3. Figure 4: No caption, no legend, colour-only, vague reference & \textbf{Figure Standards:} Numbered, Caption \textbf{below}, Legend mandatory, Colour-blind safe (patterns + labels), Referenced \textbf{by number} before appearance. & Add: ``\textbf{Figure 4:} Logical L3 Topology showing OSPF Areas (Area 0 Core, Area 1--10 Branches), VRF Lite separation (Red/Blue), and SD-WAN Overlay tunnels (dashed).'' Add legend: Solid=Physical, Dashed=Overlay, Red=Voice VLAN, Blue=Data VLAN. Text: ``Figure 4 illustrates the logical topology...'' \\
\hline
4. Testing only ping success & \textbf{Validation Depth:} Must test Negative cases, Failure scenarios, Performance metrics, Security enforcement. & Add Test Cases: ATP-02: Core Link Failure $\rightarrow$ Convergence $< 200$ ms (capture). ATP-03: Unauthorised VLAN hop $\rightarrow$ Drop + Log. ATP-04: iPerf TCP 10G $\rightarrow$ Throughput $> 9.5$ Gbps. Evidence: PCAP files, CLI logs. \\
\hline
5. Conclusion introduces new scope & \textbf{Conclusion Rule:} Summarise \textbf{only} what was done. New ideas $\rightarrow$ ``Future Work'' subsection. & Move Wireless Controller to ``7.2 Future Work: Integration of Catalyst 9800 WLC for 802.11ax upgrade and centralised RF management.'' \\
\hline
6. References non-authoritative & \textbf{Citation Authority:} Primary sources (RFC, IEEE Std, Vendor Config Guides, Law). No search engines, no Wikipedia. & Replace with: \texttt{[1] RFC 791: Internet Protocol. [2] IEEE 802.1Q-2018: Bridges and Bridged Networks. [3] Cisco ISR 4000 Series Config Guide, IOS-XE 17.6. [4] Law No. 2010/012 of 21 Dec 2010 on Cybersecurity in Cameroon. [5] MEF 10.3: Ethernet Services Attributes.} \\
\hline
\end{tabularx}
\end{center}

\noindent\textbf{2. Correct Executive Summary (SD-WAN Deployment)}

\begin{quote}
\textbf{Executive Summary}

This project delivers the design, simulation, and staged deployment plan for an SD-WAN overlay connecting 10 branch offices to a dual-data-centre core for \textit{Enterprise Cameroon SA}. The existing MPLS-only WAN incurred high OpEx (15M FCFA/month) and lacked application visibility, causing 40\% degradation in ERP response time during peak hours.

A Cisco Viptela (vManage/vSmart/vBond) fabric was designed over hybrid transports (MPLS + 4G/5G LTE), implementing Application-Aware Routing (AAR) with SLA-based steering for critical traffic (Voice EF, ERP AF31). Security integrates Zone-Based Firewalling at hubs and SIG (Umbrella) for direct internet access.

Validation in Cisco Modeling Labs (CML) confirmed sub-second failover ($< 300$ ms), 60\% WAN cost reduction (projected 6M FCFA/month), and ERP latency improvement from 420 ms to 95 ms (P95). The phased rollout plan (Pilot $\rightarrow$ 3 waves) ensures zero-downtime migration with automated rollback.

\textbf{Recommendation:} Approve Phase 1 (Pilot at Douala/Bafoussam branches) for Q3 2025. Future work includes Zero-Trust Network Access (ZTNA) integration for remote workforce.
\end{quote}

\noindent\textbf{Examiner's Tip:} The Executive Summary is often the \textbf{only part} a busy examiner or manager reads fully. It must stand alone: Context $\rightarrow$ Problem $\rightarrow$ Solution $\rightarrow$ Evidence $\rightarrow$ Ask. No ``I'', no config, no fluff.
\end{exoresolu}

\begin{aretenir}[Chapter ICT02: Key Exam Pointers]
\begin{itemize}
\item \textbf{Complexity} is multi-dimensional: Topology, Heterogeneity, Virtualisation, Scale/Dynamics, Tooling. Score it, don't just describe it.
\item \textbf{Operational Challenges} map to four pillars: Performance (SLOs), Availability (MTBF/MTTR, SPOF), Security (Zero Trust, Compliance), Management (Maturity Model). Quantify gaps.
\item \textbf{Design:} Prefer \textbf{Collapsed Core} for $< 2000$ ports; \textbf{Routed Access} (L3 to Access) eliminates STP, enables ECMP, speeds convergence. Always summarise routes at Distribution $\rightarrow$ Core.
\item \textbf{HA:} HSRPv2/VRRPv3 \textbf{must track uplinks}. MLAG/vPC/StackWise Virtual for active-active L2. \texttt{spanning-tree port type network} on uplinks, \texttt{edge} on access.
\item \textbf{Monitoring:} Telemetry = SNMP + NetFlow + Syslog + IP SLA. Baseline = Dynamic (time-aware). Alert = Multi-level (Critical/Warning/Info) with Correlation.
\item \textbf{Troubleshooting:} Define $\rightarrow$ Gather (Show cmds) $\rightarrow$ Approach (Divide \& Conquer) $\rightarrow$ Hypothesise $\rightarrow$ Verify. Correlate timestamps across sources.
\item \textbf{Integration:} PPDIOO / Agile. HLD $\rightarrow$ LLD $\rightarrow$ Simulation (Digital Twin) $\rightarrow$ Staged Rollout (Rollback ready) $\rightarrow$ ATP $\rightarrow$ As-Built Handover.
\item \textbf{Reporting:} Executive Summary ($< 250$ w), Third-person passive, Numbered Figures/Tables with captions, Authoritative references (RFC, IEEE, Law), Version control.
\end{itemize}
\end{aretenir}

\begin{attention}[Common GCE Traps in ICT02]
\begin{itemize}
\item \textbf{Confusing L2 vs L3 Redundancy:} HSRP/VRRP = Gateway (L3) redundancy. STP = L2 loop prevention. MLAG/vPC = L2 Multi-chassis link aggregation (no STP blocking). Do not mix them up.
\item \textbf{Forgetting HSRP Tracking:} Configuring HSRP without \texttt{track} interface = Active router stays active when WAN fails $\rightarrow$ Blackhole. \textbf{Always track.}
\item \textbf{Static Thresholds in Monitoring:} ``Alert if CPU $> 80\%$'' fails at 03:00 (backup) or 12:00 (lunch dip). \textbf{Use Baselines ($\mu \pm 3\sigma$).}
\item \textbf{ECMP over Unequal Paths:} EIGRP/OSPF install equal-cost paths even if one is 100M and other is 10G (if metrics match). \textbf{Set bandwidth/delay correctly} or use \texttt{variance} / Floating Static.
\item \textbf{QoS without Shaping on Sub-rate Links:} Applying LLQ on a 1G port connected to 50M Metro-E = Queue never fills $\rightarrow$ No prioritisation. \textbf{Parent Shaper required.}
\item \textbf{Simulation $\neq$ Reality:} Packet Tracer/GNS3 do not model ASIC queues, TCAM exhaustion, CPU punt limits. State limitations in report.
\item \textbf{Vague References:} ``Cisco website'' scores 0. Cite: ``Cisco Catalyst 9500 Config Guide, IOS-XE 17.9, Chapter: Configuring StackWise Virtual''.
\end{itemize}
\end{attention}

\begin{exercice}[Self-Assessment: ICT02 Mastery Check]
\begin{itemize}
\item A network has 3 sites, 2 vendors, 1 overlay (SD-WAN), 50 VLANs, 10 changes/week. Score complexity (5 factors) and give overall index.
\item Define SLOs for a Trading App (Latency, Jitter, Loss, Availability). Current: 99.9\% uptime, 50 ms RTT, 5 ms jitter, 0.2\% loss. Target: 99.99\%, 20 ms, 2 ms, 0.05\%. Calculate gaps.
\item Design IPv4 VLSM for 4 buildings, 6 zones each, using 172.16.0.0/16. Show Building 1 allocation.
\item HSRP Active priority 110, Standby 100. Active WAN link fails. Standby becomes Active. WAN restores. Active \textbf{does not} revert. Why? Fix.
\item Write an Alert Rule (pseudo-code) for ``BGP Peer Flap $> 3$ times in 5 min''.
\item List 3 differences between Packet Tracer and GNS3/EVE-NG for SD-WAN validation.
\item In a project report, where does the ``Wireless Controller Recommendation'' belong if not in Conclusion?
\end{itemize}
\end{exercice}

\begin{corrige}[Exercise 1]
\textbf{1. Complexity Scorecard:} Topology: Medium (3 sites, hybrid WAN). Heterogeneity: Medium (2 vendors). Virtualisation: Medium (1 overlay). Scale: Low (50 VLANs, 10 CR/wk). Tooling: Medium (2 mgmt planes). Scores: 2,2,2,1,2 $\rightarrow$ Avg = 1.8 (Medium). \textbf{2. SLO Gaps:} Avail: 99.9\% $\rightarrow$ 99.99\% (Gap 0.09\% = 47 min/yr excess). Latency: 50 $\rightarrow$ 20 ms (Gap +150\%). Jitter: 5 $\rightarrow$ 2 ms (+150\%). Loss: 0.2\% $\rightarrow$ 0.05\% (+300\%). \textbf{3. VLSM Bldg 1 (/20 = 172.16.0.0/20):} Mgmt /24 (VLAN 10), Voice /24 (20), User /23 (30), IoT /23 (40), Server /24 (50), Guest /24 (60). \textbf{4. HSRP No Revert:} Missing \texttt{preempt} on Active (Priority 110). Fix: \texttt{standby 10 preempt}. \textbf{5. Alert Rule:} IF (bgp\_peer\_flap\_count > 3 IN 5m) THEN SEV=CRIT, RUNBOOK=``Check hold-timer, MTU, keepalive''. \textbf{6. PT vs GNS3:} PT: Simulated OS, limited scale, no real vendor images, no Linux VMs. GNS3: Real images (IOS-XE, NX-OS, Linux), KVM/QEMU, hardware resource heavy, supports NETCONF/gNCI. \textbf{7. Future Work section (Section 7.2).}
\end{corrige}

\newpage

\coursec{Exercises}

\courssub{I check my knowledge}

Before tackling the integration activities, make sure the key notions of this chapter --- redundancy, availability, hierarchical addressing, loop protection and change management --- are firmly acquired. The following exercises cover exactly the vocabulary and the calculations expected at the GCE Advanced Level.

\begin{exercice}[MCQ: complexity and operations]
For each question, choose the single correct answer.
\begin{enumerate}
\item The \emph{mean time between failures} (MTBF) of a device is:
\begin{enumerate}
\item the average time needed to repair the device;
\item the average time the device works correctly between two consecutive failures;
\item the guaranteed lifetime of the device;
\item the time needed to reboot the device.
\end{enumerate}
\item A \emph{single point of failure} (SPOF) is:
\begin{enumerate}
\item a component whose failure stops the whole service;
\item a faulty network card;
\item a switch with only one VLAN;
\item a router with a single user account.
\end{enumerate}
\item In the hierarchical IPv6 addressing of an enterprise receiving a /48, the prefix typically assigned to one \emph{site} is:
\begin{enumerate}
\item /32 \item /48 \item /56 or /64 \item /128
\end{enumerate}
\item BPDU Guard is normally configured on:
\begin{enumerate}
\item trunk ports between core switches;
\item access ports where end users connect (PortFast ports);
\item the console port;
\item the management VLAN only.
\end{enumerate}
\end{enumerate}
\end{exercice}

\begin{exercice}[True or false — justify]
Say whether each statement is true or false, and justify your answer in one or two sentences.
\begin{enumerate}
\item Adding redundant links between switches always improves the network, even without a spanning-tree protocol.
\item Availability of a service increases when MTBF increases and MTTR decreases.
\item In a change management process, a rollback plan is optional when the change is small.
\item A broadcast storm can be caused by a physical loop between switches.
\end{enumerate}
\end{exercice}

\begin{exercice}[Definitions]
Define each of the following terms in two or three sentences, and give one concrete example for each:
\begin{enumerate}
\item Redundancy;
\item MTTR (Mean Time To Repair);
\item Broadcast storm;
\item Change Request (CR);
\item Network supervision (monitoring).
\end{enumerate}
\end{exercice}

\begin{exercice}[Direct calculation: availability]
A school server in Bamenda has an MTBF of $2000$ hours and an MTTR of $4$ hours.
\begin{enumerate}
\item Calculate its availability $A$ using $A = \dfrac{\text{MTBF}}{\text{MTBF}+\text{MTTR}}$.
\item Express this availability as a percentage, rounded to two decimal places.
\item How many hours of downtime does this represent over one year of $8760$ hours?
\item The school wants ``three nines'' availability ($99.9\%$). With the same MTBF, what maximum MTTR is allowed?
\end{enumerate}
\end{exercice}

\courssub{I practise}

\begin{exercice}[Identifying single points of failure]
The diagram below shows the network of a small company in Douala. The two core switches are linked by a single cable; each access switch is connected to one core switch only; the router has one link to the core and one link to the ISP.
\begin{popfigure}
\begin{center}
\begin{tikzpicture}[scale=0.95, every node/.style={font=\small}]
\node[draw, rounded corners, fill=popblueL, minimum width=2.2cm, minimum height=0.8cm] (r) at (0,3) {Router R1};
\node[draw, rounded corners, fill=popgreenL, minimum width=2cm, minimum height=0.8cm] (c1) at (-2,1.5) {Core SW1};
\node[draw, rounded corners, fill=popgreenL, minimum width=2cm, minimum height=0.8cm] (c2) at (2,1.5) {Core SW2};
\node[draw, rounded corners, minimum width=2cm, minimum height=0.8cm] (a1) at (-3,0) {Access SW-A};
\node[draw, rounded corners, minimum width=2cm, minimum height=0.8cm] (a2) at (0,0) {Access SW-B};
\node[draw, rounded corners, minimum width=2cm, minimum height=0.8cm] (a3) at (3,0) {Access SW-C};
\node[draw, ellipse, fill=popblueL, minimum width=1.8cm, minimum height=1cm] (isp) at (0,4.6) {ISP};
\draw[thick] (isp) -- (r);
\draw[thick, poporange] (r) -- (c1);
\draw[thick, poporange] (r) -- (c2);
\draw[thick] (c1) -- (c2);
\draw[thick] (c1) -- (a1);
\draw[thick] (c1) -- (a2);
\draw[thick] (c2) -- (a3);
\end{tikzpicture}
\end{center}
\legende{Network of the company (single uplinks everywhere).}
\end{popfigure}
\begin{enumerate}
\item List all the single points of failure (links and devices) in this network.
\item For each SPOF, state which users or services are affected if it fails.
\item Propose a redundancy solution for each SPOF (for example VRRP between two routers, LACP link aggregation between the core switches, dual uplinks on access switches).
\item The current router has MTBF $= 4000$ h and MTTR $= 8$ h. Compute its availability, then explain how a second router with VRRP improves the \emph{effective} availability of the gateway service.
\end{enumerate}
\end{exercice}

\begin{exercice}[IPv6 addressing plan]
A Cameroonian bank receives the prefix \texttt{2001:db8:acad::/48} from its provider. It must address $12$ branches, each branch having $8$ VLANs.
\begin{enumerate}
\item How many bits are available for subnetting between /48 and /64?
\item Propose a hierarchical plan: choose a site prefix length (for example /56) and show that it can number at least $12$ sites.
\item Inside one site, show that the remaining bits can number at least $8$ VLANs of /64 each.
\item Write, in hexadecimal, the prefix of site number $3$ and the prefix of VLAN number $5$ inside that site (numbering starts at $1$).
\item Give two advantages of this hierarchical plan for routing and management.
\end{enumerate}
\end{exercice}

\begin{exercice}[Reading a broadcast storm]
A Wireshark capture on a school LAN shows, during $10$ seconds, $180\,000$ broadcast frames (ARP and unknown unicast flooding), while normal traffic is about $200$ broadcast frames per second. The CPU of the access switches reaches $95\%$ and users complain that ``the network is frozen''.
\begin{enumerate}
\item What is the most likely physical cause of these symptoms? Justify.
\item Which protocol should normally prevent this situation, and how does it work in one sentence?
\item A technician temporarily disabled this protocol on two switches ``to speed up convergence''. Explain why this was dangerous.
\item Propose a correct configuration approach: which protocol version (STP or RSTP), which ports should get PortFast, and where should BPDU Guard be enabled? Justify each choice.
\end{enumerate}
\end{exercice}

\begin{exercice}[Supervision indicators]
The IT team of a Douala ISP supervises a fibre link with the following measurements over one month ($30$ days, i.e. $720$ hours):
\begin{itemize}
\item $3$ failures occurred, lasting respectively $2$ h, $30$ min and $1$ h;
\item the link carries on average $620$ Mbit/s on a $1$ Gbit/s capacity.
\end{itemize}
\begin{enumerate}
\item Compute the MTBF, the MTTR and the availability of the link for that month.
\item Compute the average utilisation rate of the link. Is the link overloaded? Justify.
\item Give two concrete supervision actions (tools or alerts) the team should put in place to detect the next failure faster.
\end{enumerate}
\end{exercice}

\courssub{I prepare for the GCE}

\begin{exercice}[GCE-style: redundancy and availability — 20 marks]
A hospital in Yaoundé runs a patient records system on servers connected to two core switches. The current architecture has: one router connected to one ISP, one link between the two core switches, and each access switch connected to a single core switch. Downtime of the records system is estimated to cost the hospital $150\,000$ FCFA per hour.
\begin{enumerate}
\item Identify four single points of failure in this architecture and, for each, state the impact of its failure. \hfill (4 marks)
\item Propose a redundancy solution for each SPOF, naming the appropriate technology (VRRP/HSRP, LACP, dual-homing, dual ISP) and justifying your choice. \hfill (8 marks)
\item The gateway router has MTBF $= 3000$ h and MTTR $= 6$ h. Calculate its availability and the expected downtime cost per year ($8760$ h). \hfill (4 marks)
\item Explain, in at most six lines, why adding redundancy increases the \emph{complexity} of the network and name two operational practices that help manage this complexity. \hfill (4 marks)
\end{enumerate}
\end{exercice}

\begin{exercice}[GCE-style: design and addressing — 20 marks]
The Ministry of Secondary Education wants to interconnect $10$ regional delegations. Each delegation has $6$ services (VLANs): Administration, Exams, Finance, Teachers, Students, Guests. The provider allocates the IPv6 prefix \texttt{2001:db8:educ::/48}.
\begin{enumerate}
\item Draw (or describe precisely) a hierarchical topology with a central site, regional delegations, and redundant links where you judge them necessary. Justify two of your redundancy choices. \hfill (6 marks)
\item Design the addressing plan: choose the site prefix length, show it covers $10$ sites, and show each site can hold at least $6$ VLANs of /64. Write the prefix of the Exams VLAN (VLAN $3$) of delegation $4$. \hfill (8 marks)
\item Give two reasons why a flat addressing plan (no site hierarchy) would make routing and troubleshooting harder. \hfill (3 marks)
\item State three supervision indicators the Ministry should monitor on the inter-delegation links, and for each one the threshold that should trigger an alert. \hfill (3 marks)
\end{enumerate}
\end{exercice}

\begin{exercice}[GCE-style: change management essay — 20 marks]
The core switches of a university campus network in Buea must receive a firmware update that fixes a security vulnerability. The update requires a reboot of each switch. The network carries e-learning, VoIP telephony and the payment system of the campus restaurant.

Write a structured essay (about one page) presenting a complete \textbf{Change Request} for this operation. Your essay must include:
\begin{enumerate}
\item the justification and objectives of the change; \hfill (3 marks)
\item a risk analysis (what can go wrong, which services are affected); \hfill (4 marks)
\item a step-by-step implementation plan, including the choice of a maintenance window and the order in which redundant switches are updated; \hfill (5 marks)
\item a rollback plan in case of failure; \hfill (4 marks)
\item the communication plan towards stakeholders (users, management, technicians) before, during and after the change. \hfill (4 marks)
\end{enumerate}
\end{exercice}



\newpage

\coursec{Solutions to the exercises}

\coursec{Activités d'intégration : exercices corrigés}

\begin{corrige}[Exercice 1 --- QCM : complexité et exploitation]
\begin{enumerate}
\item \textbf{Réponse (b).} Le \cle{MTBF} (Mean Time Between Failures) est le temps moyen de \emph{bon fonctionnement} entre deux pannes consécutives. Le temps moyen de réparation est un autre indicateur, le \cle{MTTR} ; la durée de vie garantie est fixée par le constructeur ; le temps de redémarrage n'est pas pertinent ici.
\item \textbf{Réponse (a).} Un \cle{point unique de défaillance} (SPOF) est tout composant (équipement, liaison, alimentation) dont la panne \emph{seule} interrompt le service global. Une carte réseau défectueuse ou un routeur mono-utilisateur ne sont des SPOF que si tout le service en dépend.
\item \textbf{Réponse (c).} Avec une allocation /48, un \emph{site} reçoit généralement un /56 (ou un unique /64 pour un tout petit site). Le /48 est l'allocation de l'entreprise entière, le /32 correspond au niveau registre/fournisseur, et le /128 identifie un seul hôte.
\item \textbf{Réponse (b).} Le \cle{BPDU Guard} se configure sur les ports \emph{d'accès} où PortFast est activé (ports pour PC et imprimantes). Si un commutateur --- qui émet des BPDU --- y est branché, le port est immédiatement mis en erreur (shutdown), évitant la boucle. Les ports trunk entre commutateurs doivent \emph{recevoir} des BPDU, donc BPDU Guard ne doit jamais y être activé.
\end{enumerate}
\end{corrige}

\begin{corrige}[Exercice 2 --- Vrai ou faux --- justifier]
\begin{enumerate}
\item \textbf{Faux.} Sans protocole d'arbre couvrant, un lien redondant crée une \emph{boucle physique} : les trames de broadcast circulent indéfiniment et sont dupliquées à chaque passage (tempête de broadcast), les CPU des commutateurs saturent et le réseau s'effondre. La redondance n'est bénéfique que si STP/RSTP (ou l'agrégation de liens) contrôle le chemin supplémentaire.
\item \textbf{Vrai.} Puisque
\[ A = \frac{\text{MTBF}}{\text{MTBF}+\text{MTTR}}, \]
augmenter MTBF ou diminuer MTTR rend MTTR relativement plus petit devant MTBF, donc $A$ tend vers $1$ (c'est-à-dire vers $100\%$).
\item \textbf{Faux.} Tout changement peut échouer, même ``mineur'' --- une erreur de configuration sur une ligne peut isoler un site entier. Un \cle{plan de retour arrière} (rollback : restauration de la configuration ou du firmware précédent) est une section obligatoire de \emph{toute} Demande de Changement (Change Request), quelle que soit l'ampleur du changement.
\item \textbf{Vrai.} Sur une boucle, une trame de broadcast est réinjectée sur la boucle par chaque commutateur et dupliquée sans fin ; le nombre de trames croît exponentiellement jusqu'à saturer les liaisons et les CPU des commutateurs --- c'est précisément une tempête de broadcast.
\end{enumerate}
\end{corrige}

\begin{corrige}[Exercice 3 --- Définitions]
\begin{enumerate}
\item \textbf{Redondance} : duplication d'un composant critique (liaison, commutateur, routeur, alimentation) pour que le service survive à la panne d'un élément. \emph{Exemple} : deux commutateurs cœurs reliés par deux câbles agrégés, de sorte qu'une coupure de l'un n'isole pas le réseau.
\item \textbf{MTTR (Mean Time To Repair)} : temps moyen nécessaire pour détecter, diagnostiquer, réparer une panne et rétablir le service. \emph{Exemple} : si trois pannes ont duré $2$ h, $1$ h et $30$ min, $\text{MTTR} = \dfrac{2 + 1 + 0.5}{3} \approx 1.17$ h, soit environ $1$ h $10$ min.
\item \textbf{Tempête de broadcast (broadcast storm)} : multiplication incontrôlée de trames de broadcast, due le plus souvent à une boucle de commutation sans STP, qui sature la bande passante et les CPU des commutateurs. \emph{Exemple} : un câble branché entre deux ports du même commutateur de salle de classe fige le LAN de l'établissement en quelques secondes.
\item \textbf{Demande de changement (Change Request --- CR)} : document formel décrivant une modification planifiée --- justification, risques, étapes d'implémentation, plan de retour arrière, communication --- qui doit être approuvée avant exécution. \emph{Exemple} : le CR rédigé avant la mise à jour du firmware des commutateurs cœurs du campus.
\item \textbf{Supervision réseau (monitoring)} : surveillance continue des équipements et liaisons (disponibilité, trafic, CPU, erreurs) via des plateformes SNMP (Zabbix, Nagios, PRTG), avec alertes vers les administrateurs dès qu'un seuil est franchi. \emph{Exemple} : un courriel/SMS alerte le technicien quand l'utilisation de la liaison Internet de l'école dépasse $80\%$.
\end{enumerate}
\end{corrige}

\begin{corrige}[Exercice 4 --- Calcul direct : disponibilité]
\begin{enumerate}
\item Application de la formule :
\[ A = \frac{\text{MTBF}}{\text{MTBF}+\text{MTTR}} = \frac{2000}{2000+4} = \frac{2000}{2004} \approx 0.998004. \]
\item En pourcentage : $A \approx 99.80\%$.
\item L'indisponibilité est $1 - A \approx 0.001996$. Sur un an :
\[ 8760 \times 0.001996 \approx 17.5\ \text{heures d'arrêt par an}. \]
\item On impose
\[ \frac{2000}{2000+\text{MTTR}} \geq 0.999
\;\Longleftrightarrow\;
2000 + \text{MTTR} \leq \frac{2000}{0.999} \approx 2002.0
\;\Longleftrightarrow\;
\text{MTTR} \leq 2\ \text{h (environ)}. \]
L'établissement doit donc réparer le serveur en moins de $2$ heures en moyenne --- ou, plus réalistement, ajouter un \emph{serveur redondant} pour que le service survive à une panne matérielle.
\end{enumerate}
\end{corrige}

\begin{corrige}[Exercice 5 --- Identification des points uniques de défaillance]
\begin{enumerate}
\item Les SPOF sont : le routeur R1 lui-même ; la liaison FAI ; le câble unique entre SW1 et SW2 ; les uplinks SW-A--SW1, SW-B--SW1, SW-C--SW2 ; et chaque commutateur cœur \emph{pour les commutateurs d'accès qui n'y sont raccordés qu'à lui} (SW1 pour A et B, SW2 pour C).
\item Impact de chaque panne :
\begin{itemize}
\item R1 ou liaison FAI $\Rightarrow$ plus d'Internet pour \emph{toute} l'entreprise ;
\item câble SW1--SW2 $\Rightarrow$ SW-C isolé du routeur et des commutateurs A et B ;
\item uplink d'accès $\Rightarrow$ tous les utilisateurs de ce commutateur d'accès coupés ;
\item SW1 $\Rightarrow$ commutateurs d'accès A et B hors service (idem SW2 pour C).
\end{itemize}
\item Solutions de redondance :
\begin{itemize}
\item Routeur : ajouter un second routeur avec \textbf{VRRP} --- une passerelle virtuelle unique, basculement automatique, transparent pour les utilisateurs ;
\item Liaison FAI : second abonnement (\textbf{dual ISP}), par exemple fibre CAMTEL + secours 4G ;
\item Inter-cœur : second câble en \textbf{LACP} (agrégation de liens) --- plus de bande passante \emph{et} pas de risque de boucle, le faisceau étant un lien logique unique ;
\item Commutateurs d'accès : \textbf{double uplink} vers les deux cœurs, avec RSTP pour gérer les boucles.
\end{itemize}
\item Disponibilité du routeur actuel :
\[ A = \frac{4000}{4000+8} = \frac{4000}{4008} \approx 0.998004 \approx 99.80\%, \]
soit environ $17.5$ h d'indisponibilité de la passerelle par an. Avec VRRP, le service passerelle ne tombe que si \emph{les deux} routeurs sont en panne \emph{simultanément} ; la probabilité de panne simultanée est le \emph{produit} des deux indisponibilités individuelles ($\approx 0.002 \times 0.002 = 4\times10^{-6}$), donc la disponibilité effective de la passerelle dépasse largement $99.99\%$.
\end{enumerate}
\end{corrige}

\begin{corrige}[Exercice 6 --- Plan d'adressage IPv6]
\begin{enumerate}
\item Entre /48 et /64 il y a $64 - 48 = 16$ bits disponibles pour le sous-réseautage.
\item Choix : /56 par site. Le champ site a $56 - 48 = 8$ bits, soit $2^8 = 256$ sites possibles --- bien plus que les $12$ succursales requises, avec une marge confortable pour l'avenir.
\item À l'intérieur d'un site /56, $64 - 56 = 8$ bits restent pour les VLAN : $2^8 = 256$ sous-réseaux /64 --- bien plus que les $8$ VLAN requis.
\item Le champ site occupe l'octet de poids fort du quatrième hextet et le champ VLAN son octet de poids faible. En numérotant à partir de $1$ :
\begin{itemize}
\item Site $3$ : \texttt{2001:db8:acad:0300::/56} ;
\item VLAN $5$ du site $3$ : \texttt{2001:db8:acad:0305::/64}.
\end{itemize}
\item Deux avantages :
\begin{enumerate}
\item \emph{Agrégation de routes} : le cœur n'a besoin que d'une route récapitulative par site (/56), les tables de routage restent petites et convergent vite ;
\item \emph{Gestion et dépannage aisés} : le préfixe encode le site et le VLAN, un administrateur lisant une adresse sait immédiatement où se trouve la machine.
\end{enumerate}
\end{enumerate}
\end{corrige}

\begin{corrige}[Exercice 7 --- Lecture d'une tempête de broadcast]
\begin{enumerate}
\item La cause la plus probable est une \emph{boucle de commutation} : un câble relie deux commutateurs (ou deux ports) en boucle alors que l'arbre couvrant ne bloque aucun port. Les trames de broadcast circulent alors indéfiniment et sont dupliquées à chaque passage --- ce qui explique l'explosion de $200$ à $18\,000$ trames de broadcast par seconde et la saturation CPU à $95\%$.
\item Le \textbf{Spanning Tree Protocol} (STP, ou sa version rapide RSTP) : il détecte les boucles et bloque logiquement les ports redondants, conservant un arbre sans boucle tout en gardant des chemins de secours en cas de panne.
\item La désactivation de STP a supprimé la \emph{seule} protection contre les boucles : dès qu'un câble redondant ou mal branché a créé une boucle, rien ne l'a bloquée et la tempête a démarré immédiatement. ``Accélérer la convergence'' n'est jamais une raison valable pour retirer la protection anti-boucle.
\item Bonne démarche :
\begin{itemize}
\item Utiliser \textbf{RSTP} (802.1w) partout : il converge en environ une seconde (contre $30$--$50$ s pour STP classique), donc aucune raison de le désactiver ;
\item Activer \textbf{PortFast} uniquement sur les \emph{ports d'accès} vers terminaux (PC, imprimantes), pour qu'ils montent immédiatement sans participer au calcul de topologie ;
\item Activer \textbf{BPDU Guard} sur ces mêmes ports PortFast : si un utilisateur y branche un commutateur non autorisé, le port est mis en erreur au lieu de créer une boucle.
\end{itemize}
\end{enumerate}
\end{corrige}

\begin{corrige}[Exercice 8 --- Indicateurs de supervision]
\begin{enumerate}
\item Temps d'arrêt total : $2 + 0.5 + 1 = 3.5$ h, donc temps de fonctionnement total $= 720 - 3.5 = 716.5$ h.
\[ \text{MTTR} = \frac{3.5}{3} \approx 1.17\ \text{h} \quad (\text{environ } 1\ \text{h}\ 10\ \text{min}); \]
\[ \text{MTBF} = \frac{716.5}{3} \approx 238.8\ \text{h}; \]
\[ A = \frac{238.8}{238.8 + 1.17} \approx 0.9951 \approx 99.51\% \quad \left(\text{équivalent à } 1 - \frac{3.5}{720}\right). \]
\item Utilisation moyenne :
\[ \frac{620}{1000} = 62\%. \]
La liaison n'est \emph{pas} surchargée : elle est légèrement au-dessus du seuil de confort d'exploitation (environ $50\%$--$60\%$ pour une liaison de production), il faut donc surveiller la capacité et planifier une montée en débit, mais il reste de la marge.
\item Deux actions concrètes :
\begin{enumerate}
\item Déployer une plateforme de supervision SNMP (Zabbix, Nagios ou PRTG) interrogeant l'interface chaque minute, avec alerte automatique courriel/SMS dès que l'interface tombe ;
\item Configurer une alerte d'utilisation à $80\%$ et une alerte de taux d'erreurs sur l'interface optique, pour détecter une dégradation \emph{avant} la panne totale.
\end{enumerate}
\end{enumerate}
\end{corrige}

\begin{corrige}[Exercice 9 --- Type GCE : redondance et disponibilité]
\textbf{Barème indicatif (20 points).}
\begin{enumerate}
\item \textbf{(4 points --- 1 par SPOF correct avec son impact.)} N'importe lesquels de :
\begin{itemize}
\item le routeur unique : plus de passerelle $\Rightarrow$ le système de dossiers devient inaccessible depuis les autres sites/Internet ;
\item la liaison FAI unique : perte totale de la connectivité externe ;
\item le lien inter-cœur unique : les deux moitiés du LAN sont séparées ;
\item chaque uplink unique de commutateur d'accès : tous les utilisateurs de ce commutateur coupés ;
\item chaque commutateur cœur, pour les commutateurs d'accès qui n'y sont attachés qu'à lui.
\end{itemize}
\item \textbf{(8 points --- 2 par solution : 1 pour la technologie, 1 pour la justification.)}
\begin{itemize}
\item Routeur $\rightarrow$ second routeur avec \textbf{VRRP/HSRP} : une IP de passerelle virtuelle, basculement automatique, transparent pour les utilisateurs ;
\item FAI $\rightarrow$ \textbf{dual ISP} (ex. fibre CAMTEL + secours 4G ou seconde fibre) : protège contre la panne du fournisseur ;
\item Inter-cœur $\rightarrow$ second câble en \textbf{LACP} : plus de bande passante \emph{et} pas de risque de boucle, le faisceau étant un lien logique unique ;
\item Commutateurs d'accès $\rightarrow$ \textbf{double raccordement} (dual-homing) aux deux cœurs avec RSTP : l'uplink de secours prend le relais si le principal tombe.
\end{itemize}
\item \textbf{(4 points --- 2 pour $A$, 2 pour le coût.)}
\[ A = \frac{3000}{3000+6} = \frac{3000}{3006} \approx 0.998004 \approx 99.80\%. \]
Indisponibilité $\approx 8760 \times (1 - 0.998004) \approx 17.5$ h/an, donc un coût de
\[ 17.5 \times 150\,000 \approx 2\,620\,000\ \text{FCFA par an} \]
(accepter les réponses entre $2.6$ et $2.7$ millions FCFA).
\item \textbf{(4 points --- 2 pour l'explication, 1 par pratique.)} La redondance ajoute des équipements, liaisons et protocoles (VRRP, LACP, RSTP) qui doivent être configurés de façon cohérente et qui interagissent ; boucles, erreurs de configuration et dépannage plus difficile deviennent possibles, et il y a plus d'éléments à superviser. Deux pratiques : \emph{documentation rigoureuse} (schémas et configurations à jour) et \emph{gestion des changements} (toute modification via une Demande de Changement approuvée avec plan de retour arrière) ; la \emph{supervision centralisée} est aussi acceptable.
\end{enumerate}
\textbf{Attendus de l'examinateur :} vocabulaire précis (SPOF, VRRP, LACP, dual-homing), formule de disponibilité correcte avec unités, réponses justifiées en une ou deux phrases --- pas de listes sèches.
\end{corrige}

\begin{corrige}[Exercice 10 --- Type GCE : conception et adressage]
\textbf{Barème indicatif (20 points).}
\begin{enumerate}
\item \textbf{(6 points --- 4 pour une topologie trois niveaux cohérente, 2 pour les deux choix de redondance justifiés.)} Attendu : un site central (siège du ministère) avec deux routeurs/commutateurs cœurs ; chaque délégation régionale raccordée au centre par une liaison WAN principale ; les délégations les plus critiques (ou toutes, budget permitting) avec une liaison de secours (second fournisseur ou 4G). Justifications, ex. : deux routeurs cœurs avec VRRP au centre car le centre est le hub par où passent toutes les délégations --- sa panne isole tout le monde ; liaisons WAN de secours car une coupure de fibre ferait taire une délégation entière, inacceptable en période d'examens.
\item \textbf{(8 points --- 2 pour le choix du préfixe site, 2 pour le nombre de sites, 2 pour le nombre de VLAN, 2 pour le préfixe écrit.)} Choix : /56 par site. $56 - 48 = 8$ bits $\Rightarrow 2^8 = 256$ sites $\geq 10$. Dans un site, $64 - 56 = 8$ bits $\Rightarrow 256$ VLAN de /64 $\geq 6$. VLAN Examens (VLAN $3$) de la délégation $4$ :
\[ \texttt{2001:db8:educ:0403::/64}. \]
\item \textbf{(3 points --- 1.5 par raison.)} Avec un plan plat : (i) pas d'agrégation de routes possible, donc tables de routage grosses et convergence lente ; (ii) une adresse n'indique plus son site, localiser une machine défaillante oblige à consulter une base externe au lieu de lire simplement le préfixe.
\item \textbf{(3 points --- 1 par indicateur avec seuil réaliste.)} N'importe lesquels de :
\begin{itemize}
\item disponibilité/état de la liaison : alerte immédiate sur interface down ;
\item utilisation bande passante : alerte au-dessus de $80\%$ soutenu ;
\item latence (aller-retour) : alerte au-dessus d'une valeur convenue, ex. $100$ ms sur le backbone national ;
\item taux d'erreurs/pertes paquets : alerte au-dessus de $1\%$.
\end{itemize}
\end{enumerate}
\textbf{Attendus de l'examinateur :} description ou schéma hiérarchique clair, arithmétique binaire correcte ($2^8 = 256$), préfixes hexadécimaux bien formés, seuils réalistes et non arbitraires.
\end{corrige}

\begin{corrige}[Exercice 11 --- Type GCE : dissertation sur la gestion des changements]
\textbf{Guide de notation (20 points).} Attribuer les points pour une Demande de Changement (CR) structurée, réaliste, contenant les éléments ci-dessous ; la dissertation doit être organisée (titres ou paragraphes nets), en anglais correct, d'environ une page.
\begin{enumerate}
\item \textbf{Justification et objectifs (3 points)} : le firmware actuel présente une vulnérabilité de sécurité connue permettant à un attaquant de prendre le contrôle des commutateurs cœurs ; l'objectif est de patcher tous les cœurs avec interruption de service minimale et aucune perte de données.
\item \textbf{Analyse des risques (4 points)} : le redémarrage coupe l'e-learning, la VoIP et les paiements au restaurant ; le nouveau firmware peut échouer au démarrage ou être incompatible avec la configuration existante ; une erreur humaine est possible pendant l'opération. Le risque est réduit par la redondance (le second cœur continue de fonctionner) et en travaillant hors heures de cours.
\item \textbf{Plan d'implémentation (5 points)} : programmer une fenêtre de maintenance (ex. samedi 22:00--02:00) quand l'usage est minimal ; sauvegarder toutes les configurations et télécharger l'ancien et le nouveau firmware ; mettre à jour le cœur \emph{standby} d'abord, le vérifier (VLAN, routage, état VRRP), basculer le trafic dessus, puis mettre à jour le second cœur ; tester les trois services après chaque étape ; documenter chaque action.
\item \textbf{Plan de retour arrière (4 points)} : si le nouveau firmware échoue, recharger l'ancien firmware sauvegardé et la configuration de secours sur le commutateur concerné, vérifier le rétablissement du service, et reporter le changement ; définir une échéance de décision claire (ex. si le commutateur n'est pas stable à 01:00, on fait le rollback).
\item \textbf{Plan de communication (4 points)} : \emph{avant} --- courriel/affichage au personnel et aux étudiants et validation de la direction, avec date et impact attendu ; \emph{pendant} --- canal de statut pour les techniciens et point de contact pour les urgences ; \emph{après} --- message de confirmation que les services sont restaurés, plus rapport d'incident/changement à la direction.
\end{enumerate}
\textbf{Attendus de l'examinateur :} les cinq sections du CR doivent toutes être présentes et dans un ordre logique ; la fenêtre de maintenance et la stratégie ``cœur standby d'abord'' sont des détails professionnels clés ; un CR sans plan de retour arrière ou sans communication ne peut pas avoir la note maximale, quelle que soit sa qualité technique.
\end{corrige}

\newpage

\coursec{Summary sheet}

\begin{popfigure}
\begin{tikzpicture}[grow cyclic, mindmap, concept color=popblue,
  level 1 concept/.append style={level distance=3.6cm, sibling angle=51.4, font=\small},
  level 2 concept/.append style={level distance=2.2cm, font=\scriptsize},
  every node/.style={concept, text=white},
  scale=0.92, transform shape]
\node[concept, font=\small\bfseries]{Computer Networking: Complexity \& Challenges}
  child[concept color=poporange]{ node {Factors of complexity}
    child { node {Scale \& growth} }
    child { node {Heterogeneity} }
    child { node {Mobility \& cloud} }
  }
  child[concept color=popgreen]{ node {Performance}
    child { node {Bandwidth, delay} }
    child { node {QoS, congestion} }
  }
  child[concept color=poppurple]{ node {Availability}
    child { node {Redundancy} }
    child { node {Backup \& failover} }
  }
  child[concept color=popgold]{ node {Security}
    child { node {CIA triad} }
    child { node {Firewall, VPN} }
  }
  child[concept color=poppink]{ node {Management}
    child { node {SNMP monitoring} }
    child { node {Documentation} }
  }
  child[concept color=popgreen!70!black]{ node {Design \& deployment}
    child { node {Needs analysis} }
    child { node {Addressing plan} }
  }
  child[concept color=popblue!60!popdark]{ node {Integration activities}
    child { node {Case study} }
    child { node {Simulation \& report} }
  };
\end{tikzpicture}
\legende{Mind map of the chapter: from the sources of network complexity to the design, supervision and integration activities.}
\end{popfigure}

\begin{bilan}
\textbf{1. Key definitions}\par
\begin{itemize}
\item \cle{Network complexity}: the degree of difficulty in designing, operating and troubleshooting a network, caused by its \textbf{size} (number of hosts), its \textbf{heterogeneity} (different hardware, operating systems, protocols), its \textbf{topology} and its constant \textbf{evolution} (new users, cloud services, mobile devices).
\item \cle{Performance}: the ability of the network to deliver data quickly and reliably; measured by \textbf{bandwidth} (bit/s), \textbf{throughput} (effective data rate), \textbf{delay/latency} (ms) and \textbf{packet loss} (\%).
\item \cle{Availability}: the proportion of time the service is operational: $A = \dfrac{\text{uptime}}{\text{uptime}+\text{downtime}} \times 100\ \%$. High availability is obtained through \textbf{redundancy} (duplicate links, servers, power supplies).
\item \cle{Security}: protection based on the \textbf{CIA triad}: \emph{Confidentiality} (only authorised persons read the data), \emph{Integrity} (data is not altered), \emph{Availability} (services remain accessible).
\item \cle{Network management / supervision}: continuous monitoring of devices and links (using protocols such as \textbf{SNMP}), fault detection, configuration, accounting and performance analysis.
\item \cle{QoS} (Quality of Service): techniques that prioritise sensitive traffic (voice, video) over ordinary traffic.
\end{itemize}

\textbf{2. Factors of complexity in modern networks}\par
\begin{itemize}
\item Growth of the number of connected devices (computers, phones, IoT sensors).
\item Mix of technologies: wired (fibre, UTP), wireless (Wi-Fi, 4G), cloud services.
\item Interconnection of many protocols and vendors; compatibility problems.
\item Mobility of users and remote work (VPN connections).
\item Increasing cyber-threats (viruses, phishing, denial of service).
\end{itemize}

\textbf{3. Operational challenges}\par
\begin{center}
\begin{tabularx}{\linewidth}{|l|X|}\hline
\rowcolor{popblueL} \textbf{Challenge} & \textbf{Typical solutions} \\ \hline
Performance & Upgrade bandwidth, QoS, reduce congestion, segment the LAN with switches/VLANs \\ \hline
Availability & Redundant links and servers, UPS, regular backups, failover mechanisms \\ \hline
Security & Firewall, antivirus, strong passwords, VPN, user authentication, updates \\ \hline
Management & SNMP monitoring tools, log files, documentation, trained administrators \\ \hline
\end{tabularx}
\end{center}

\textbf{4. Method: designing and deploying an enterprise network}\par
\begin{enumerate}
\item \textbf{Analyse the needs}: number of users, services (file sharing, Internet, VoIP), budget, building plan.
\item \textbf{Choose the topology} (star is the most common in LANs) and the equipment (routers, switches, access points, servers).
\item \textbf{Plan the IP addressing}: choose network classes, subnet masks; divide into subnets/VLANs per department.
\item \textbf{Deploy and configure}: cabling, device configuration, security rules (firewall, access rights).
\item \textbf{Test and supervise}: connectivity tests (ping, traceroute), monitoring with SNMP, write the documentation.
\end{enumerate}

\textbf{5. Common mistakes to avoid}\par
\begin{itemize}
\item Confusing \emph{bandwidth} (capacity of the link) with \emph{throughput} (what is actually achieved).
\item Believing that one firewall alone guarantees security: security also needs updates, passwords and user training.
\item Forgetting redundancy and backups: a single point of failure can stop the whole enterprise.
\item Designing a network without an addressing plan or documentation: maintenance becomes impossible.
\item Ignoring the human factor: most security incidents come from user errors (weak passwords, phishing).
\end{itemize}

\textbf{6. GCE examination tips}\par
\begin{itemize}
\item Define each term precisely before explaining: examiners reward exact definitions (availability, QoS, CIA triad).
\item In a case study, always follow the design steps in order: \emph{needs $\rightarrow$ topology $\rightarrow$ addressing $\rightarrow$ security $\rightarrow$ supervision}.
\item Know how to compute availability in \% and interpret a bandwidth or delay value.
\item For integration activities: present a clear diagram, justify every technical choice, and conclude with limits and possible improvements.
\item Use local examples naturally (a school network in Yaound\'e, a mobile money platform, a cybercaf\'e) to illustrate, but keep the technical vocabulary exact.
\end{itemize}
\end{bilan}

\findecours

\end{document}
